% Lectura y comentario : derechos de las mujeres y de las minorías — Espagnol Tle A — Cours signé OnBuch+
% Programme officiel MINESEC. Compiler avec : tectonic E03-derechos-mujeres-minorias.tex
\newcommand{\DOCMATIERE}{Espagnol}
\newcommand{\DOCNIVEAU}{Tle A}
\newcommand{\DOCMODULE}{Módulo 1 : Vie familiale et intégration sociale}
\newcommand{\DOCLECON}{E03}
\newcommand{\DOCTITRE}{Lectura y comentario : derechos de las mujeres y de las minorías}
% OnBuch+ — préambule « cours signés » ENRICHI (compilé avec tectonic / XeLaTeX)
% Système visuel « Pop » d'OnBuch+ : fond crème (jamais blanc), bordures encre
% épaisses, ombres « dures » décalées, titres Archivo Black en capitales.
% Version MATHÉMATIQUES : graphes (pgfplots/tikz), boîtes pédagogiques
% (définition, propriété, méthode, attention, le savais-tu, exercice résolu,
% exercice, TP, bilan), page de garde et signature OnBuch+ ancrée en bas.
%
% Macros attendues dans main.tex AVANT \input{preamble} :
%   \DOCMATIERE  \DOCNIVEAU  \DOCMODULE  \DOCLECON (numéro)  \DOCTITRE
\documentclass[11pt]{article}

\usepackage{fontspec}
\usepackage[spanish,french]{babel} % français = langue principale ; l'espagnol via \esp{..} / espagnol
\usepackage[a4paper,margin=2cm,top=3.4cm,bottom=2.8cm,headheight=1.4cm,footskip=1.3cm]{geometry}
\usepackage{amsmath,amssymb,amsfonts}
\usepackage{mathtools} % \xmapsto, \coloneqq... souvent utilisés par les modèles
\usepackage{cancel} % simplifications barrées (bilans de forces)
% \abs{z} (module, valeur absolue) et \norm{u} (norme), utilisés par les modèles
\providecommand{\abs}{}\let\abs\relax\DeclarePairedDelimiter{\abs}{\lvert}{\rvert}
\providecommand{\norm}{}\let\norm\relax\DeclarePairedDelimiter{\norm}{\lVert}{\rVert}
\usepackage[table]{xcolor} % [table] : fournit aussi \rowcolors (alternance de lignes)
\usepackage{pagecolor}
\usepackage{tikz}
\usetikzlibrary{arrows.meta,positioning,calc,shapes.geometric,shapes.misc,shapes.arrows,decorations.pathmorphing,decorations.pathreplacing,decorations.markings,patterns,shadows,mindmap,trees,fit,backgrounds,matrix,intersections,angles,quotes}
\usepackage{pgfplots}
\usepackage[version=4]{mhchem} % équations nucléaires \ce{^{235}_{92}U}
\usepackage[european,siunitx]{circuitikz} % schémas électriques
\pgfplotsset{compat=1.18}
\usepgfplotslibrary{fillbetween,groupplots} % aires entre courbes (calcul intégral)
\usepackage{enumitem}
\usepackage{fancyhdr}
% [most] charge toutes les bibliothèques tcolorbox (listings, minted, raster,
% documentation...) et gonfle inutilement la mémoire TeX sur un document long
% (~300 boîtes) : on ne charge que ce qui est réellement utilisé.
\usepackage[skins,breakable]{tcolorbox}
\usepackage{graphicx}
\usepackage{colortbl}
\usepackage{booktabs}
\usepackage{tabularx}
\usepackage{diagbox} % case d'en-tête diagonale des tableaux à double entrée
% Colonnes extensibles centrée / gauche / droite (C, L, R), souvent utilisées par les modèles
\newcolumntype{C}{>{\centering\arraybackslash}X}
\newcolumntype{L}{>{\raggedright\arraybackslash}X}
\newcolumntype{R}{>{\raggedleft\arraybackslash}X}
\usepackage{array}
\usepackage{multicol}
\usepackage{siunitx}
% parse-numbers=false : accepte aussi \SI{2,5 \times 10^{-3}}{...}
\sisetup{locale=FR,output-decimal-marker={,},group-minimum-digits=4,parse-numbers=false}
\DeclareSIUnit\ohm{\ensuremath{\symup{\Omega}}} % U+2126 absent de la police
\DeclareSIUnit\division{div} % réglages d'oscilloscope (ms/div, V/div)
\DeclareSIUnit\tr{tr} % tours (vitesse de rotation, ex. \SI{1500}{\tr\per\minute})
\DeclareSIUnit\molar{M} % concentration molaire (ex. \SI{3}{\molar}, \SI{1}{\micro\molar})
\DeclareSIUnit\an{an} % année (le modèle confond parfois avec \year, un compteur TeX) — ex. \SI{5}{\kilo\meter\per\an}
\DeclareSIUnit\FCFA{FCFA} % franc CFA (ex. \SI{500}{\FCFA\per\kilo\gram})
\DeclareSIUnit\degreeBrix{\textdegree Brix} % teneur en sucre d'un sirop/jus (réfractométrie)
\DeclareSIUnit\month{mois} % le modèle utilise parfois \month, un compteur TeX, comme unité de durée
\usepackage{qrcode}
% \degree : commande fréquemment utilisée par les modèles pour un angle ou une
% température en dehors de \SI{}{} (non fournie par les paquets chargés).
\providecommand{\degree}{^{\circ}}
% \qed (fin de démonstration), utilisé par les modèles sans amsthm
\providecommand{\qed}{\hfill$\square$}
% \barre : double barre des valeurs interdites dans un tableau de variations (tkz-tab)
\providecommand{\barre}{\ensuremath{\|}}
% environnement proof (amsthm non chargé) utilisé par les modèles
\ifdefined\proof\else\newenvironment{proof}[1][Démonstration]{\par\noindent\textit{#1.}\ }{\qed\par}\fi
\usepackage{lastpage}
\usepackage{hyperref}
\hypersetup{hidelinks}

% ============================================================
% Palette « Pop » OnBuch+ — mêmes hex que la classe P (plus_ui.dart)
% ============================================================
\definecolor{poporange}{HTML}{F59321}
\definecolor{popdark}{HTML}{E07A0C}
\definecolor{popcream}{HTML}{FFF6E8}
\definecolor{popink}{HTML}{1C1714}
\definecolor{poppurple}{HTML}{7A5AE0}
\definecolor{popgreen}{HTML}{2E9E6A}
\definecolor{poppink}{HTML}{E0457B}
\definecolor{popblue}{HTML}{2F7DE1}
\definecolor{popgold}{HTML}{C98A12}
\definecolor{popmuted}{HTML}{8A7F76}
\definecolor{popline}{HTML}{EADFD2}
\definecolor{popgray}{HTML}{8A7F76}
\colorlet{popgrey}{popgray}
% Alias tolérants (le modèle nomme parfois les couleurs différemment)
\colorlet{popwhite}{white}
\colorlet{popblack}{popink}
\colorlet{popred}{poppink}
\colorlet{popyellow}{popgold}
% Teintes claires (fonds de tableaux/encadrés) — le modèle nomme parfois
% les couleurs avec un suffixe L (ex. popblueL) sans les avoir définies.
\colorlet{poporangeL}{poporange!20}
\colorlet{popdarkL}{popdark!20}
\colorlet{poppurpleL}{poppurple!20}
\colorlet{popgreenL}{popgreen!20}
\colorlet{poppinkL}{poppink!20}
\colorlet{popblueL}{popblue!20}
\colorlet{popgoldL}{popgold!20}
\colorlet{popredL}{popred!20}
\colorlet{popgrayL}{popgray!20}
% Teintes claires pour fonds de tableaux / zones
\colorlet{popblueL}{popblue!10!white}
\colorlet{poporangeL}{poporange!14!white}
\colorlet{popgreenL}{popgreen!12!white}
\colorlet{poppurpleL}{poppurple!12!white}
\colorlet{poppinkL}{poppink!12!white}
\colorlet{popgoldL}{popgold!14!white}

\pagecolor{popcream}

% ============================================================
% Typographie
% ============================================================
\defaultfontfeatures{Ligatures=TeX}
\setmainfont{PlusJakartaSans-Medium.ttf}[Path=../../../fonts/,BoldFont=PlusJakartaSans-Bold.ttf]
% Maths en sans-serif (Fira Math) avec chiffres et lettres droites de Plus
% Jakarta, pour que les formules se fondent dans le texte.
\usepackage[math-style=ISO,bold-style=ISO]{unicode-math}
\setmathfont{FiraMath-Regular.otf}[Path=../../../fonts/]
\setmathfont{PlusJakartaSans-Medium.ttf}[Path=../../../fonts/,range={up/num,up/{Latin,latin},\mathup}]
\usepackage{icomma} % virgule décimale sans espace en mode math
% Caractères mathématiques Unicode absents des polices de texte (Plus Jakarta,
% Archivo Black) : rendus par leur équivalent mathématique, en texte comme en
% formule — sinon ils s'affichent en carré vide (ex. « Arithmétique dans ℤ »).
\usepackage{newunicodechar}
\newunicodechar{ℝ}{\ensuremath{\mathbb{R}}}
\newunicodechar{ℤ}{\ensuremath{\mathbb{Z}}}
\newunicodechar{ℂ}{\ensuremath{\mathbb{C}}}
\newunicodechar{ℕ}{\ensuremath{\mathbb{N}}}
\newunicodechar{ℚ}{\ensuremath{\mathbb{Q}}}
\newunicodechar{𝔻}{\ensuremath{\mathbb{D}}}
\newunicodechar{α}{\ensuremath{\alpha}}
\newunicodechar{β}{\ensuremath{\beta}}
\newunicodechar{γ}{\ensuremath{\gamma}}
\newunicodechar{δ}{\ensuremath{\delta}}
\newunicodechar{ε}{\ensuremath{\varepsilon}}
\newunicodechar{θ}{\ensuremath{\theta}}
\newunicodechar{λ}{\ensuremath{\lambda}}
\newunicodechar{μ}{\ensuremath{\mu}}
\newunicodechar{σ}{\ensuremath{\sigma}}
\newunicodechar{τ}{\ensuremath{\tau}}
\newunicodechar{φ}{\ensuremath{\varphi}}
\newunicodechar{ψ}{\ensuremath{\psi}}
\newunicodechar{ω}{\ensuremath{\omega}}
\newunicodechar{Σ}{\ensuremath{\Sigma}}
% majuscules produites par \MakeUppercase dans les titres (α-aminés -> Α)
\newunicodechar{Α}{\ensuremath{\alpha}}
\newunicodechar{Β}{\ensuremath{\beta}}
\newunicodechar{Γ}{\ensuremath{\Gamma}}
\newunicodechar{Δ}{\ensuremath{\Delta}}
\newunicodechar{Ε}{\ensuremath{\varepsilon}}
\newunicodechar{Θ}{\ensuremath{\Theta}}
\newunicodechar{Λ}{\ensuremath{\Lambda}}
\newunicodechar{Μ}{\ensuremath{\mu}}
\newunicodechar{Τ}{\ensuremath{\tau}}
\newunicodechar{Φ}{\ensuremath{\Phi}}
\newunicodechar{Ψ}{\ensuremath{\Psi}}
\newunicodechar{Ω}{\ensuremath{\Omega}}
\newunicodechar{↦}{\ensuremath{\mapsto}}
\newunicodechar{∈}{\ensuremath{\in}}
\newunicodechar{∉}{\ensuremath{\notin}}
\newunicodechar{∧}{\ensuremath{\wedge}}
\newunicodechar{⇔}{\ensuremath{\Leftrightarrow}}
\newunicodechar{⟺}{\ensuremath{\Longleftrightarrow}}
\newunicodechar{⇒}{\ensuremath{\Rightarrow}}
\newunicodechar{⟹}{\ensuremath{\Longrightarrow}}
\newunicodechar{∘}{\ensuremath{\circ}}
\newunicodechar{∪}{\ensuremath{\cup}}
\newunicodechar{∩}{\ensuremath{\cap}}
\newunicodechar{≡}{\ensuremath{\equiv}}
\newunicodechar{⊂}{\ensuremath{\subset}}
\newunicodechar{⊥}{\ensuremath{\perp}}
\newunicodechar{∃}{\ensuremath{\exists}}
\newunicodechar{∀}{\ensuremath{\forall}}
\newunicodechar{∛}{\ensuremath{\sqrt[3]{\,}}}
\newunicodechar{∜}{\ensuremath{\sqrt[4]{\,}}}
\newunicodechar{⃗}{\ensuremath{{}^{\to}}}
\newunicodechar{ⁿ}{\ifmmode^{n}\else\textsuperscript{\ensuremath{n}}\fi}
\newunicodechar{ˣ}{\ifmmode^{x}\else\textsuperscript{\ensuremath{x}}\fi}
\newunicodechar{ᵇ}{\ifmmode^{b}\else\textsuperscript{\ensuremath{b}}\fi}
\newunicodechar{ʳ}{\ifmmode^{r}\else\textsuperscript{\ensuremath{r}}\fi}
\newunicodechar{ᵘ}{\ifmmode^{u}\else\textsuperscript{\ensuremath{u}}\fi}
\newunicodechar{ʸ}{\ifmmode^{y}\else\textsuperscript{\ensuremath{y}}\fi}
\newunicodechar{ᵃ}{\ifmmode^{a}\else\textsuperscript{\ensuremath{a}}\fi}
\newunicodechar{ᵉ}{\ifmmode^{e}\else\textsuperscript{\ensuremath{e}}\fi}
\newunicodechar{ᵏ}{\ifmmode^{k}\else\textsuperscript{\ensuremath{k}}\fi}
\newunicodechar{ᵖ}{\ifmmode^{p}\else\textsuperscript{\ensuremath{p}}\fi}
\newunicodechar{ᴺ}{\ifmmode^{N}\else\textsuperscript{\ensuremath{N}}\fi}
\newunicodechar{ᶜ}{\ifmmode^{c}\else\textsuperscript{\ensuremath{c}}\fi}
\newunicodechar{ʰ}{\ifmmode^{h}\else\textsuperscript{\ensuremath{h}}\fi}
\newunicodechar{ᵠ}{\ifmmode^{q}\else\textsuperscript{\ensuremath{q}}\fi}
\newunicodechar{ₙ}{\ifmmode_{n}\else\textsubscript{\ensuremath{n}}\fi}
\newunicodechar{ₐ}{\ifmmode_{a}\else\textsubscript{\ensuremath{a}}\fi}
\newunicodechar{ᵢ}{\ifmmode_{i}\else\textsubscript{\ensuremath{i}}\fi}
\newunicodechar{ᵣ}{\ifmmode_{r}\else\textsubscript{\ensuremath{r}}\fi}
\newunicodechar{ₖ}{\ifmmode_{k}\else\textsubscript{\ensuremath{k}}\fi}
\newunicodechar{ᵦ}{\ifmmode_{\beta}\else\textsubscript{\ensuremath{\beta}}\fi}
\newunicodechar{ₓ}{\ifmmode_{x}\else\textsubscript{\ensuremath{x}}\fi}
\newfontfamily\popdisplay{ArchivoBlack-Regular.ttf}[Path=../../../fonts/]
\newfontfamily\poplogo{SpaceGrotesk-Bold.ttf}[Path=../../../fonts/]
\newfontfamily\popsemi{PlusJakartaSans-SemiBold.ttf}[Path=../../../fonts/]
\newfontfamily\popxbold{PlusJakartaSans-ExtraBold.ttf}[Path=../../../fonts/]
\newfontfamily\popxboldls{PlusJakartaSans-ExtraBold.ttf}[Path=../../../fonts/,LetterSpace=3.0]

\setlist[enumerate]{leftmargin=1.7em,itemsep=5pt,topsep=4pt}
\setlist[itemize]{leftmargin=1.7em,itemsep=4pt,topsep=4pt,label={\color{poporange}\textbullet}}
\setlength{\parindent}{0pt}
\setlength{\parskip}{5pt}
\renewcommand{\arraystretch}{1.25}

% pgfplots : style Pop par défaut
\pgfplotsset{
  every axis/.append style={axis line style={popink,line width=1pt},tick style={popink},
    grid style={popline,line width=0.5pt},label style={font=\small},tick label style={font=\footnotesize}},
  popcurve/.style={poporange,line width=1.8pt,no markers,smooth},
}
% Même style disponible dans un \draw TikZ simple (hors pgfplots)
\tikzset{popcurve/.style={poporange,line width=1.8pt}}
\tikzset{popgrid/.style={popline,very thin}}
\colorlet{popcurve}{poporange} % le nom du style est parfois utilisé comme couleur

% ============================================================
% Pill / squiggle
% ============================================================
\newcommand{\poppill}[3]{%
  \tikz[baseline=-0.55ex]\node[draw=popink,line width=1.2pt,rounded corners=2.6pt,%
    fill=#2,inner xsep=6.5pt,inner ysep=3pt]{%
    \popxboldls\fontsize{7}{7}\selectfont\color{#3}\MakeUppercase{#1}};%
}
\newcommand{\popsquiggle}[1]{%
  \tikz[baseline=-0.4ex]{
    \draw[#1,line width=2.6pt,line cap=round]
      plot [smooth] coordinates {(0,0.09) (0.34,0.01) (0.68,0.21) (1.02,0.01) (1.36,0.10)};
  }%
}
% Mot-clé surligné (marqueur orange sous le texte)
\newcommand{\cle}[1]{{\popxbold\color{popdark}#1}}

% ============================================================
% En-tête / pied
% ============================================================
\pagestyle{fancy}
\fancyhf{}
\renewcommand{\headrulewidth}{0pt}
\renewcommand{\footrulewidth}{0pt}
\lhead{\raisebox{-0.15ex}{\poplogo\fontsize{13}{13}\selectfont\color{popink}OnBuch\color{poporange}+}}
\rhead{\poppill{\DOCMATIERE\ \textbullet\ \DOCNIVEAU\ \textbullet\ Leçon \DOCLECON}{white}{popink}}
\lfoot{\popxbold\fontsize{7.6}{7.6}\selectfont\color{popmuted}\MakeUppercase{Cours signé OnBuch+}\ \textbullet\ {\popsemi\DOCTITRE}}
\rfoot{\tikz[baseline=-0.6ex]\node[rounded corners=8.5pt,fill=popink,minimum height=17pt,minimum width=17pt,inner xsep=5pt,inner sep=0]{\popxbold\fontsize{8}{8}\selectfont\color{popcream}\thepage\,/\,\pageref*{LastPage}};}

% ============================================================
% Titres
% ============================================================
\newcommand{\chaptitre}[2]{%
  \begin{tcolorbox}[enhanced,colback=poporange,colframe=popink,boxrule=2.6pt,arc=11pt,
      drop shadow=popink,left=15pt,right=15pt,top=12pt,bottom=11pt,before skip=3pt,after skip=18pt]
    \poppill{#2}{popink}{popcream}\par\vspace{9pt}
    {\raggedright\hyphenpenalty=10000\exhyphenpenalty=10000\popdisplay\fontsize{22}{24}\selectfont\color{popcream}\MakeUppercase{#1}\par}
  \end{tcolorbox}
}
% Section numérotée : pastille orange avec le numéro + titre Archivo
\newcounter{popsec}
\newcommand{\coursec}[1]{%
  \stepcounter{popsec}\setcounter{popsub}{0}%
  \par\vspace{16pt}\noindent
  \begin{minipage}{\linewidth}
  \tikz[baseline=-0.7ex]\node[draw=popink,line width=1.6pt,fill=poporange,rounded corners=4pt,minimum size=22pt,inner sep=0]{\popdisplay\fontsize{12}{12}\selectfont\color{popcream}\thepopsec};%
  \hspace{8pt}{\popdisplay\fontsize{15}{17}\selectfont\color{popink}\MakeUppercase{#1}}\par
  \vspace{4pt}\noindent{\color{poporange}\rule{36pt}{3.6pt}}
  \end{minipage}\par\nopagebreak\vspace{8pt}
}
\newcounter{popsub}
\newcommand{\courssub}[1]{%
  \stepcounter{popsub}%
  \par\vspace{9pt}\noindent{\popxbold\fontsize{11.5}{13}\selectfont\color{popink}{\color{poporange}\thepopsec.\thepopsub}\hspace{6pt}#1}\par\nopagebreak\vspace{4pt}
}

% ============================================================
% Boîtes pédagogiques — même recette : fond blanc, bordure épaisse colorée,
% ombre dure encre, badge plein. Titre optionnel [#1].
% ============================================================
\tcbset{popbase/.style={breakable,enhanced,colback=white,boxrule=2.3pt,arc=9pt,
  shadow={3pt}{-3pt}{0pt}{popink},left=14pt,right=14pt,top=11pt,bottom=11pt,before skip=11pt,after skip=15pt,
  fontupper=\popsemi\fontsize{10.6}{14.5}\selectfont\color{popink},
  fontlower=\popsemi\fontsize{10.6}{14.5}\selectfont\color{popink}}}
% Coche dessinée (le glyphe ✓ n'existe pas dans les polices du document)
\newcommand{\popcheck}[1]{\tikz[baseline=-0.5ex]\draw[#1,line width=1.6pt,line cap=round,line join=round](-0.13,0.0)--(-0.04,-0.09)--(0.13,0.1);}
\providecommand{\coche}{\popcheck{popgreen}}
\providecommand{\resultat}[1]{\par\smallskip\noindent\fcolorbox{popgreen}{popgreenL}{\parbox{\dimexpr\linewidth-2\fboxsep-2\fboxrule\relax}{\textbf{Résultat.} #1}}\par\smallskip}
\newcommand{\popboxhead}[3]{\poppill{#1}{#2}{white}\ifx\relax#3\relax\else\hspace{6pt}{\popxbold\fontsize{10.6}{12}\selectfont\color{popink}#3}\fi\par\vspace{7pt}}

\newtcolorbox{aretenir}[1][]{popbase,colframe=popgreen,before upper={\popboxhead{À retenir}{popgreen}{#1}}}
% Texte en espagnol (document de lecture, dialogue, poème, extrait) : cadre
% d'étude. Un titre optionnel (source, auteur) s'affiche dans l'en-tête.
\newtcolorbox{textebox}[1][]{popbase,colframe=popblue,before upper={\popboxhead{Texto}{popblue}{#1}\begin{otherlanguage}{spanish}},after upper={\end{otherlanguage}}}
% Marques ✓ / ✗ (forme correcte / fautive) souvent utilisées par les modèles
\providecommand{\cross}{\textcolor{poppink}{\ensuremath{\times}}}
\providecommand{\xmark}{\cross}\providecommand{\crossmark}{\cross}\providecommand{\cmark}{\ensuremath{\checkmark}}
% Ligne de réponse / justification à compléter (inventée par les modèles)
\providecommand{\justification}[1]{\par\noindent\textit{Justification~:} \dotfill\par}
% \textipa (package tipa non chargé) : la police ne couvre pas l'API -> simple passage
\providecommand{\textipa}[1]{#1}
% Mot ou phrase en espagnol à l'intérieur d'un paragraphe français
\newcommand{\esp}[1]{\foreignlanguage{spanish}{\emph{#1}}}
\newtcolorbox{exemplebox}[1][]{popbase,colframe=poporange,before upper={\popboxhead{Exemple}{poporange}{#1}}}
\newtcolorbox{definition}[1][]{popbase,colframe=popblue,before upper={\popboxhead{Définition}{popblue}{#1}}}
\newtcolorbox{propriete}[1][]{popbase,colframe=poppurple,before upper={\popboxhead{Propriété}{poppurple}{#1}}}
\newtcolorbox{methode}[1][]{popbase,colframe=popgold,before upper={\popboxhead{Méthode}{popgold}{#1}}}
\newtcolorbox{attention}[1][]{popbase,colframe=poppink,before upper={\popboxhead{Attention}{poppink}{#1}}}
\newtcolorbox{experience}[1][]{popbase,colframe=popblue,colback=popblueL,before upper={\popboxhead{Expérience}{popblue}{#1}}}
% « Le savais-tu ? » : carte violette pleine (contexte, culture, Cameroun)
\newtcolorbox{savaistu}[1][]{popbase,colback=poppurple,colframe=popink,
  fontupper=\popsemi\fontsize{10.4}{14.2}\selectfont\color{white},
  before upper={\poppill{Le savais-tu\,?}{white}{poppurple}\ifx\relax#1\relax\else\hspace{6pt}{\popxbold\color{white}#1}\fi\par\vspace{7pt}}}
% Exercice résolu : énoncé en haut, \tcblower puis la solution sur fond crème
\newcounter{popexo}\newcounter{popexr}
\newtcolorbox{exoresolu}[1][]{popbase,colframe=poporange,colbacklower=popcream,
  segmentation style={popink,line width=1.2pt,dashed},
  before upper={\stepcounter{popexr}\popboxhead{Exercice résolu \thepopexr}{poporange}{#1}},
  before lower={\poppill{Solution}{popink}{popcream}\par\vspace{6pt}}}
% Exercice (non résolu en place) — la correction est dans la partie Corrigés
\newtcolorbox{exercice}[1][]{popbase,colframe=popink,
  before upper={\stepcounter{popexo}\popboxhead{Exercice \thepopexo}{popink}{#1}}}
% Corrigé
\newtcolorbox{corrige}[1][]{popbase,colframe=popgreen,colback=popgreenL,
  before upper={\popboxhead{Corrigé}{popgreen}{#1}}}
% Objectifs / prérequis / bilan
\newtcolorbox{objectifs}{popbase,colframe=popink,colback=poporangeL,
  before upper={\popboxhead{Objectifs de la leçon}{popink}{}}}
\newtcolorbox{prerequis}{popbase,colframe=popink,colback=popblueL,
  before upper={\popboxhead{Prérequis}{popblue}{}}}
\newtcolorbox{bilan}{popbase,colframe=popink,colback=white,boxrule=2.8pt,
  before upper={\popboxhead{Fiche bilan}{poporange}{L'essentiel à connaître pour l'examen}}}
% Figure encadrée avec légende
\newtcolorbox{popfigure}[1][]{enhanced,breakable=false,colback=white,colframe=popline,boxrule=1.4pt,arc=8pt,
  left=10pt,right=10pt,top=10pt,bottom=8pt,before skip=10pt,after skip=12pt,halign=center,
  lower separated=false,halign lower=center,
  fontlower=\popsemi\fontsize{9}{11.5}\selectfont\color{popmuted},
  #1}
\newcommand{\legende}[1]{\par\vspace{6pt}{\popsemi\fontsize{9}{11.5}\selectfont\color{popmuted}#1}}

% ============================================================
% Signature OnBuch+
% ============================================================
\newcommand{\poplogoinline}[1]{{\poplogo\fontsize{#1}{#1}\selectfont\color{popink}OnBuch\color{poporange}+}}
\newcommand{\signatureonbuch}{%
  \begin{tcolorbox}[enhanced,colback=popink,colframe=popink,boxrule=2.2pt,arc=10pt,
      drop shadow=poporange,left=14pt,right=14pt,top=10pt,bottom=10pt,before skip=0pt,after skip=0pt]
    \tikz[baseline=-0.7ex]\node[circle,fill=popgreen,draw=popcream,line width=1.3pt,minimum size=22pt,inner sep=0]{\popcheck{white}};%
    \hspace{9pt}\begin{minipage}[c]{0.62\linewidth}
      {\popxbold\fontsize{12}{13}\selectfont\color{popcream}Signé\ \textbullet\ {\poplogo OnBuch\color{poporange}+}}\par\vspace{2pt}
      {\raggedright\popsemi\fontsize{8}{10.5}\selectfont\color{popline}\MakeUppercase{Équipe pédagogique OnBuch+}\\\MakeUppercase{Conforme au programme officiel MINESEC}\par}
    \end{minipage}\hfill
    {\poplogo\fontsize{22}{22}\selectfont\color{popcream}OnBuch\color{poporange}+}
  \end{tcolorbox}%
}

% ============================================================
% Page de garde
% #1 titre · #2 matière · #3 niveau · #4 module · #5 n° leçon · #6 durée ·
% #7 description / objectifs (texte ou itemize)
% ============================================================
\newcommand{\pagegarde}[7]{%
  \thispagestyle{empty}
  \newgeometry{a4paper,margin=1.7cm,top=1.6cm,bottom=1.4cm}%
  \begin{tikzpicture}[remember picture,overlay]
    \fill[poporange!18] ([xshift=-3.2cm,yshift=-3.6cm]current page.north east) circle (4.6cm);
    \fill[poppurple!14] ([xshift=2.4cm,yshift=7.6cm]current page.south west) circle (3.8cm);
  \end{tikzpicture}%
  {\poplogo\fontsize{30}{30}\selectfont\color{popink}OnBuch\color{poporange}+}\hfill
  \raisebox{0.6ex}{\poppill{Cours signé \textbullet\ Premium}{popink}{popcream}}\par
  \vspace{1pt}\popsquiggle{poppurple}\par
  \vspace{8pt}
  \begin{tcolorbox}[enhanced,colback=poporange,colframe=popink,boxrule=2.8pt,arc=13pt,
      drop shadow=popink,left=18pt,right=18pt,top=12pt,bottom=13pt,after skip=10pt]
    \poppill{#2 \textbullet\ #3}{popink}{popcream}\hspace{5pt}\poppill{#4}{white}{popink}\par\vspace{8pt}
    {\popdisplay\fontsize{14}{14}\selectfont\color{popink}LEÇON #5}\par\vspace{4pt}
    {\raggedright\hyphenpenalty=10000\exhyphenpenalty=10000\popdisplay\fontsize{25}{27}\selectfont\color{popcream}\MakeUppercase{#1}\par}
    \vspace{8pt}
    {\popxbold\fontsize{9.5}{9.5}\selectfont\color{popink}\MakeUppercase{Durée conseillée : #6}}
  \end{tcolorbox}
  \begin{tcolorbox}[enhanced,colback=white,colframe=popink,boxrule=2.2pt,arc=10pt,
      drop shadow=popink,left=16pt,right=16pt,top=10pt,bottom=10pt,after skip=8pt]
    \poppill{Dans cette leçon}{poppurple}{white}\par\vspace{5pt}
    \popsemi\fontsize{9.3}{12.6}\selectfont\color{popink}#7
  \end{tcolorbox}
  \noindent\begin{minipage}[t]{0.487\linewidth}
    \begin{tcolorbox}[enhanced,colback=popgreen,colframe=popink,boxrule=2.2pt,arc=10pt,
        drop shadow=popink,left=11pt,right=11pt,top=10pt,bottom=10pt,height=3.1cm,valign=center]
      \tcbox[colback=white,colframe=popink,boxrule=1.3pt,arc=5pt,boxsep=0pt,nobeforeafter,
        left=3pt,right=3pt,top=3pt,bottom=3pt,valign=center]{\qrcode[height=1.9cm]{https://play.google.com/store/apps/details?id=com.luvvix.onbuch&hl=fr}}%
      \hspace{8pt}\begin{minipage}[c]{0.5\linewidth}
        {\popdisplay\fontsize{11}{12.5}\selectfont\color{white}\MakeUppercase{Télécharge OnBuch}}\par\vspace{4pt}
        {\raggedright\popsemi\fontsize{8.4}{11}\selectfont\color{white}Gratuit \textbullet\ Google Play\\Tuteur IA, annales, concours, résultats\par}
      \end{minipage}
    \end{tcolorbox}
  \end{minipage}\hfill
  \begin{minipage}[t]{0.487\linewidth}
    \begin{tcolorbox}[enhanced,colback=poppurple,colframe=popink,boxrule=2.2pt,arc=10pt,
        drop shadow=popink,left=11pt,right=11pt,top=10pt,bottom=10pt,height=3.1cm,valign=center]
      \tikz[baseline=-0.7ex]\node[circle,fill=white,draw=popink,line width=1.3pt,minimum size=1.6cm,inner sep=0]{\popdisplay\fontsize{24}{24}\selectfont\color{poppurple}+};%
      \hspace{8pt}\begin{minipage}[c]{0.6\linewidth}
        {\popdisplay\fontsize{11}{12.5}\selectfont\color{white}\MakeUppercase{Passe à OnBuch+}}\par\vspace{4pt}
        {\raggedright\popsemi\fontsize{8.4}{11}\selectfont\color{white}Tous les cours signés, Tuteur IA illimité, fiches PDF — onglet OnBuch+ de l'app\par}
      \end{minipage}
    \end{tcolorbox}
  \end{minipage}
  \vspace{8pt}
  \popcontenu
  \vfill
  \signatureonbuch
  \restoregeometry
  \newpage
}

% Rangée « Ce cours contient » (page de garde)
\newcommand{\poptile}[3]{% #1 couleur #2 gros texte #3 légende
  \begin{tcolorbox}[enhanced,colback=#1,colframe=popink,boxrule=2pt,arc=9pt,shadow={2.5pt}{-2.5pt}{0pt}{popink},
      left=6pt,right=6pt,top=9pt,bottom=9pt,halign=center,valign=center,height=1.75cm,width=\linewidth,nobeforeafter]
    {\popdisplay\fontsize{12.5}{13}\selectfont\color{white}\MakeUppercase{#2}}\par\vspace{3pt}
    {\centering\popsemi\fontsize{7.8}{9.5}\selectfont\color{white}#3\par}
  \end{tcolorbox}}
\newcommand{\popcontenu}{%
  \par\noindent{\popxboldls\fontsize{7.5}{7.5}\selectfont\color{popmuted}CE COURS SIGNÉ CONTIENT}\par\vspace{7pt}
  \noindent\begin{minipage}[t]{0.235\linewidth}\poptile{popblue}{Cours}{complet, illustré, pas à pas}\end{minipage}\hfill
  \begin{minipage}[t]{0.235\linewidth}\poptile{poporange}{Méthodes}{et exercices résolus}\end{minipage}\hfill
  \begin{minipage}[t]{0.235\linewidth}\poptile{poppink}{Exercices}{type examen + corrigés}\end{minipage}\hfill
  \begin{minipage}[t]{0.235\linewidth}\poptile{popgreen}{Fiche bilan}{l'essentiel à retenir}\end{minipage}\par}

% Sommaire
\newtcolorbox{sommaire}{enhanced,colback=white,colframe=popink,boxrule=2.2pt,arc=10pt,
  drop shadow=popink,left=18pt,right=18pt,top=13pt,bottom=13pt,before skip=6pt,after skip=20pt,
  before upper={\poppill{Sommaire}{poppurple}{white}\par\vspace{9pt}\popsemi\fontsize{10.6}{14.5}\selectfont\color{popink}}}

% Signature de fin de cours — poussée tout en bas de la dernière page
\newcommand{\findecours}{%
  \par\vfill\nopagebreak
  \begin{center}\popsquiggle{poporange}\end{center}\vspace{4pt}
  \signatureonbuch
}
\AtBeginDocument{\def\llbracket{\mathopen{[\mkern-3.5mu[}}\def\rrbracket{\mathclose{]\mkern-3.5mu]}}} % intervalles d'entiers (glyphe ⟦ absent de la police)
% Glyphes absents de Fira Math : redéfinis à partir de glyphes disponibles
\AtBeginDocument{%
  \def\setminus{\mathbin{\backslash}}\let\smallsetminus\setminus
  \def\vdots{\mathord{\vbox{\baselineskip3.2pt\lineskiplimit0pt\kern4pt\hbox{.}\hbox{.}\hbox{.}}}}%
  \def\ddots{\mathinner{\mkern1mu\raise6.5pt\hbox{.}\mkern2mu\raise3.6pt\hbox{.}\mkern2mu\raise0.7pt\hbox{.}\mkern1mu}}%
  \def\triangle{\mathord{\scalebox{0.9}{$\symup{\Delta}$}}}%
  \def\nabla{\mathord{\raisebox{\depth}{\scalebox{1}[-1]{$\symup{\Delta}$}}}}%
  \def\checkmark{\popcheck{popdark}}%
  \def\star{\mathbin{\ast}}\def\bigstar{\mathord{\ast}}%
  \def\diamond{\mathbin{\scalebox{0.6}{$\Diamond$}}}%
  \def\bigcup{\mathop{\mathchoice{\raisebox{-0.3em}{\scalebox{1.7}{$\cup$}}}{\scalebox{1.25}{$\cup$}}{\cup}{\cup}}\displaylimits}%
  \def\bigcap{\mathop{\mathchoice{\raisebox{-0.3em}{\scalebox{1.7}{$\cap$}}}{\scalebox{1.25}{$\cap$}}{\cap}{\cap}}\displaylimits}%
  \def\lesseqqgtr{\mathrel{\lessgtr}}\def\bigoplus{\mathop{\oplus}\displaylimits}%
}
\providecommand{\ssi}{si et seulement si} % abréviation fréquente
\AtBeginDocument{\providecommand{\wideparen}{\overparen}} % arc de cercle

\begin{document}

\pagegarde{Lectura y comentario : derechos de las mujeres y de las minorías}{Espagnol}{Tle A}{Módulo 1 : Vie familiale et intégration sociale}{E03}{2 h}{Cette leçon propose la lecture commentée d'un texte original sur les droits des femmes et des minorités. Elle fournit le lexique thématique indispensable, la méthode complète du commentaire de texte au Bac et une production guidée de revendication.
\vspace{4pt}
\begin{multicols}{2}\small
\begin{itemize}
\item À la fin de cette leçon, je sais lire et comprendre globalement un texte espagnol sur les droits des femmes et des minorités
\item je sais dégager les idées directrices et le point de vue de l'auteur
\item je sais utiliser le lexique des droits humains (derechos humanos, minoría, discriminación, sufragio, igualdad de oportunidades, ley, reivindicar) dans un écrit
\item je sais répondre à des questions de compréhension en phrases complètes et justifiées
\item je sais exprimer mon opinion et une revendication avec les structures adaptées (me parece que, es necesario que + subjuntivo, hay que, reivindicamos)
\item je sais construire un commentaire de texte en trois parties (introduction, développement, conclusion)
\end{itemize}\end{multicols}}

\begin{sommaire}
\begin{multicols}{2}
\begin{enumerate}
\item Texte support et première lecture
\item Compréhension détaillée du texte
\item Lexique thématique : los derechos humanos
\item Exprimer l'opinion et la revendication
\item Méthode du commentaire de texte
\item Production guidée : un texto reivindicativo
\item Activité d'intégration
\item Méthodes et exercices résolus
\item Exercices
\item Corrigés des exercices
\item Fiche bilan
\end{enumerate}
\end{multicols}
\end{sommaire}

\begin{objectifs}
\begin{itemize}
\item À la fin de cette leçon, je sais lire et comprendre globalement un texte espagnol sur les droits des femmes et des minorités
\item je sais dégager les idées directrices et le point de vue de l'auteur
\item je sais utiliser le lexique des droits humains (derechos humanos, minoría, discriminación, sufragio, igualdad de oportunidades, ley, reivindicar) dans un écrit
\item je sais répondre à des questions de compréhension en phrases complètes et justifiées
\item je sais exprimer mon opinion et une revendication avec les structures adaptées (me parece que, es necesario que + subjuntivo, hay que, reivindicamos)
\item je sais construire un commentaire de texte en trois parties (introduction, développement, conclusion)
\item je sais rédiger un court texte engagé (lettre, article, discours) en réemployant le lexique de la leçon
\end{itemize}
\end{objectifs}

\begin{prerequis}
\begin{itemize}
\item Le présent de l'indicatif et le présent du subjonctif (Première)
\item Les structures d'opinion de base : pienso que, creo que, en mi opinión
\item Le lexique de la famille et de la société (Module 1, leçons E01-E02)
\item La technique de lecture d'un texte : titre, source, repérage des connecteurs
\end{itemize}
\end{prerequis}

\newpage

\coursec{Texte support et première lecture}

En mars 2018, des milliers de femmes sont descendues dans les rues de Madrid et de Mexico pour réclamer l'égalité salariale : pancartes, slogans, grèves symboliques. Au même moment, au Cameroun, des associations défendent le droit des filles à l'éducation dans les régions du Nord et celui des populations minoritaires à préserver leurs langues et leurs terres. Partout, la même question se pose : comment lire, comprendre et commenter un texte qui \cle{revendique} des droits ?

Au Baccalauréat, l'épreuve d'espagnol vous demandera précisément cela : lire un texte de presse ou un texte littéraire, en dégager les idées directrices, identifier le point de vue de l'auteur et exprimer le vôtre. Cette leçon vous donne les clés pour analyser un texte engagé en espagnol. La problématique qui nous guidera est la suivante : \textbf{comment un article de presse défend-il les droits des femmes et des minorités, et quels outils de lecture nous permettent de le comprendre ?}

\courssub{Présentation du texte}

Avant même de lire, un bon lecteur observe le « paratexte » : le titre, le nom de l'auteur, la source, la date. Ces éléments donnent des indices précieux sur le contenu et l'intention du texte.

\begin{definition}[Fiche d'identité du texte]
\begin{itemize}
\item \textbf{Titre :} \esp{Las mujeres y las minorías: una lucha por la igualdad} (« Les femmes et les minorités : une lutte pour l'égalité »).
\item \textbf{Auteur :} Carmen Ndong, journaliste (nom fictif créé pour ce cours).
\item \textbf{Source :} article de presse publié dans le journal imaginaire \emph{El Día de Malabo} (Guinée équatoriale, pays hispanophone d'Afrique centrale).
\item \textbf{Genre :} article d'opinion (\esp{artículo de opinión}).
\item \textbf{Thème annoncé :} les droits des femmes et des minorités, la lutte contre la discrimination.
\end{itemize}
\end{definition}

\begin{savaistu}[Un pays hispanophone en Afrique]
La Guinée équatoriale, voisine du Cameroun, est le seul pays d'Afrique subsaharienne dont l'espagnol est langue officielle. Un élève camerounais qui maîtrise l'espagnol peut donc communiquer à Malabo ou Bata sans interprète ! C'est aussi pour cela que les textes de presse hispanophones parlent parfois de réalités très proches des nôtres.
\end{savaistu}

\courssub{Première lecture : la méthode}

On ne lit jamais un texte de presse « à froid », ligne par ligne, dès la première rencontre. La première lecture est une lecture \cle{globale} : elle vise à saisir le sens général, pas chaque mot.

\begin{methode}[Réussir une première lecture efficace]
\begin{enumerate}
\item \textbf{Lire le titre} et formuler une hypothèse : de quoi va parler le texte ?
\item \textbf{Lire la source et l'auteur} : article de presse ? poème ? discours ? Cela change la manière d'analyser.
\item \textbf{Lire le premier paragraphe} : il contient presque toujours le sujet et la thèse de l'auteur.
\item \textbf{Lire le dernier paragraphe} : il contient la conclusion, souvent l'appel ou la proposition de l'auteur.
\item \textbf{Lire ensuite le texte entier} en silence, sans dictionnaire, en soulignant les mots transparents.
\item \textbf{Reformuler en une phrase} : « Ce texte parle de... »
\end{enumerate}
\end{methode}

En classe, cette lecture globale se fait en deux temps : d'abord une \textbf{lecture silencieuse} individuelle (deux minutes), puis une \textbf{lecture à voix haute} par l'enseignant, qui modèle la prononciation et l'intonation. Pendant la lecture silencieuse, soulignez les mots que vous reconnaissez sans dictionnaire.

\courssub{Le texte}

\begin{textebox}[Las mujeres y las minorías: una lucha por la igualdad (Carmen Ndong, \emph{El Día de Malabo})]
A lo largo de la historia, las mujeres y las minorías han tenido que reivindicar derechos que los demás consideraban naturales. El sufragio femenino, conquistado en España en 1931, es un ejemplo histórico de esta lucha. Hoy, la discriminación sigue presente en el trabajo, en la educación y en la vida política. Sin embargo, las leyes avanzan : la igualdad de oportunidades es ya un principio reconocido por los derechos humanos. Pero una ley no basta : hace falta cambiar las mentalidades. Por eso, las asociaciones siguen reivindicando una sociedad más justa, donde nadie sea excluido por su sexo, su origen o su lengua.
\end{textebox}

\textbf{Glossaire (mots difficiles) :}

\begin{center}
\begin{tabularx}{\linewidth}{|X|X|}\hline
\rowcolor{popblueL} \textbf{Español} & \textbf{Français} \\ \hline
\esp{a lo largo de la historia} & au fil de l'histoire \\ \hline
\esp{la lucha} & la lutte \\ \hline
\esp{reivindicar} & revendiquer, réclamer \\ \hline
\esp{el sufragio femenino} & le droit de vote des femmes \\ \hline
\esp{conquistado} & conquis, obtenu (participe passé de \esp{conquistar}) \\ \hline
\esp{hace falta} & il faut, il est nécessaire de \\ \hline
\esp{las mentalidades} & les mentalités \\ \hline
\esp{excluido} & exclu \\ \hline
\esp{la lengua} & la langue \\ \hline
\end{tabularx}
\end{center}

\textbf{Traduction du texte :} « Au fil de l'histoire, les femmes et les minorités ont dû revendiquer des droits que les autres considéraient comme naturels. Le droit de vote des femmes, conquis en Espagne en 1931, est un exemple historique de cette lutte. Aujourd'hui, la discrimination reste présente au travail, dans l'éducation et dans la vie politique. Cependant, les lois avancent : l'égalité des chances est déjà un principe reconnu par les droits humains. Mais une loi ne suffit pas : il faut changer les mentalités. C'est pourquoi les associations continuent de revendiquer une société plus juste, où personne ne soit exclu à cause de son sexe, de son origine ou de sa langue. »

\courssub{Repérage des mots transparents}

Un \cle{mot transparent} est un mot dont la forme est très proche du français : on le comprend sans dictionnaire. C'est le premier allié du lecteur francophone. Dans notre texte, repérez et soulignez :

\begin{center}
\begin{tabularx}{\linewidth}{|X|X|X|}\hline
\rowcolor{popblueL} \textbf{Español} & \textbf{Français} & \textbf{Remarque} \\ \hline
\esp{la discriminación} & la discrimination & même racine latine \\ \hline
\esp{la igualdad} & l'égalité & attention : \esp{igualdad} et non « igualité » \\ \hline
\esp{las oportunidades} & les opportunités, les chances & \esp{igualdad de oportunidades} = égalité des chances \\ \hline
\esp{la minoría} & la minorité & accent sur le \emph{í} \\ \hline
\esp{la educación} & l'éducation & les mots en \esp{-ción} = français en « -tion » \\ \hline
\esp{la política} & la politique & \esp{la vida política} = la vie politique \\ \hline
\esp{natural} & naturel & adjectif identique \\ \hline
\esp{el principio} & le principe & ne pas confondre avec \esp{el príncipe} (le prince) \\ \hline
\end{tabularx}
\end{center}

\begin{propriete}[La famille des mots en -ción]
Les noms français en « -tion » correspondent presque toujours à des noms espagnols en \esp{-ción}, tous \textbf{féminins} : \esp{la discriminación}, \esp{la educación}, \esp{la asociación}, \esp{la población}, \esp{la integración}. Ils prennent un accent écrit sur le \emph{ó} au singulier, qui disparaît au pluriel : \esp{las discriminaciones}.
\end{propriete}

\begin{attention}[Faux amis et pièges de traduction]
\begin{itemize}
\item \esp{la lucha} signifie « la lutte » (combat pour une cause), pas « la luxure » ni « la honte ».
\item \esp{reivindicar} = revendiquer ; ne pas traduire par « revendiquer » au sens de « se vanter ».
\item \esp{los derechos humanos} : on dit « droits de l'homme » ou « droits humains » en français, jamais « droits des humains ».
\item \esp{excluido por su sexo} : \esp{sexo} = sexe (masculin/féminin) ; attention au registre soutenu du texte.
\item \esp{hace falta} est impersonnel : \esp{Hace falta cambiar} = « il faut changer ». Ne dites jamais \esp{hace falta que cambiar}.
\end{itemize}
\end{attention}

\courssub{Hypothèses de sens et vérification globale}

À partir du titre seul — \esp{Las mujeres y las minorías: una lucha por la igualdad} — on peut formuler des hypothèses avant la lecture :

\begin{itemize}
\item Le texte parlera probablement de \textbf{deux groupes} : les femmes et les minorités.
\item Le mot \esp{lucha} annonce un texte \textbf{engagé}, qui parle d'un combat, pas d'une simple description.
\item \esp{Igualdad} annonce le thème central : l'égalité.
\end{itemize}

Après la lecture, on vérifie ces hypothèses avec les trois questions fondamentales du commentaire :

\begin{exemplebox}[Les trois questions de la lecture globale]
\begin{itemize}
\item \textbf{De quoi parle le texte ?} De la lutte des femmes et des minorités pour l'égalité des droits.
\item \textbf{Qui parle ?} Une journaliste, Carmen Ndong, qui donne son opinion : c'est un article d'opinion, pas un simple reportage.
\item \textbf{À qui s'adresse-t-elle ?} Aux lecteurs du journal, c'est-à-dire au grand public, pour le convaincre et le sensibiliser.
\end{itemize}
\end{exemplebox}

\begin{exoresolu}[Vérifier sa compréhension globale]
\textbf{Énoncé.} Répondez en français, en vous appuyant sur le texte : 1) Quel exemple historique de conquête des droits l'auteur cite-t-elle ? 2) Dans quels trois domaines la discrimination est-elle encore présente selon elle ? 3) Que faut-il changer, d'après la conclusion ?
\tcblower
\textbf{Solution détaillée.}
\begin{enumerate}
\item L'auteur cite \esp{el sufragio femenino, conquistado en España en 1931} : le droit de vote des femmes, obtenu en Espagne en 1931. C'est l'exemple historique de la lutte.
\item La discrimination reste présente \esp{en el trabajo, en la educación y en la vida política} : au travail, dans l'éducation et dans la vie politique. On repère l'énumération de trois éléments, procédé typique de l'argumentation.
\item D'après la conclusion, \esp{hace falta cambiar las mentalidades} : il faut changer les mentalités, car une loi ne suffit pas (\esp{una ley no basta}).
\end{enumerate}
\textbf{Méthode :} pour chaque réponse, citez d'abord la phrase espagnole exacte du texte, puis traduisez-la. C'est la démarche exigée au Baccalauréat.
\end{exoresolu}

\courssub{Carte mentale du thème}

Pour mémoriser le lexique et organiser les idées du texte, voici une carte mentale du champ lexical des droits :

\begin{popfigure}
\begin{tikzpicture}[mindmap, grow cyclic, every node/.style={concept, text=white, align=center}, root concept/.append style={concept color=popblue, font=\small\bfseries}, level 1/.append style={level distance=3.4cm, sibling angle=90}]
\node[root concept]{derechos de las mujeres y minorías}
  child[concept color=poporange]{ node {mujeres} }
  child[concept color=popgreen]{ node {minorías étnicas} }
  child[concept color=poppurple]{ node {leyes} }
  child[concept color=poppink]{ node {reivindica\-ciones} };
\end{tikzpicture}
\legende{Carte mentale : le champ lexical des droits autour du thème du texte.}
\end{popfigure}

Chaque branche pourra être enrichie au fil de la leçon : \esp{sufragio, igualdad salarial} pour « mujeres » ; \esp{lengua, origen, exclusión} pour « minorías étnicas » ; \esp{igualdad de oportunidades, derechos humanos} pour « leyes » ; \esp{manifestación, asociación, exigir} pour « reivindicaciones ».

\begin{savaistu}[Le droit de vote des femmes en Espagne]
En Espagne, les femmes ont obtenu le droit de vote (\esp{sufragio femenino}) en 1931, sous la Seconde République, et l'ont exercé pour la première fois aux élections de 1933. C'était bien avant certains pays d'Amérique latine : en Équateur, le vote féminin date de 1929, mais au Paraguay il faudra attendre 1961 ! En France, rappelons-le, les femmes votent pour la première fois en 1945.
\end{savaistu}

\begin{exercice}[À vous de lire]
Relisez le texte en silence, puis répondez par écrit en trois phrases en français : 1) Quelle est la thèse de l'auteur (ce qu'elle défend) ? 2) Relevez deux mots transparents supplémentaires non présents dans le tableau. 3) Le texte se termine par une subordonnée introduite par \esp{donde} : que exprime cette dernière phrase ?
\end{exercice}

\begin{corrige}[Exercice « À vous de lire »]
1) La thèse de l'auteur : les droits des femmes et des minorités ont été conquis par la lutte, les lois progressent, mais il faut encore changer les mentalités pour construire une société plus juste. 2) Mots transparents possibles : \esp{asociaciones} (associations), \esp{histórico} (historique), \esp{presente} (présente), \esp{reconocido} (reconnu). 3) La phrase finale exprime le \textbf{but} ou l'\textbf{idéal} visé : une société où personne n'est exclu à cause de son sexe, de son origine ou de sa langue. Notez le subjonctif \esp{sea} après \esp{donde}, qui marque ici un but non encore atteint — nous y reviendrons dans la section sur la revendication.
\end{corrige}

\begin{aretenir}[L'essentiel de la première lecture]
\begin{itemize}
\item Avant de lire : observer le titre, la source, l'auteur ; formuler des hypothèses.
\item Pendant : lire d'abord le premier et le dernier paragraphe, puis le texte entier en soulignant les mots transparents.
\item Après : répondre aux trois questions — de quoi parle le texte ? qui parle ? à qui ?
\item Les mots en \esp{-ción} sont féminins et correspondent au français « -tion ».
\item Vocabulaire clé : \esp{reivindicar} (revendiquer), \esp{la lucha} (la lutte), \esp{el sufragio} (le droit de vote), \esp{la igualdad de oportunidades} (l'égalité des chances), \esp{hace falta} (il faut).
\end{itemize}
\end{aretenir}

\coursec{Compréhension détaillée du texte}

Dans la section précédente, nous avons découvert le texte support \esp{« Una ley no basta »} et effectué une première lecture globale : nous savons déjà qu'il s'agit d'un article d'opinion sur les droits des femmes et des minorités. Il s'agit maintenant de passer à la \cle{lecture analytique} : comprendre le texte dans le détail, repérer ses idées directrices, dégager le point de vue de l'auteur et apprendre à justifier chaque réponse par une citation précise. C'est exactement la compétence exigée à l'épreuve d'espagnol du Baccalauréat A.

\courssub{Le texte support : relecture guidée}

Relisons le texte intégralement, cette fois avec un glossaire, avant de répondre aux questions.

\begin{textebox}[« Una ley no basta » — artículo de opinión]
A lo largo de la historia, muchas personas han tenido que reivindicar derechos que hoy nos parecen evidentes. Las mujeres y las minorías han sido, en casi todas las sociedades, las primeras víctimas de la discriminación: se les negaba la educación, el voto o el simple derecho a trabajar con el mismo salario.

El sufragio femenino es un ejemplo histórico de esta lucha. En muchos países, las mujeres votaron por primera vez en el siglo XX, después de décadas de manifestaciones, de escritos y de sacrificios. Lo que hoy parece un derecho natural fue ayer una conquista difícil.

Sin embargo, la discriminación persiste. En el mundo del trabajo, muchas mujeres cobran menos que los hombres por el mismo empleo. En algunas regiones, las minorías étnicas o lingüísticas no tienen acceso a los mismos servicios. En la vida cotidiana, las bromas y los prejuicios siguen hiriendo a quienes son diferentes.

Es verdad que existen leyes que protegen la igualdad de oportunidades. Pero una ley no basta: una ley puede prohibir la discriminación, pero no puede cambiar por sí sola las mentalidades. El cambio verdadero ocurre en la escuela, en la familia y en los medios de comunicación.

Por eso, reivindicar los derechos humanos es tarea de todos. Los jóvenes de hoy, en España, en América Latina, en Guinea Ecuatorial o en Camerún, tienen en sus manos la posibilidad de construir una sociedad más justa. La historia demuestra que cada generación puede avanzar: la nuestra no debe ser la excepción.
\end{textebox}

\textbf{Glossaire (mots difficiles) :}

\begin{center}
\begin{tabularx}{\linewidth}{|X|X|}
\hline
\rowcolor{popblueL} \textbf{Español} & \textbf{Français} \\
\hline
a lo largo de la historia & au cours de l'histoire \\
\hline
reivindicar & revendiquer, réclamer \\
\hline
se les negaba & on leur refusait (imparfait) \\
\hline
el voto / el sufragio & le vote / le suffrage (droit de vote) \\
\hline
una conquista difícil & une conquête difficile, un acquis \\
\hline
cobrar menos & être payé moins, gagner moins \\
\hline
las bromas y los prejuicios & les plaisanteries et les préjugés \\
\hline
herir (hiriendo) & blesser (en blessant) \\
\hline
por sí sola & à elle seule \\
\hline
los medios de comunicación & les médias \\
\hline
tarea de todos & affaire de tous, devoir de tous \\
\hline
avanzar & progresser, avancer \\
\hline
\end{tabularx}
\end{center}

\courssub{Questions de compréhension globale}

Avant d'entrer dans les détails, on vérifie la compréhension d'ensemble : le \cle{thème}, le \cle{genre} du texte et le \cle{destinataire}.

\begin{exemplebox}[Questions globales et réponses modèles]
\textbf{1. ¿Cuál es el tema del texto?}

\esp{El tema del texto es la lucha por los derechos de las mujeres y de las minorías.} (« Le thème du texte est la lutte pour les droits des femmes et des minorités. »)

\textbf{2. ¿Qué tipo de texto es?}

\esp{Es un artículo de opinión, probablemente publicado en un periódico, porque el autor defiende una tesis y se dirige directamente al lector.} (« C'est un article d'opinion, probablement publié dans un journal, car l'auteur défend une thèse et s'adresse directement au lecteur. »)

\textbf{3. ¿A quién se dirige el autor?}

\esp{El autor se dirige a todos los ciudadanos, y en particular a los jóvenes, como muestra la frase « los jóvenes de hoy... tienen en sus manos la posibilidad de construir una sociedad más justa ».} (« L'auteur s'adresse à tous les citoyens, et en particulier aux jeunes... »)
\end{exemplebox}

\begin{attention}[Bien identifier le genre du texte]
Ne confondez pas \esp{artículo de opinión} (article d'opinion, l'auteur défend une thèse) et \esp{artículo informativo} (article d'information, l'auteur rapporte des faits de manière neutre). Ici, la présence de jugements de valeur (\esp{« una conquista difícil »}, \esp{« una sociedad más justa »}) et d'un appel final prouve qu'il s'agit d'un texte \cle{argumentatif} et \cle{engagé}. Autre piège : \esp{género} signifie « genre » (littéraire ou sexuel), pas « général ».
\end{attention}

\courssub{Questions de compréhension fine}

Ces questions portent sur des informations précises, localisées dans un paragraphe. La règle d'or : répondre par une phrase complète et justifier par une \cle{citation courte} entre guillemets.

\begin{exemplebox}[Questions fines et réponses modèles]
\textbf{1. ¿Quiénes han tenido que reivindicar sus derechos?}

\esp{Las mujeres y las minorías han tenido que reivindicar sus derechos.} Justification : \esp{« Las mujeres y las minorías han sido, en casi todas las sociedades, las primeras víctimas de la discriminación. »} (« Les femmes et les minorités ont été, dans presque toutes les sociétés, les premières victimes de la discrimination. »)

\textbf{2. ¿Qué ejemplo histórico de esta lucha se cita en el texto?}

\esp{El ejemplo histórico citado es el sufragio femenino.} Justification : \esp{« El sufragio femenino es un ejemplo histórico de esta lucha. »} (« Le suffrage féminin est un exemple historique de cette lutte. »)

\textbf{3. ¿Dónde persiste la discriminación hoy en día?}

\esp{La discriminación persiste en el mundo del trabajo, en el acceso a los servicios para las minorías y en la vida cotidiana.} Justification : \esp{« muchas mujeres cobran menos que los hombres por el mismo empleo »} et \esp{« las bromas y los prejuicios siguen hiriendo a quienes son diferentes »}. (« ...beaucoup de femmes gagnent moins que les hommes pour le même emploi » ; « les plaisanteries et les préjugés continuent de blesser ceux qui sont différents ».)

\textbf{4. Según el autor, ¿basta una ley para acabar con la discriminación?}

\esp{No, una ley no basta porque no puede cambiar por sí sola las mentalidades.} Justification : \esp{« Una ley no basta: una ley puede prohibir la discriminación, pero no puede cambiar por sí sola las mentalidades. »} (« Une loi ne suffit pas : une loi peut interdire la discrimination, mais elle ne peut pas changer à elle seule les mentalités. »)
\end{exemplebox}

\begin{methode}[Répondre à une question de compréhension au Bac]
\begin{enumerate}
\item \textbf{Repérer} le paragraphe concerné grâce aux mots-clés de la question (\esp{ejemplo histórico} $\rightarrow$ paragraphe 2).
\item \textbf{Rédiger} une phrase complète en espagnol, avec sujet + verbe + complément ; ne jamais répondre par un seul mot.
\item \textbf{Justifier} par une citation courte entre guillemets, introduite par \esp{como dice el texto} ou \esp{según el autor}.
\item \textbf{Traduire} la citation en français si la consigne le demande (\esp{traduce la cita}).
\item \textbf{Vérifier} que la réponse est dans la langue exigée : question en espagnol $\rightarrow$ réponse en espagnol.
\end{enumerate}
\end{methode}

\begin{attention}[La citation justificative : courte et pertinente]
L'erreur la plus fréquente consiste à recopier tout le paragraphe comme justification : le correcteur y voit une absence d'analyse. Ne citez \textbf{que le segment pertinent} : quelques mots ou une proposition suffisent. Par exemple, pour la question sur l'exemple historique, citez uniquement \esp{« El sufragio femenino es un ejemplo histórico de esta lucha »} et non les trois phrases du paragraphe 2. Autre piège : n'oubliez jamais les guillemets ; une phrase recopiée sans guillemets est considérée comme du plagiat du texte, pas comme une justification.
\end{attention}

\courssub{Repérage des idées directrices : une idée par paragraphe}

Un texte argumentatif bien construit développe \cle{une idée directrice par paragraphe}. Pour la trouver, on repère la phrase la plus générale du paragraphe (souvent la première ou la dernière) et on la résume en quelques mots.

\begin{center}
\begin{tabularx}{\linewidth}{|l|X|X|}
\hline
\rowcolor{popblueL} \textbf{§} & \textbf{Idea directriz (español)} & \textbf{Idée directrice (français)} \\
\hline
1 & \esp{Muchas personas han tenido que reivindicar sus derechos.} & Constat historique : des droits longtemps refusés \\
\hline
2 & \esp{El sufragio femenino es un ejemplo histórico de esta lucha.} & Exemple : la conquête du vote des femmes \\
\hline
3 & \esp{La discriminación persiste.} & Les discriminations actuelles \\
\hline
4 & \esp{Una ley no basta.} & La loi ne suffit pas à changer les mentalités \\
\hline
5 & \esp{Reivindicar los derechos humanos es tarea de todos.} & Appel à l'action, surtout des jeunes \\
\hline
\end{tabularx}
\end{center}

\begin{popfigure}
\begin{center}
\begin{tikzpicture}[
  node distance=0.55cm and 0.9cm,
  box/.style={rectangle, rounded corners=3pt, draw=popblue, fill=popblueL, text width=4.6cm, align=center, font=\small, inner sep=4pt},
  lbl/.style={font=\footnotesize\bfseries, text=popdark},
  arr/.style={-{Stealth[length=2.5mm]}, thick, poporange}
]
\node[box] (p1) {\textbf{§1} Constat historique : mujeres y minorías, víctimas de la discriminación};
\node[box, below=of p1] (p2) {\textbf{§2} Exemple : el sufragio femenino, una conquista difícil};
\node[box, below=of p2] (p3) {\textbf{§3} Aujourd'hui : la discriminación persiste (trabajo, servicios, vida cotidiana)};
\node[box, below=of p3] (p4) {\textbf{§4} Limites : una ley no basta, hay que cambiar las mentalidades};
\node[box, below=of p4, draw=popgreen, fill=popgreenL] (p5) {\textbf{§5} Appel à l'action : los jóvenes pueden construir una sociedad más justa};
\draw[arr] (p1) -- (p2);
\draw[arr] (p2) -- (p3);
\draw[arr] (p3) -- (p4);
\draw[arr] (p4) -- (p5);
\node[lbl, left=0.25cm of p1.west, anchor=east] {pasado};
\node[lbl, left=0.25cm of p3.west, anchor=east] {presente};
\node[lbl, left=0.25cm of p5.west, anchor=east] {futuro};
\end{tikzpicture}
\end{center}
\legende{Organigramme de la structure du texte : du constat historique à l'appel à l'action.}
\end{popfigure}

Remarquez la progression chronologique du texte : \esp{pasado} (§1--2) $\rightarrow$ \esp{presente} (§3--4) $\rightarrow$ \esp{futuro} (§5). Cette architecture \cle{chronologique} est classique dans les articles engagés : elle montre un cheminement du constat vers l'espoir.

\courssub{Dégager le point de vue de l'auteur}

Le \cle{point de vue de l'auteur} se déduit de l'ensemble du texte : son vocabulaire, ses connecteurs, ses procédés d'insistance.

\begin{propriete}[Indices d'un texte engagé]
\begin{itemize}
\item \textbf{Lexique valorisant la lutte} : \esp{lucha}, \esp{conquista}, \esp{sacrificios}, \esp{reivindicar}.
\item \textbf{Lexique dévalorisant la discrimination} : \esp{víctimas}, \esp{prejuicios}, \esp{hiriendo}.
\item \textbf{Concession puis réfutation} : \esp{« Es verdad que existen leyes... Pero una ley no basta »} — l'auteur reconnaît un argument adverse pour mieux le dépasser.
\item \textbf{Appel direct au lecteur} : \esp{« los jóvenes de hoy... tienen en sus manos »}.
\item \textbf{Phrase finale sentencieuse} : \esp{« la nuestra no debe ser la excepción »}.
\end{itemize}
\end{propriete}

\begin{exemplebox}[Le point de vue de l'auteur : réponse modèle]
\esp{El autor defiende un punto de vista comprometido: denuncia la discriminación y anima a los jóvenes a luchar por la igualdad. Su tono es optimista pero realista: cree en el progreso (« cada generación puede avanzar »), pero reconoce que la discriminación persiste y que una ley no basta.} (« L'auteur défend un point de vue engagé : il dénonce la discrimination et encourage les jeunes à lutter pour l'égalité. Son ton est optimiste mais réaliste : il croit au progrès, tout en reconnaissant que la discrimination persiste et qu'une loi ne suffit pas. »)
\end{exemplebox}

\begin{attention}[Optimiste mais réaliste : nuancer, toujours]
Au Bac, une analyse du type « l'auteur est optimiste » est jugée insuffisante si elle ignore le §3 (\esp{« la discriminación persiste »}). La bonne formule est \cle{nuancée} : \esp{optimista pero realista}. Méfiez-vous aussi du faux ami : \esp{comprometido} signifie « engagé » (politiquement, socialement), et non « compromis » ; « compromettre » se dit \esp{comprometer}.
\end{attention}

\courssub{Reformulation : exprimer la thèse avec ses propres mots}

Dernière étape de la compréhension : être capable de \cle{reformuler} la thèse de l'auteur sans recopier ses phrases. C'est la preuve que le texte est réellement compris. On utilise des synonymes et des périphrases.

\begin{center}
\begin{tabularx}{\linewidth}{|X|X|}
\hline
\rowcolor{popblueL} \textbf{Palabras del texto} & \textbf{Reformulación posible} \\
\hline
\esp{reivindicar derechos} & \esp{luchar por los derechos / exigir derechos} \\
\hline
\esp{una ley no basta} & \esp{la ley no es suficiente / no sirve solo con legislar} \\
\hline
\esp{cambiar las mentalidades} & \esp{transformar las actitudes / educar a la sociedad} \\
\hline
\esp{una sociedad más justa} & \esp{un mundo más igualitario} \\
\hline
\esp{tarea de todos} & \esp{responsabilidad de cada ciudadano} \\
\hline
\end{tabularx}
\end{center}

\begin{exemplebox}[Reformulation complète de la thèse]
\esp{Según el autor, la igualdad real entre hombres y mujeres y entre mayorías y minorías no se consigue únicamente con leyes: hace falta educar a la sociedad y comprometerse personalmente, sobre todo cuando se es joven.} (« Selon l'auteur, l'égalité réelle entre hommes et femmes et entre majorités et minorités ne s'obtient pas uniquement par des lois : il faut éduquer la société et s'engager personnellement, surtout quand on est jeune. »)
\end{exemplebox}

\begin{exoresolu}[Question type Bac : compréhension fine + justification]
Énoncé : \esp{¿Por qué dice el autor que « una ley no basta »? Justifica tu respuesta con una cita del texto.}

\tcblower

\textbf{Solution détaillée.}

Étape 1 : le mot-clé \esp{ley} oriente vers le paragraphe 4.

Étape 2 : réponse en phrase complète :

\esp{El autor dice que una ley no basta porque una ley puede prohibir la discriminación, pero no puede cambiar las mentalidades; el cambio verdadero ocurre en la escuela, en la familia y en los medios de comunicación.}

Étape 3 : justification par une citation courte :

\esp{Como dice el texto: « una ley puede prohibir la discriminación, pero no puede cambiar por sí sola las mentalidades ».}

Étape 4 : traduction de la citation (si demandée) : « une loi peut interdire la discrimination, mais elle ne peut pas changer à elle seule les mentalités ».

\textbf{Barème-type} : réponse correcte et complète (1 point) ; citation pertinente entre guillemets (0,5 point) ; correction linguistique (0,5 point).
\end{exoresolu}

\begin{exercice}[À vous de jouer]
Répondez en espagnol, avec une citation courte justificative :
\begin{enumerate}
\item \esp{¿Qué derechos se negaban a las mujeres y a las minorías? (§1)}
\item \esp{¿Cuándo votaron por primera vez las mujeres en muchos países? (§2)}
\item \esp{¿Dónde ocurre el cambio verdadero, según el autor? (§4)}
\item Reformulez avec vos propres mots : \esp{« Reivindicar los derechos humanos es tarea de todos »}.
\end{enumerate}
\end{exercice}

\begin{corrige}[Exercice : compréhension fine]
\begin{enumerate}
\item \esp{Se les negaba la educación, el voto y el derecho a trabajar con el mismo salario.} Citation : \esp{« se les negaba la educación, el voto o el simple derecho a trabajar con el mismo salario »}.
\item \esp{En muchos países, las mujeres votaron por primera vez en el siglo XX.} Citation : \esp{« las mujeres votaron por primera vez en el siglo XX »}.
\item \esp{El cambio verdadero ocurre en la escuela, en la familia y en los medios de comunicación.} Citation : \esp{« El cambio verdadero ocurre en la escuela, en la familia y en los medios de comunicación »}.
\item Reformulation possible : \esp{Defender los derechos humanos es responsabilidad de cada ciudadano.} (« Défendre les droits humains est la responsabilité de chaque citoyen. »)
\end{enumerate}
\end{corrige}

\begin{aretenir}[L'essentiel de la compréhension détaillée]
\begin{itemize}
\item Compréhension globale : thème (droits des femmes et des minorités), genre (article d'opinion engagé), destinataire (tous les citoyens, surtout les jeunes).
\item Compréhension fine : phrase complète + citation \textbf{courte} entre guillemets + traduction si demandée.
\item Une idée directrice par paragraphe ; le texte progresse du passé vers le futur.
\item Point de vue de l'auteur : engagé, \esp{optimista pero realista}.
\item Savoir reformuler la thèse avec des synonymes : c'est la preuve ultime de la compréhension.
\end{itemize}
\end{aretenir}

\coursec{Lexique thématique : los derechos humanos}

Après avoir compris le texte de la section précédente dans son ensemble — ses idées directrices, son point de vue, son ton —, nous allons maintenant nous arrêter sur ses mots. Un commentaire de texte au Bac ne peut pas se contenter de paraphraser : il faut \emph{nommer} précisément les réalités dont parle l'auteur. Or le thème des droits des femmes et des minorités possède en espagnol un vocabulaire riche, codifié, et truffé de faux amis pour le francophone. Cette section vous donne les clés lexicales indispensables, classées en quatre champs, puis vous entraîne à les réemployer immédiatement.

\courssub{Observer d'abord : un texte pour faire émerger le lexique}

Lisons ce court article de presse (texte original). Soulignez mentalement tous les mots qui appartiennent au thème des droits.

\begin{textebox}[Un artículo de prensa : « La larga marcha hacia la igualdad »]
La historia de los derechos humanos es también la historia de las mujeres y de las minorías que tuvieron que luchar para conquistar lo que otros ya poseían. Durante siglos, las mujeres no tenían derecho al sufragio: no podían votar ni ser elegidas. La igualdad de oportunidades era una palabra vacía, y la discriminación formaba parte de la vida cotidiana, en la escuela, en el trabajo y hasta dentro de la familia.

Hoy, la ley protege a las minorías y prohíbe el racismo y el sexismo, pero la desigualdad sigue existiendo. En muchos países, las mujeres reivindican la igualdad salarial: a trabajo igual, salario igual. Los colectivos de ciudadanos se manifiestan en las calles para exigir justicia y defender la libertad de todos. Los jóvenes, especialmente, protestan contra la exclusión social y los prejuicios que todavía dividen a la sociedad.

En África central, de Camerún a Guinea Ecuatorial, las asociaciones luchan contra la discriminación y por la integración social de los más desfavorecidos. Todos tienen derecho a la educación, a la salud y al trabajo digno: esto no es un regalo, es una conquista que hay que defender cada día. La equidad no llega sola: se construye con leyes justas, con escuelas abiertas y con ciudadanos que no aceptan la injusticia en silencio.
\end{textebox}

\textbf{Glossaire :} \esp{el sufragio} = le droit de vote ; \esp{elegida} = élue ; \esp{una palabra vacía} = une expression creuse ; \esp{prohibir} = interdire ; \esp{reivindicar} = revendiquer ; \esp{se manifiestan} = manifestent ; \esp{exigir} = exiger ; \esp{los prejuicios} = les préjugés ; \esp{los desfavorecidos} = les défavorisés ; \esp{una conquista} = une conquête, un acquis ; \esp{la equidad} = l'équité.

\textbf{Questions de repérage :} (1) Quels droits les femmes ne possédaient-elles pas autrefois ? (2) Que revendiquent-elles aujourd'hui ? (3) Quels verbes expriment l'action de protester ? (4) Selon l'auteur, l'équité arrive-t-elle toute seule ?

\courssub{Premier champ lexical : los derechos (les droits)}

\begin{definition}[Derecho, derechos humanos]
\esp{\cle{El derecho}} = le droit (ce qui est reconnu à chacun : \esp{tener derecho a la educación}). \esp{\cle{Los derechos humanos}} = les droits de l'homme / les droits humains. Attention au genre : \esp{el derecho} est masculin, contrairement à \esp{\cle{la ley}} (la loi).
\end{definition}

Observons les mots du texte : \esp{derecho}, \esp{ley}, \esp{sufragio}, \esp{igualdad}, \esp{libertad}, \esp{justicia}. Ce sont les noms \emph{positifs} du thème, ceux qui désignent ce que la société doit garantir.

\begin{center}
\begin{tabularx}{\linewidth}{|X|X|X|}
\hline
\rowcolor{popblueL} Español & Français & Remarque \\
\hline
\esp{\cle{el derecho}} & le droit & \esp{tener derecho a} + nom/infinitif \\
\hline
\esp{\cle{los derechos humanos}} & les droits humains & toujours au pluriel dans ce sens \\
\hline
\esp{\cle{la ley}} & la loi & féminin ! \esp{la ley protege...} \\
\hline
\esp{\cle{el sufragio}} & le suffrage, le droit de vote & \esp{sufragio universal} = suffrage universel \\
\hline
\esp{\cle{la igualdad}} & l'égalité & de l'adjectif \esp{igual} \\
\hline
\esp{\cle{la libertad}} & la liberté & \esp{libertad de expresión} \\
\hline
\esp{\cle{la justicia}} & la justice & \esp{justicia social} \\
\hline
\end{tabularx}
\end{center}

\begin{exemplebox}[Exemples traduits]
\esp{La ley protege a las minorías.} « La loi protège les minorités. »\par
\esp{Todos tienen derecho a la educación.} « Tous ont droit à l'éducation. »\par
\esp{Las mujeres conquistaron el sufragio después de una larga lucha.} « Les femmes ont conquis le droit de vote après une longue lutte. »\par
\esp{En Yaoundé, los estudiantes hablan de libertad de expresión.} « À Yaoundé, les étudiants parlent de liberté d'expression. »
\end{exemplebox}

\courssub{Deuxième champ lexical : la discriminación}

\begin{definition}[Minoría]
\esp{\cle{La minoría}} : grupo de personas que se distingue del resto de la sociedad por su origen, su lengua o su religión. En français : groupe minoritaire, minorité. Exemple : \esp{las minorías lingüísticas} = les minorités linguistiques.
\end{definition}

Ce champ regroupe les mots \emph{négatifs} : ce que l'on combat. \esp{\cle{La discriminación}} (la discrimination) est le nom central ; il se construit avec \esp{contra} : \esp{la discriminación contra las mujeres} = la discrimination envers les femmes. On trouve aussi \esp{\cle{la exclusión}} (l'exclusion), \esp{\cle{la desigualdad}} (l'inégalité), \esp{\cle{el prejuicio}} (le préjugé), \esp{\cle{el racismo}} (le racisme), \esp{\cle{el sexismo}} (le sexisme).

\begin{attention}[Faux amis et pièges de genre]
\begin{itemize}
\item \esp{\cle{el racismo}}, \esp{\cle{el sexismo}}, \esp{\cle{el prejuicio}} sont \textbf{masculins} : \esp{el sexismo existe}, pas \esp{la sexismo}.
\item \esp{\cle{la desigualdad}} est le contraire de \esp{\cle{la igualdad}} : préfixe \esp{des-} comme en français « dés- » (\esp{desigual} = inégal).
\item \esp{discriminación positiva} = discrimination positive (terme identique au français).
\end{itemize}
\end{attention}

\begin{exemplebox}[Exemples traduits]
\esp{El racismo y el sexismo son formas de discriminación.} « Le racisme et le sexisme sont des formes de discrimination. »\par
\esp{Los prejuicios dividen a la sociedad.} « Les préjugés divisent la société. »\par
\esp{La exclusión social afecta a los más pobres.} « L'exclusion sociale touche les plus pauvres. »
\end{exemplebox}

\courssub{Troisième champ lexical : la reivindicación (la revendication)}

C'est le champ des \emph{verbes d'action}, le plus important pour vos productions écrites. Observons comment ils se construisent :

\begin{center}
\begin{tabularx}{\linewidth}{|X|X|X|}
\hline
\rowcolor{popgreenL} Verbo & Construcción & Français \\
\hline
\esp{\cle{reivindicar}} & + nom & revendiquer \\
\hline
\esp{\cle{luchar}} & \esp{por} + nom / \esp{contra} + nom & lutter pour / contre \\
\hline
\esp{\cle{protestar}} & \esp{contra} + nom & protester contre \\
\hline
\esp{\cle{manifestarse}} & (verbe pronominal) & manifester \\
\hline
\esp{\cle{exigir}} & + nom & exiger \\
\hline
\esp{\cle{defender}} (e$\to$ie) & + nom & défendre \\
\hline
\esp{\cle{conquistar}} & \esp{un derecho} & conquérir (un droit) \\
\hline
\end{tabularx}
\end{center}

\begin{propriete}[Constructions à mémoriser]
\begin{itemize}
\item \esp{luchar por la igualdad} = lutter pour l'égalité ; \esp{luchar contra la discriminación} = lutter contre la discrimination.
\item \esp{manifestarse} est pronominal : \esp{los ciudadanos \textbf{se} manifiestan} = les citoyens manifestent. Au passé : \esp{se manifestaron ayer}.
\item \esp{defender} est irrégulier (e$\to$ie) : \esp{defiendo, defiendes, defiende, defendemos, defendéis, defienden}.
\item \esp{exigir} fait \esp{exijo} à la première personne (g$\to$j).
\end{itemize}
\end{propriete}

\begin{exemplebox}[Exemples traduits]
\esp{Las mujeres reivindican la igualdad salarial.} « Les femmes revendiquent l'égalité salariale. »\par
\esp{Los jóvenes se manifiestan para exigir justicia.} « Les jeunes manifestent pour exiger la justice. »\par
\esp{Defendemos los derechos de las minorías.} « Nous défendons les droits des minorités. »\par
\esp{Las asociaciones luchan contra la exclusión en Douala.} « Les associations luttent contre l'exclusion à Douala. »
\end{exemplebox}

\begin{attention}[Reivindicar : un verbe noble]
\esp{\cle{Reivindicar}} ne signifie pas seulement « réclamer bruyamment ». C'est le verbe exact pour « revendiquer un droit », au sens noble et positif : \esp{reivindicar la igualdad} = revendiquer l'égalité. Ne le remplacez pas par \esp{reclamar}, qui signifie plutôt « réclamer, se plaindre » (une réclamation commerciale = \esp{una reclamación}). Autre piège capital : ne confondez jamais \esp{\cle{el sufragio}} (le droit de vote) et \esp{\cle{el sufrimiento}} (la souffrance) ! Écrire \esp{las mujeres lucharon por el sufrimiento} dirait exactement le contraire de ce que vous voulez dire.
\end{attention}

\courssub{Quatrième champ lexical : la igualdad y la integración}

Ce champ exprime le but à atteindre : \esp{\cle{la igualdad de oportunidades}} (l'égalité des chances), \esp{\cle{la igualdad salarial}} (l'égalité salariale), \esp{\cle{la equidad}} (l'équité), \esp{\cle{la integración social}} (l'intégration sociale).

\begin{savaistu}[« Igualdad de oportunidades » au Bac]
\esp{\cle{La igualdad de oportunidades}} est l'équivalent exact de « l'égalité des chances ». Cette expression revient très fréquemment dans les sujets du Bac, dans les textes de presse comme dans les sujets d'expression écrite. Apprenez-la comme un bloc, avec son article : \esp{la igualdad de oportunidades}. Notez la nuance : \esp{igualdad} désigne l'égalité de principe (mêmes droits), tandis que \esp{equidad} désigne l'équité, c'est-à-dire la justice concrète qui donne à chacun ce dont il a besoin.
\end{savaistu}

\begin{exemplebox}[Expressions utiles pour rédiger]
\esp{\cle{tener derecho a}} + nom/infinitif : \esp{Tenemos derecho a estudiar.} « Nous avons le droit d'étudier. »\par
\esp{\cle{luchar contra la discriminación}} : « lutter contre la discrimination »\par
\esp{\cle{estar a favor de / en contra de}} : \esp{Estoy a favor de la igualdad salarial.} « Je suis pour l'égalité salariale. » \esp{Estamos en contra del racismo.} « Nous sommes contre le racisme. »\par
\esp{\cle{a trabajo igual, salario igual}} : « à travail égal, salaire égal » (slogan très fréquent).
\end{exemplebox}

\courssub{Carte mentale et tableau récapitulatif}

\begin{popfigure}
\begin{center}
\begin{tikzpicture}[mindmap, grow cyclic, every node/.style={concept, align=center, font=\small},
root concept/.append style={concept color=popblue, font=\small\bfseries},
level 1/.append style={level distance=3.4cm, sibling angle=90},
level 2/.append style={level distance=2.4cm, sibling angle=45}]
\node[root concept]{Los derechos humanos}
  child[concept color=popgreen] { node {Los derechos}
    child { node[concept color=popgreenL, text=popdark, align=center]{el derecho\\la ley} }
    child { node[concept color=popgreenL, text=popdark, align=center]{el sufragio\\la libertad} }
  }
  child[concept color=poporange] { node {La discriminación}
    child { node[concept color=poporangeL, text=popdark, align=center]{el racismo\\el sexismo} }
    child { node[concept color=poporangeL, text=popdark, align=center]{el prejuicio\\la exclusión} }
  }
  child[concept color=poppurple] { node {La reivindicación}
    child { node[concept color=poppurpleL, text=popdark, align=center]{reivindicar\\exigir} }
    child { node[concept color=poppurpleL, text=popdark, align=center]{luchar por\\manifestarse} }
  }
  child[concept color=popgold] { node {La igualdad}
    child { node[concept color=popgoldL, text=popdark, align=center]{igualdad de\\oportunidades} }
    child { node[concept color=popgoldL, text=popdark, align=center]{equidad\\integración} }
  };
\end{tikzpicture}
\end{center}
\legende{Carte mentale des quatre champs lexicaux du thème des droits.}
\end{popfigure}

\begin{center}
\begin{tabularx}{\linewidth}{|l|X|X|}
\hline
\rowcolor{popblueL} Campo & Español & Français \\
\hline
Droits & \esp{el derecho, la ley, el sufragio, la libertad, la justicia} & le droit, la loi, le suffrage, la liberté, la justice \\
\hline
Discrimination & \esp{la minoría, la exclusión, la desigualdad, el prejuicio, el racismo, el sexismo} & la minorité, l'exclusion, l'inégalité, le préjugé, le racisme, le sexisme \\
\hline
Revendication & \esp{reivindicar, luchar por, protestar, manifestarse, exigir, defender, conquistar} & revendiquer, lutter pour, protester, manifester, exiger, défendre, conquérir \\
\hline
Égalité & \esp{la igualdad de oportunidades, la igualdad salarial, la equidad, la integración social} & l'égalité des chances, l'égalité salariale, l'équité, l'intégration sociale \\
\hline
\end{tabularx}
\end{center}

\courssub{Réemploi immédiat : exercice résolu}

\begin{exoresolu}[Complétez avec le lexique étudié]
Complétez les phrases avec : \esp{sufragio}, \esp{reivindican}, \esp{derecho}, \esp{discriminación}, \esp{igualdad de oportunidades}, \esp{se manifiestan}.\par
1. Todos los ciudadanos tienen \ldots{} a la educación.\par
2. Las mujeres \ldots{} la igualdad salarial en muchas empresas.\par
3. El \ldots{} universal permite votar a todos los adultos.\par
4. Hay que luchar contra la \ldots{} racial.\par
5. Los estudiantes \ldots{} en las calles de la capital.\par
6. La \ldots{} significa que todos pueden acceder a los mismos empleos.
\tcblower
\textbf{Solution détaillée :}\par
1. \esp{Todos los ciudadanos tienen \textbf{derecho} a la educación.} — construction \esp{tener derecho a} + nom.\par
2. \esp{Las mujeres \textbf{reivindican} la igualdad salarial en muchas empresas.} — sujet pluriel, troisième personne : \esp{reivindican}.\par
3. \esp{El \textbf{sufragio} universal permite votar a todos los adultos.} — attention : pas \esp{sufrimiento} !\par
4. \esp{Hay que luchar contra la \textbf{discriminación} racial.} — \esp{luchar contra} + nom.\par
5. \esp{Los estudiantes \textbf{se manifiestan} en las calles de la capital.} — verbe pronominal : ne pas oublier \esp{se}.\par
6. \esp{La \textbf{igualdad de oportunidades} significa que todos pueden acceder a los mismos empleos.} — expression figée à apprendre en bloc.
\end{exoresolu}

\begin{exercice}[À vous de jouer]
Traduisez en espagnol : (1) « La loi protège les minorités. » (2) « Les femmes revendiquent l'égalité des chances. » (3) « Nous luttons contre le racisme et le sexisme. » (4) « Tous ont droit à l'éducation. » (5) « Les jeunes manifestent pour défendre la liberté. »
\end{exercice}

\begin{corrige}[Exercice : traduction]
(1) \esp{La ley protege a las minorías.} — \esp{ley} est féminin : \esp{la ley}.\par
(2) \esp{Las mujeres reivindican la igualdad de oportunidades.} — pas de préposition après \esp{reivindicar}.\par
(3) \esp{Luchamos contra el racismo y el sexismo.} — articles masculins : \esp{el racismo}, \esp{el sexismo}.\par
(4) \esp{Todos tienen derecho a la educación.} — \esp{tener derecho a}, sans article devant \esp{derecho}.\par
(5) \esp{Los jóvenes se manifiestan para defender la libertad.} — pronominal + \esp{para} + infinitif.
\end{corrige}

\begin{aretenir}[L'essentiel du lexique]
Quatre champs à mobiliser dans tout écrit sur ce thème : les \cle{droits} (\esp{derecho, ley, sufragio, libertad, justicia}), la \cle{discrimination} (\esp{discriminación, minoría, exclusión, prejuicio, racismo, sexismo}), la \cle{revendication} (\esp{reivindicar, luchar por/contra, manifestarse, exigir, defender, conquistar un derecho}) et l'\cle{égalité} (\esp{igualdad de oportunidades, igualdad salarial, equidad, integración social}). Trois constructions à ne jamais oublier : \esp{tener derecho a}, \esp{luchar contra la discriminación}, \esp{estar a favor de / en contra de}. Et deux pièges mortels : \esp{sufragio} $\neq$ \esp{sufrimiento} ; \esp{reivindicar} $\neq$ \esp{reclamar}.
\end{aretenir}

\coursec{Exprimer l'opinion et la revendication}

Dans la section précédente, nous avons construit le \cle{lexique des droits humains} : \esp{los derechos humanos}, \esp{la minoría}, \esp{la discriminación}, \esp{el sufragio}, \esp{la igualdad de oportunidades}, \esp{la ley}, \esp{reivindicar}. Posséder du vocabulaire ne suffit cependant pas : au commentaire de texte comme à l'expression écrite du Baccalauréat, l'examinateur attend que vous \textbf{preniez position} et que vous \textbf{revendiquiez}. Dire \esp{pienso que la ley protege a las minorías} ou \esp{es necesario que cambiemos las mentalidades}, c'est déjà commenter. Cette section vous donne les outils grammaticaux exacts pour le faire sans faute : le grand partage entre \cle{indicatif} et \cle{subjonctif}, les structures impersonnelles de nécessité, les verbes de revendication et les connecteurs de l'argumentation.

\courssub{Observer avant de théoriser : un éditorial revendicatif}

Lisons d'abord un court éditorial original. Soulignez mentalement tous les verbes qui suivent \esp{que}.

\begin{textebox}[« Voces que piden justicia », éditorial original]
En mi opinión, nuestra sociedad ha avanzado mucho, pero creo que todavía queda un largo camino. Pienso que las mujeres camerunesas y españolas tienen más oportunidades que sus abuelas; sin embargo, no creo que la igualdad sea una realidad completa. Me parece que las leyes existen, pero considero que no siempre se aplican.

No pienso que el sufragio femenino sea suficiente si las mujeres no acceden a los puestos de decisión. Además, hace falta que los medios de comunicación den más voz a las minorías. Es necesario que cambiemos las mentalidades desde la escuela: un niño que respeta a sus compañeras será un adulto que respeta a sus colegas.

Por eso, exigimos que el gobierno garantice la igualdad de oportunidades y pedimos que la ley castigue toda forma de discriminación. Hay que luchar contra los prejuicios cada día, en la familia, en el mercado, en el trabajo. En conclusión, reivindicamos que los derechos humanos sean para todos, sin excepción, porque una sociedad justa no es un regalo: es una conquista colectiva.
\end{textebox}

\textbf{Glossaire :} \esp{todavía queda un largo camino} = il reste encore un long chemin ; \esp{se aplican} = elles sont appliquées (verbe \esp{aplicarse}, passif réfléchi) ; \esp{los puestos de decisión} = les postes de décision ; \esp{den más voz} = donnent davantage la parole (subjonctif présent de \esp{dar}) ; \esp{castigue} = sanctionne (subjonctif présent de \esp{castigar}) ; \esp{los prejuicios} = les préjugés ; \esp{una conquista colectiva} = une conquête collective.

\textbf{Questions d'observation :}
\begin{enumerate}
\item Relevez trois verbes à l'indicatif après une structure d'opinion affirmative.
\item Relevez trois verbes au subjonctif et identifiez la structure qui les déclenche.
\item Quels connecteurs l'auteur utilise-t-il pour organiser son argumentation ?
\end{enumerate}

\textbf{Observation guidée.} Après \esp{creo que}, \esp{pienso que}, \esp{me parece que}, on trouve \esp{tienen}, \esp{existen}, \esp{queda} : des \textbf{indicatifs}. Après \esp{no creo que}, \esp{es necesario que}, \esp{exigimos que}, on trouve \esp{sea}, \esp{cambiemos}, \esp{garantice} : des \textbf{subjonctifs}. Ce n'est pas un hasard : c'est LA grande loi de l'expression de l'opinion en espagnol.

\courssub{Exprimer son opinion : l'indicatif après les verbes d'opinion affirmatifs}

\begin{propriete}[Opinion affirmative $\rightarrow$ indicatif]
Après les verbes d'opinion employés à la forme \textbf{affirmative} et suivis de \esp{que} — \esp{pienso que}, \esp{creo que}, \esp{me parece que}, \esp{considero que}, \esp{estoy seguro/a de que}, \esp{es verdad que}, \esp{es cierto que} —, le verbe de la subordonnée se met à l'\cle{indicatif}, car le locuteur présente le contenu comme un fait réel, acquis ou certain.
Les expressions figées sans \esp{que} comme \esp{en mi opinión}, \esp{a mi juicio} introduisent une proposition principale (donc indicatif) ou s'emploient en incise.
\end{propriete}

\begin{exemplebox}[Opinion affirmative + indicatif]
\begin{itemize}
\item \esp{Creo que la ley protege a las minorías.} = « Je crois que la loi protège les minorités. »
\item \esp{Pienso que las mujeres tienen las mismas capacidades que los hombres.} = « Je pense que les femmes ont les mêmes capacités que les hommes. »
\item \esp{Me parece que el texto denuncia la discriminación.} = « Il me semble que le texte dénonce la discrimination. »
\item \esp{En mi opinión, el sufragio fue una victoria histórica.} = « À mon avis, le suffrage fut une victoire historique. »
\item \esp{Considero que el autor tiene razón.} = « Je considère que l'auteur a raison. »
\end{itemize}
\end{exemplebox}

\begin{attention}[Faux ami : « me parece que »]
\esp{Me parece que} signifie « il me semble que » et se construit avec l'\textbf{indicatif} : \esp{Me parece que el autor exagera.} Ne le confondez pas avec le français « il me paraît que » (même sens, même construction), ni avec \esp{parecerse a} (« ressembler à »). Et surtout : \esp{me parece que} ne prend \textbf{JAMAIS} le subjonctif à la forme affirmative.
\end{attention}

\courssub{Opinion négative et nécessité impersonnelle : le règne du subjonctif}

\begin{propriete}[Opinion négative ou nécessité impersonnelle $\rightarrow$ subjonctif]
Deux grandes familles de structures exigent le \cle{subjonctif} dans la subordonnée introduite par \esp{que} :
\begin{enumerate}
\item Les verbes d'opinion à la forme \textbf{négative} : \esp{no pienso que}, \esp{no creo que}, \esp{no me parece que}, \esp{no considero que} — car le contenu est présenté comme douteux, nié, non réel.
\item Les tournures \textbf{impersonnelles de nécessité, d'importance ou de convenance} : \esp{es necesario que}, \esp{es importante que}, \esp{es imprescindible que}, \esp{es fundamental que}, \esp{hace falta que}, \esp{conviene que} — car elles expriment un but à atteindre, une volonté, pas un fait constaté.
\end{enumerate}
\end{propriete}

\begin{exemplebox}[Opinion négative et nécessité + subjonctif]
\begin{itemize}
\item \esp{No creo que la ley sea suficiente.} = « Je ne crois pas que la loi soit suffisante. »
\item \esp{No pienso que la discriminación haya desaparecido.} = « Je ne pense pas que la discrimination ait disparu. » (subjonctif parfait : \esp{haya} + participe)
\item \esp{Es necesario que cambiemos las mentalidades.} = « Il faut que nous changions les mentalités. »
\item \esp{Es importante que los jóvenes conozcan sus derechos.} = « Il est important que les jeunes connaissent leurs droits. »
\item \esp{Hace falta que el Estado proteja a las minorías.} = « Il faut que l'État protège les minorités. »
\end{itemize}
\end{exemplebox}

\begin{popfigure}
\begin{center}
\begin{tikzpicture}[
  font=\small,
  racine/.style={rectangle, rounded corners, draw=popdark, fill=popblueL, thick, align=center, minimum width=4.2cm, minimum height=1cm},
  branche/.style={rectangle, rounded corners, thick, align=center, minimum width=5.4cm, minimum height=1.1cm},
  fleche/.style={-{Stealth[length=3mm]}, thick}
]
\node[racine, align=center] (r) at (0,0){\textbf{Exprimer l'opinion}\\ \esp{...que + verbe}};
\node[branche, draw=popgreen, fill=popgreenL, align=center] (g) at (-4.6,-2.2){\textbf{Opinion affirmative}\\ \esp{creo que / pienso que}\\ $\Rightarrow$ \textbf{INDICATIF}};
\node[branche, draw=poporange, fill=poporangeL, align=center] (d) at (4.6,-2.2){\textbf{Opinion négative ou nécessité}\\ \esp{no creo que / es necesario que}\\ $\Rightarrow$ \textbf{SUBJONCTIF}};
\draw[fleche, popgreen] (r.south west) -- (g.north);
\draw[fleche, poporange] (r.south east) -- (d.north);
\end{tikzpicture}
\end{center}
\legende{Le grand partage : opinion affirmative $\rightarrow$ indicatif ; opinion négative ou nécessité impersonnelle $\rightarrow$ subjonctif.}
\end{popfigure}

\courssub{Les structures de revendication}

Revendiquer, c'est demander un changement. L'espagnol dispose de trois outils complémentaires.

\textbf{1. \esp{Hay que + infinitif} : la nécessité générale, sans sujet précis.}

\begin{propriete}[\esp{Hay que} + infinitif]
\esp{Hay que} (de l'impersonnel \esp{haber}) est suivi de l'\textbf{infinitif} et exprime une obligation générale qui ne vise personne en particulier : \esp{Hay que luchar contra la discriminación.} = « Il faut lutter contre la discrimination. » C'est l'équivalent exact du français « il faut + infinitif ». \textbf{Jamais de \esp{que} ni de subjonctif après \esp{hay que}.}
\end{propriete}

\textbf{2. \esp{Es necesario que / hace falta que + subjonctif} : la nécessité visant un sujet précis.}

Dès que vous voulez dire \textbf{QUI} doit agir, l'infinitif ne suffit plus : il faut une subordonnée au subjonctif. C'est le point de comparaison capital avec le français.

\begin{center}
\begin{tabularx}{\linewidth}{|X|X|X|}
\hline
\rowcolor{popblueL} \textbf{Français} & \textbf{Español} & \textbf{Remarque} \\
\hline
Il faut lutter. & \esp{Hay que luchar.} & Infinitif des deux côtés, sens général. \\
\hline
Il faut que je lutte. & \esp{Es necesario que luche.} / \esp{Hace falta que luche.} & Sujet précis $\rightarrow$ subjonctif obligatoire. \\
\hline
Il faut que nous changions les mentalités. & \esp{Es necesario que cambiemos las mentalidades.} & Le subjonctif espagnol rend le subjonctif français. \\
\hline
Il est important que les femmes votent. & \esp{Es importante que las mujeres voten.} & \esp{voten} = subjonctif présent de \esp{votar}. \\
\hline
\end{tabularx}
\end{center}

\textbf{3. Les verbes de revendication : \esp{exigir}, \esp{pedir}, \esp{reivindicar}.}

\begin{propriete}[Verbes de volonté + \esp{que} + subjonctif]
Les verbes exprimant une demande, une exigence ou une revendication — \esp{exigir que}, \esp{pedir que}, \esp{reivindicar que}, \esp{reclamar que}, \esp{querer que}, \esp{desear que} — sont suivis de \esp{que} + \textbf{subjonctif}, car le contenu demandé n'est pas encore réalisé. Attention à l'irrégularité de \esp{pedir} (changement \esp{e} $\rightarrow$ \esp{i}) : \esp{pido, pides, pide, pedimos, pedís, piden} (indicatif) ; \esp{pida, pidas, pida, pidamos, pidáis, pidan} (subjonctif).
\end{propriete}

\begin{exemplebox}[La revendication en action]
\begin{itemize}
\item \esp{Exigimos que el gobierno respete los derechos humanos.} = « Nous exigeons que le gouvernement respecte les droits humains. »
\item \esp{Las mujeres piden que se aplique la ley de igualdad.} = « Les femmes demandent que la loi sur l'égalité soit appliquée. » (passif réfléchi \esp{se aplique})
\item \esp{Los estudiantes reivindican que haya más becas.} = « Les étudiants revendiquent qu'il y ait plus de bourses. » (\esp{haya} = subjonctif de \esp{hay})
\item \esp{Hay que educar a los niños en el respeto.} = « Il faut éduquer les enfants dans le respect. »
\end{itemize}
\end{exemplebox}

\begin{attention}[Les deux erreurs classiques des francophones]
\begin{enumerate}
\item \textbf{Oublier le subjonctif après \esp{es necesario que}} : on entend souvent *\esp{es necesario que cambiamos las mentalidades} — \textbf{FAUX}. La forme correcte est \esp{es necesario que cambiemos} (subjonctif présent de \esp{cambiar}). Astuce : dès que \esp{es necesario que} apparaît, vérifiez immédiatement la voyelle thématique du verbe : \esp{-ar} $\rightarrow$ \esp{-e} (\esp{cambiemos}), \esp{-er/-ir} $\rightarrow$ \esp{-a} (\esp{coma, viva, haga, sea, vaya}).
\item \textbf{Mettre le subjonctif après \esp{hay que}} : *\esp{hay que luchemos} est impossible. \esp{Hay que} + \textbf{infinitif uniquement} : \esp{Hay que luchar.} Si vous voulez un sujet précis, passez à \esp{es necesario que luchemos}.
\end{enumerate}
\end{attention}

\courssub{Les connecteurs de l'argumentation}

Une opinion sans connecteurs est une liste ; avec eux, elle devient un raisonnement. Voici la boîte à outils indispensable du commentaire et de l'essai.

\begin{center}
\begin{tabularx}{\linewidth}{|X|X|X|}
\hline
\rowcolor{popblueL} \textbf{Conector} & \textbf{Français} & \textbf{Fonction} \\
\hline
\esp{en primer lugar} / \esp{primeramente} & en premier lieu, d'abord & introduire le premier argument \\
\hline
\esp{en segundo lugar} & en second lieu & structurer le deuxième point \\
\hline
\esp{además} / \esp{por añadidura} & de plus, en outre & ajouter un argument \\
\hline
\esp{sin embargo} / \esp{no obstante} & cependant, pourtant & concession, opposition \\
\hline
\esp{por eso} / \esp{por tanto} / \esp{por consiguiente} & c'est pourquoi, donc, par conséquent & conséquence \\
\hline
\esp{en conclusión} / \esp{en definitiva} / \esp{para terminar} & en conclusion, en définitive & conclure \\
\hline
\end{tabularx}
\end{center}

\begin{exemplebox}[Une mini-argumentation connectée]
\esp{En primer lugar, creo que el texto defiende la igualdad. Además, me parece que el autor denuncia la discriminación con ejemplos concretos. Sin embargo, no creo que las leyes sean suficientes. Por eso, es necesario que la sociedad entera se movilice. En conclusión, hay que seguir luchando.}

= « D'abord, je crois que le texte défend l'égalité. De plus, il me semble que l'auteur dénonce la discrimination avec des exemples concrets. Cependant, je ne crois pas que les lois soient suffisantes. C'est pourquoi il faut que la société entière se mobilise. En conclusion, il faut continuer à lutter. »
\end{exemplebox}

\begin{methode}[Nuancer son opinion au Bac en trois temps]
\begin{enumerate}
\item \textbf{Affirmez} avec une structure d'opinion + indicatif : \esp{Creo que el texto defiende los derechos de las mujeres.}
\item \textbf{Nuancez} avec \esp{sin embargo} + opinion négative + subjonctif : \esp{Sin embargo, no creo que la situación haya mejorado del todo.}
\item \textbf{Revendiquez} avec une structure impersonnelle ou un verbe de volonté + subjonctif : \esp{Por eso, es necesario que cambiemos las mentalidades.} / \esp{Exigimos que se respete la igualdad.}
\end{enumerate}
Ce triptyque \emph{opinion $\rightarrow$ concession ($\esp{sin embargo}$) $\rightarrow$ revendication} est la colonne vertébrale de tout paragraphe argumentatif réussi au commentaire de texte.
\end{methode}

\begin{exoresolu}[Indicatif ou subjonctif ?]
Complétez avec la forme correcte du verbe entre parenthèses, puis traduisez :
\begin{enumerate}
\item \esp{Creo que las minorías (tener) \ldots\ derechos.}
\item \esp{No pienso que la ley (ser) \ldots\ suficiente.}
\item \esp{Es necesario que nosotros (luchar) \ldots\ contra la discriminación.}
\item \esp{Hay que (respetar) \ldots\ a todas las personas.}
\item \esp{Las mujeres exigen que se les (pagar) \ldots\ el mismo salario.}
\end{enumerate}
\tcblower
\textbf{Solutions détaillées :}
\begin{enumerate}
\item \esp{Creo que las minorías \textbf{tienen} derechos.} — Opinion affirmative (\esp{creo que}) $\rightarrow$ indicatif présent (\esp{tener} $\rightarrow$ \esp{tienen}). « Je crois que les minorités ont des droits. »
\item \esp{No pienso que la ley \textbf{sea} suficiente.} — Opinion négative (\esp{no pienso que}) $\rightarrow$ subjonctif présent de \esp{ser} (irrégulier : \esp{sea, seas, sea, seamos, seáis, sean}). « Je ne pense pas que la loi soit suffisante. »
\item \esp{Es necesario que nosotros \textbf{luchemos} contra la discriminación.} — Nécessité impersonnelle (\esp{es necesario que}) $\rightarrow$ subjonctif présent de \esp{luchar} (verbe en \esp{-ar} $\rightarrow$ \esp{-emos} : \esp{luchemos}). « Il faut que nous luttions contre la discrimination. »
\item \esp{Hay que \textbf{respetar} a todas las personas.} — \esp{Hay que} + \textbf{infinitif}, jamais de subjonctif. « Il faut respecter toutes les personnes. »
\item \esp{Las mujeres exigen que se les \textbf{pague} el mismo salario.} — Verbe de revendication (\esp{exigen que}) $\rightarrow$ subjonctif présent de \esp{pagar}. Changement orthographique \esp{g} $\rightarrow$ \esp{gu} devant \esp{e} pour garder le son [g] : \esp{pague}. « Les femmes exigent qu'on leur paie le même salaire. » (Passif réfléchi impersonnel \esp{se les pague}).
\end{enumerate}
\textbf{Clé de vérification :} structure affirmative $\rightarrow$ indicatif ; structure négative, impersonnelle de nécessité ou verbe de volonté $\rightarrow$ subjonctif ; \esp{hay que} $\rightarrow$ infinitif.
\end{exoresolu}

\begin{savaistu}[Droits des femmes : regards croisés Cameroun -- Guinée équatoriale]
La \textbf{Guinée équatoriale}, seul pays africain hispanophone, a ratifié la \emph{Convención sobre la Eliminación de todas las Formas de Discriminación contra la Mujer} (CEDAW) en 1984. Le \emph{Plan Nacional de Desarrollo Horizonte 2020} incluait des objectifs d'égalité, mais les ONG (comme \emph{Mujeres de Guinea Ecuatorial}) dénoncent la persistance des violences basées sur le genre et la sous-représentation politique.
Au \textbf{Cameroun}, la Constitution de 1996 (révisée 2008) garantit l'égalité (art. 1$^{\text{er}}$). La loi n°2011/011 sur la cybercriminalité punit le harcèlement en ligne. Pourtant, le Code de la famille (ordonnance 81-02) contient encore des dispositions discriminatoires (autorisation maritale pour le travail, gestion des biens). Le \emph{MINPROFF} (Ministère de la Promotion de la Femme et de la Famille) pilote la \emph{Stratégie Nationale de Lutte contre les Violences Basées sur le Genre}. Le chemin vers l'\esp{igualdad real} reste long des deux côtés de la frontière.
\end{savaistu}

\begin{experience}[Débat oral : « La place des femmes dans la société camerounaise »]
\textbf{Consigne :} En groupes de 4, préparez un débat de 5 minutes. Deux élèves défendent l'idée que « La situation des femmes s'est nettement améliorée au Cameroun », deux autres soutiennent que « Les inégalités structurelles persistent ».
\textbf{Contraintes linguistiques :}
\begin{itemize}
\item Utilisez au moins 3 connecteurs du tableau (\esp{en primer lugar, además, sin embargo, por eso, en conclusión}).
\item Intégrez au moins 2 opinions affirmatives (indicatif), 2 opinions négatives (subjonctif), 1 structure \esp{hay que + infinitif} et 1 structure \esp{es necesario que + subjonctif}.
\item Employez 5 mots du lexique de la leçon (\esp{discriminación, sufragio, igualdad de oportunidades, ley, reivindicar, minoría}).
\end{itemize}
\textbf{Exemple d'amorce :} \esp{``En primer lugar, creo que la ley protege a las mujeres. Sin embargo, no pienso que se aplique en el mundo rural. Por eso, hay que educar a las familias...''}
\end{experience}

\begin{aretenir}[L'essentiel de la section]
\begin{itemize}
\item Opinion \textbf{affirmative} (\esp{creo que, pienso que, me parece que, considero que, en mi opinión}) $\rightarrow$ \textbf{indicatif}.
\item Opinion \textbf{négative} (\esp{no creo que, no pienso que, no me parece que}) $\rightarrow$ \textbf{subjonctif}.
\item Nécessité impersonnelle (\esp{es necesario que, es importante que, hace falta que, es imprescindible que}) $\rightarrow$ \textbf{subjonctif}.
\item \esp{Hay que} $\rightarrow$ \textbf{infinitif} uniquement (jamais de \esp{que}).
\item Verbes de revendication (\esp{exigir, pedir, reivindicar, reclamar, querer}) + \esp{que} $\rightarrow$ \textbf{subjonctif}.
\item Connecteurs : \esp{en primer lugar, además, sin embargo, por eso, en conclusión}.
\item Stratégie Bac : opinion $\rightarrow$ concession (\esp{sin embargo}) $\rightarrow$ revendication.
\end{itemize}
\end{aretenir}

\begin{exercice}[À vous de revendiquer]
Traduisez en espagnol :
\begin{enumerate}
\item Je pense que l'égalité des chances est un droit fondamental.
\item Je ne crois pas que la discrimination ait disparu de nos sociétés.
\item Il faut que le gouvernement protège les minorités.
\item Il faut lutter chaque jour contre les préjugés.
\item Nous demandons que les femmes aient les mêmes salaires que les hommes.
\end{enumerate}
Puis rédigez quatre phrases de votre cru sur les droits des femmes au Cameroun, en utilisant au moins une opinion affirmative, une opinion négative, une structure impersonnelle et un connecteur de concession (\esp{sin embargo} ou \esp{no obstante}).
\end{exercice}

\begin{corrige}[Exercice « À vous de revendiquer »]
\textbf{Traductions :}
\begin{enumerate}
\item \esp{Creo que la igualdad de oportunidades es un derecho fundamental.} (Opinion affirmative $\rightarrow$ indicatif \esp{es}).
\item \esp{No creo que la discriminación haya desaparecido de nuestras sociedades.} (Opinion négative $\rightarrow$ subjonctif parfait \esp{haya desaparecido}).
\item \esp{Es necesario que el gobierno proteja a las minorías.} (Nécessité impersonnelle $\rightarrow$ subjonctif \esp{proteja} de \esp{proteger}).
\item \esp{Hay que luchar cada día contra los prejuicios.} (\esp{Hay que} + infinitif).
\item \esp{Pedimos / Exigimos / Reivindicamos que las mujeres tengan los mismos salarios que los hombres.} (Volonté $\rightarrow$ subjonctif \esp{tengan} de \esp{tener}, irrégulier).
\end{enumerate}

\textbf{Production écrite exemple (non unique) :}
\esp{En mi opinión, las mujeres camerunesas han ganado más visibilidad en la política y la empresa. Sin embargo, no creo que la igualdad salarial sea una realidad en el sector informal. Por eso, es necesario que el Estado fortalezca la ley laboral. Hay que seguir reivindicando los derechos de las madres trabajadoras.}
\end{corrige}

\coursec{Méthode du commentaire de texte}

Dans la section précédente, nous avons appris à exprimer l'opinion et la revendication avec des structures comme \esp{En mi opinión}, \esp{Considero que} ou \esp{Es necesario reivindicar}. Ces outils vont maintenant nous servir dans un exercice complet et décisif au baccalauréat : le \cle{commentaire de texte} (\esp{comentario de texto}). En effet, commenter un texte, c'est d'abord le comprendre, puis l'analyser, et enfin donner son point de vue personnel : exactement ce que nous venons de nous entraîner à faire.

\courssub{Nature de l'épreuve au baccalauréat}

\begin{definition}[Qu'est-ce qu'un comentario de texto ?]
Le \cle{commentaire de texte} est un exercice écrit en espagnol qui consiste à :
\begin{itemize}
\item \textbf{comprendre} le texte : identifier son thème, sa source, ses idées principales ;
\item \textbf{analyser} le texte : dégager les idées directrices, étudier le point de vue de l'auteur, son ton et ses arguments ;
\item \textbf{donner son point de vue} : exprimer un avis personnel argumenté, en espagnol correct.
\end{itemize}
Au baccalauréat, le commentaire porte généralement sur un texte de presse ou un texte littéraire lié au programme, comme notre texte sur \esp{los derechos de las mujeres y de las minorías}.
\end{definition}

\begin{attention}[Ce qu'un commentaire n'est pas]
Un commentaire n'est \textbf{ni un résumé ni une traduction}. Beaucoup d'élèves francophones se contentent de répéter ce que dit le texte avec d'autres mots : c'est insuffisant. Il faut \textbf{analyser}, c'est-à-dire répondre aux questions \emph{pourquoi} et \emph{comment} : pourquoi l'auteur écrit-il cela ? Comment s'y prend-il pour convaincre ? Quel est son ton ? Rapporter ce que dit le texte sans l'expliquer, c'est faire un résumé, et le résumé est sanctionné au Bac.
\end{attention}

\courssub{Le plan du commentaire en quatre étapes}

Le commentaire de texte suit toujours la même architecture. Mémorisez cet organigramme : c'est votre feuille de route le jour de l'examen.

\begin{popfigure}
\begin{tikzpicture}[
node distance=0.55cm,
etape/.style={rectangle, rounded corners=4pt, draw=popblue, fill=popblueL, text width=10.5cm, align=center, inner sep=6pt, font=\small},
fleche/.style={-{Stealth[length=3mm]}, thick, poporange}
]
\node[etape] (intro) {\textbf{1. Introducción} : presentación del texto (título, fuente, tema) + problemática};
\node[etape, below=of intro] (des) {\textbf{2. Desarrollo} : 2 o 3 ideas directrices, ilustradas con citas traducidas};
\node[etape, below=of des] (autor) {\textbf{3. Opinión del autor} : tono (comprometido, crítico, optimista), argumentos, procedimientos};
\node[etape, below=of autor] (conc) {\textbf{4. Conclusión} : balance + opinión personal argumentada};
\draw[fleche] (intro) -- (des);
\draw[fleche] (des) -- (autor);
\draw[fleche] (autor) -- (conc);
\end{tikzpicture}
\legende{Organigramme du plan de commentaire de texte}
\end{popfigure}

\courssub{Étape 1 : l'introduction}

L'introduction est courte (trois à cinq lignes). Elle doit présenter le texte et annoncer la \cle{problématique} (\esp{problemática}), c'est-à-dire la grande question que pose le texte.

\begin{methode}[Rédiger l'introduction en trois phrases]
\begin{enumerate}
\item \textbf{Présenter le texte} : \esp{Este texto, titulado « ... », es un artículo de prensa / un fragmento de novela, publicado en ...} (Ce texte, intitulé « ... », est un article de presse / un extrait de roman, publié dans ...)
\item \textbf{Annoncer le thème} : \esp{El texto trata de los derechos de las mujeres y de las minorías.} (Le texte traite des droits des femmes et des minorités.)
\item \textbf{Formuler la problématique} : \esp{El autor plantea el problema de la discriminación y se pregunta cómo conseguir la igualdad de oportunidades.} (L'auteur pose le problème de la discrimination et se demande comment obtenir l'égalité des chances.)
\end{enumerate}
\end{methode}

\begin{attention}[Faux amis dans l'introduction]
\begin{itemize}
\item \esp{tratar de} + nom = « traiter de » ; mais \esp{tratar de} + infinitif = « essayer de ». Ne confondez pas : \esp{El texto trata de la igualdad} (le texte traite de l'égalité) / \esp{Trata de comprender el texto} (essaie de comprendre le texte).
\item \esp{plantear} = « poser (un problème) », et non « planter ».
\item N'écrivez jamais \esp{el texto es sobre} : dites \esp{el texto trata de} ou \esp{el texto habla de}.
\end{itemize}
\end{attention}

\courssub{Étape 2 : le développement — les idées directrices}

Le développement est le cœur du commentaire. Vous devez dégager \textbf{deux ou trois idées directrices} (\esp{ideas directrices}) et les illustrer par des \textbf{citations} (\esp{citas}) tirées du texte, puis traduites ou expliquées.

\begin{propriete}[Règle de la citation]
Une citation doit être :
\begin{itemize}
\item \textbf{courte} : quelques mots ou une phrase, entre guillemets ;
\item \textbf{exacte} : recopiée fidèlement, avec les accents ;
\item \textbf{intégrée} dans votre phrase : \esp{El autor afirma que « las mujeres siguen ganando menos que los hombres ».} ;
\item \textbf{exploitée} : chaque citation doit être suivie d'une explication ou d'une traduction. Une citation jetée sans commentaire ne vaut rien.
\end{itemize}
\end{propriete}

Pour annoncer chaque idée, utilisez des connecteurs d'ordre :

\begin{center}
\begin{tabularx}{\linewidth}{|X|X|X|}
\hline
\rowcolor{popblueL} Español & Français & Emploi \\
\hline
\esp{En primer lugar} & En premier lieu & première idée \\
\hline
\esp{En segundo lugar} & En second lieu & deuxième idée \\
\hline
\esp{Además} & De plus & ajouter un argument \\
\hline
\esp{Por ejemplo} & Par exemple & introduire une citation \\
\hline
\esp{Sin embargo} & Cependant & nuancer, opposer \\
\hline
\esp{Por último} & Enfin & dernière idée \\
\hline
\end{tabularx}
\end{center}

\courssub{Étape 3 : le point de vue de l'auteur}

Après les idées, il faut analyser \textbf{comment} l'auteur défend sa cause. Trois questions vous guident :

\begin{itemize}
\item \textbf{Le ton} (\esp{el tono}) : est-il \esp{comprometido} (engagé), \esp{crítico} (critique), \esp{optimista} (optimiste), \esp{pesimista} (pessimiste), \esp{irónico} (ironique) ?
\item \textbf{Les arguments} (\esp{los argumentos}) : l'auteur utilise-t-il des faits, des exemples, des comparaisons, des questions rhétoriques ?
\item \textbf{Les procédés} (\esp{los procedimientos}) : emploi du « nous » inclusif (\esp{debemos luchar}), du lexique de la revendication (\esp{reivindicar, denunciar, exigir}), de la répétition, des contrastes (\esp{por un lado... por otro lado}).
\end{itemize}

\begin{exemplebox}[Phrases utiles pour analyser le point de vue de l'auteur]
\esp{El autor denuncia la discriminación que sufren las mujeres.} « L'auteur dénonce la discrimination que subissent les femmes. »

\esp{El tono del texto es comprometido y optimista.} « Le ton du texte est engagé et optimiste. »

\esp{Para convencer al lector, el autor utiliza preguntas retóricas.} « Pour convaincre le lecteur, l'auteur utilise des questions rhétoriques. »

\esp{El autor defiende la igualdad de oportunidades para todos.} « L'auteur défend l'égalité des chances pour tous. »
\end{exemplebox}

\courssub{Étape 4 : la conclusion}

La conclusion fait le \textbf{bilan} (\esp{balance}) de l'analyse en une ou deux phrases, puis donne votre \textbf{avis personnel argumenté}. C'est ici que vous réemployez les structures d'opinion de la section précédente :

\begin{center}
\begin{tabularx}{\linewidth}{|X|X|}
\hline
\rowcolor{popblueL} Español & Français \\
\hline
\esp{En conclusión, este texto defiende...} & En conclusion, ce texte défend... \\
\hline
\esp{En mi opinión, este texto es muy actual.} & À mon avis, ce texte est très actuel. \\
\hline
\esp{Me parece que la lucha por la igualdad debe continuar.} & Il me semble que la lutte pour l'égalité doit continuer. \\
\hline
\esp{Estoy de acuerdo con el autor porque...} & Je suis d'accord avec l'auteur parce que... \\
\hline
\esp{Considero que todavía queda mucho por hacer.} & Je considère qu'il reste encore beaucoup à faire. \\
\hline
\end{tabularx}
\end{center}

\begin{attention}[Pièges de la conclusion]
\begin{itemize}
\item Après \esp{me parece que}, \esp{creo que}, \esp{pienso que} (avis positif), on utilise l'\textbf{indicatif} : \esp{Creo que el texto \textbf{es} interesante.} Le subjonctif n'apparaît qu'à la forme négative : \esp{No creo que el problema \textbf{esté} resuelto.}
\item \esp{actual} signifie « actuel, d'actualité », et non « réel ». Pour « réel », dites \esp{real}.
\item Ne terminez jamais par une simple phrase comme \esp{Y ya está.} : la conclusion doit être rédigée et soignée.
\end{itemize}
\end{attention}

\courssub{Commentaire modèle, commenté pas à pas}

Reprenons le texte support de la leçon. En voici un extrait pour mémoire :

\begin{textebox}[Los derechos de las mujeres: una lucha que continúa]
En muchos países del mundo, las mujeres siguen luchando por sus derechos. Aunque las leyes reconocen la igualdad, la realidad es diferente: en el trabajo, en la familia y en la política, las mujeres sufren todavía discriminación. Las minorías, por su parte, reclaman el respeto de su cultura y de su lengua. Debemos reivindicar la igualdad de oportunidades para todos, porque los derechos humanos no tienen fronteras ni género. La lucha continúa, y cada voz cuenta.
\end{textebox}

\textbf{Glossaire} : \esp{seguir} + gérondif = continuer à ; \esp{aunque} = bien que ; \esp{reclamar} = réclamer ; \esp{por su parte} = de leur côté ; \esp{cada voz cuenta} = chaque voix compte.

Voici maintenant un commentaire modèle, rédigé comme au baccalauréat, avec en marge l'analyse de chaque étape.

\begin{exemplebox}[Comentario modelo]
\textbf{Introduction} — \esp{Este texto, titulado « Los derechos de las mujeres: una lucha que continúa », es un artículo de opinión. El texto trata de los derechos de las mujeres y de las minorías. El autor plantea el problema de la discriminación y se pregunta cómo conseguir una verdadera igualdad.}

\emph{(Trois phrases : présentation, thème, problématique. Le plan est respecté.)}

\textbf{Développement} — \esp{En primer lugar, el autor denuncia la discriminación que sufren las mujeres: afirma que « en el trabajo, en la familia y en la política, las mujeres sufren todavía discriminación », es decir, que la ley no basta para cambiar la realidad. En segundo lugar, el texto defiende también los derechos de las minorías, que « reclaman el respeto de su cultura y de su lengua ». Por último, el autor lanza un llamamiento: « debemos reivindicar la igualdad de oportunidades para todos ».}

\emph{(Trois idées directrices, chacune illustrée par une citation courte, exacte et expliquée : \esp{es decir} introduit la traduction-explication.)}

\textbf{Point de vue de l'auteur} — \esp{El tono del texto es comprometido pero optimista. El autor utiliza el « nosotros » inclusivo (« debemos reivindicar ») para asociar al lector a la lucha, y termina con una frase breve y esperanzadora: « cada voz cuenta ».}

\emph{(Analyse du ton et de deux procédés : le « nous » inclusif et la chute brève. On répond au « comment », pas seulement au « quoi ».)}

\textbf{Conclusion} — \esp{En conclusión, este texto denuncia la discriminación y defiende la igualdad de oportunidades. En mi opinión, este texto es muy actual, porque en mi país las mujeres y las minorías todavía luchan por sus derechos. Considero que la educación es la mejor arma para cambiar las mentalidades.}

\emph{(Bilan en une phrase, puis avis personnel argumenté avec \esp{en mi opinión} et \esp{considero que} + indicatif.)}
\end{exemplebox}

\begin{savaistu}[Le saviez-tu ?]
Au baccalauréat, le correcteur valorise toute \textbf{citation du texte correctement intégrée et traduite} dans le commentaire. Une citation bien choisie prouve que vous avez réellement lu et compris le texte. À l'inverse, un commentaire sans aucune citation donne l'impression que l'élève parle « à côté » du texte. Entraînez-vous donc à citer au moins trois passages dans chaque commentaire.
\end{savaistu}

\begin{exoresolu}[Construire le squelette d'un commentaire]
\textbf{Énoncé} — À partir du texte support \esp{« Los derechos de las mujeres: una lucha que continúa »}, complétez le squelette de commentaire suivant avec les éléments manquants : 1) la problématique, 2) une citation pour l'idée « las minorías », 3) le ton du texte, 4) une phrase d'avis personnel.

\esp{Este texto trata de los derechos de las mujeres y de las minorías. El autor plantea el problema de ... (1)}

\esp{En segundo lugar, el texto defiende los derechos de las minorías: « ... » (2)}

\esp{El tono del texto es ... (3)}

\esp{En mi opinión, ... (4)}

\tcblower
\textbf{Solution détaillée}

1) \esp{El autor plantea el problema de la discriminación y de la desigualdad entre hombres y mujeres.} — La problématique reprend le thème sous forme de problème à résoudre.

2) \esp{« reclaman el respeto de su cultura y de su lengua »} — Citation courte, exacte, entre guillemets, directement liée à l'idée annoncée.

3) \esp{El tono del texto es comprometido y optimista.} — Deux adjectifs de ton justifiés par le texte : \esp{debemos reivindicar} (engagé) et \esp{cada voz cuenta} (optimiste).

4) \esp{En mi opinión, este texto es muy actual porque la igualdad todavía no es una realidad en todas partes.} — Avis personnel avec \esp{en mi opinión} + indicatif, et une justification introduite par \esp{porque}.
\end{exoresolu}

\begin{exercice}[À vous de commenter]
Rédigez en espagnol un commentaire complet (15 à 20 lignes) du texte support \esp{« Los derechos de las mujeres: una lucha que continúa »}, en suivant le plan en quatre étapes. Vous citerez au moins trois passages du texte et vous terminerez par un avis personnel argumenté.
\end{exercice}

\begin{corrige}[Exercice : À vous de commenter]
Éléments de correction (commentaire possible) :

\textbf{Introducción} : \esp{Este texto, titulado « Los derechos de las mujeres: una lucha que continúa », es un artículo de opinión. Trata de los derechos de las mujeres y de las minorías. El autor plantea el problema de la discriminación y se pregunta cómo alcanzar la igualdad de oportunidades.}

\textbf{Desarrollo} : \esp{En primer lugar, el autor denuncia que « las mujeres sufren todavía discriminación » en el trabajo, en la familia y en la política, aunque las leyes reconocen la igualdad. En segundo lugar, defiende a las minorías, que « reclaman el respeto de su cultura y de su lengua ». Por último, lanza un llamamiento a la acción: « debemos reivindicar la igualdad de oportunidades para todos ».}

\textbf{Opinión del autor} : \esp{El tono es comprometido y optimista. El autor utiliza el « nosotros » inclusivo y el lexico de la reivindicación (« reivindicar », « luchar ») para movilizar al lector.}

\textbf{Conclusión} : \esp{En conclusión, este texto defiende los derechos humanos de todos. En mi opinión, es un texto muy actual, porque la discriminación existe todavía en muchos países. Considero que la educación y la ley deben ir juntas para cambiar la sociedad.}

Barème indicatif : respect du plan (4 points), citations exactes et exploitées (6 points), correction linguistique (6 points), avis personnel argumenté (4 points).
\end{corrige}

\begin{aretenir}[À retenir]
\begin{itemize}
\item Le commentaire de texte = comprendre + analyser + donner son avis. Jamais de simple résumé.
\item Plan en quatre étapes : \esp{introducción} (présentation + problématique), \esp{desarrollo} (2 ou 3 idées + citations traduites), \esp{opinión del autor} (ton, arguments, procédés), \esp{conclusión} (bilan + avis personnel).
\item Toute citation doit être courte, exacte, intégrée et expliquée.
\item Avis affirmatif : \esp{creo que} + indicatif ; avis négatif : \esp{no creo que} + subjonctif.
\end{itemize}
\end{aretenir}

\coursec{Production guidée : un texto reivindicativo}

Après avoir analysé la structure argumentative d'un texte militant dans la section précédente (\emph{Méthode du commentaire de texte}), nous passons maintenant à l'action : vous allez produire votre propre \cle{texto reivindicativo} (texte revendicatif). L'objectif est de mobiliser le lexique des droits humains, les structures d'opinion et le subjonctif de revendication pour défendre une cause concrète, dans le respect des attendus du Baccalauréat (épreuve d'expression écrite).

\courssub{La tâche : consignes et contraintes}

\begin{definition}[Texto reivindicativo]
Un \textbf{texto reivindicativo} est un écrit à visée persuasive et militante. Il dénonce une injustice, propose une solution et appelle à l'action collective. Il s'appuie sur des faits, des valeurs universelles (droits humains) et un lexique précis (\cle{reivindicar}, \cle{exigir}, \cle{denunciar}, \cle{igualdad}).
\end{definition}

\begin{propriete}[Contraintes formelles de l'exercice]
Pour valider cette production, votre texte doit respecter \textbf{impérativement} les critères suivants :
\begin{itemize}
    \item \textbf{Longueur :} 80 à 120 mots (hors salutations formelles si lettre).
    \item \textbf{Lexique obligatoire :} minimum \textbf{6 mots} de la liste de la leçon (\cle{derechos humanos}, \cle{minoría}, \cle{discriminación}, \cle{sufragio}, \cle{igualdad de oportunidades}, \cle{ley}, \cle{reivindicar}, \cle{excluido}, \cle{brecha}, \cle{empoderamiento}).
    \item \textbf{Structures d'opinion (2 minimum) :} \esp{En mi opinión...}, \esp{Considero que...}, \esp{Me parece inadmisible que...}, \esp{Estoy convencido/a de que...}.
    \item \textbf{Structure de revendication avec subjonctif (1 minimum) :} \esp{Exigimos que el gobierno garantice...}, \esp{Es necesario que se aplique la ley...}, \esp{Reivindicamos que se respeten...}.
    \item \textbf{Plan imposé :} Problème $\rightarrow$ Exemple concret $\rightarrow$ Revendication $\rightarrow$ Appel à l'action.
\end{itemize}
\end{propriete}

\courssub{Banque de phrases amorces et structures utiles}

Avant de rédiger, observons les « briques » syntaxiques qui donneront de la force à votre texte. Le passage de l'indicatif (constat) au subjonctif (volonté/exigence) est la marque de l'écrit revendicatif abouti.

\begin{center}
\begin{tabularx}{\linewidth}{|X|X|X|}
\hline
\rowcolor{popblueL}
\textbf{Fonction} & \textbf{Amorce / Structure (Español)} & \textbf{Traduction / Note grammaticale} \\
\hline
\cle{Constater / Dénoncer} & \esp{En nuestra sociedad, todavía existe la discriminación contra...} & « Dans notre société, la discrimination contre... existe encore. » (Indicatif : réalité) \\
\hline
 & \esp{Es inaceptable que las mujeres ganen menos que los hombres \textbf{por el mismo trabajo}.} & « Il est inacceptable que les femmes gagnent moins... » (\cle{Es + adj + que + SUBJONCTIF}) \\
\hline
\cle{Exprimer son opinion} & \esp{En mi opinión, la brecha salarial refleja una injusticia estructural.} & « À mon avis, l'écart salarial reflète une injustice structurelle. » \\
\hline
 & \esp{Considero que las minorías étnicas (como los mbororos o los bagyelis) están excluidas.} & « Je considère que les minorités ethniques... sont exclues. » (\cle{Considero que + INDICATIF}) \\
\hline
 & \esp{Me parece fundamental garantizar la igualdad de oportunidades.} & « Il me semble fondamental de garantir l'égalité des chances. » (\cle{Me parece + adj + INF}) \\
\hline
\cle{Revendiquer / Exiger (SUBJ)} & \esp{\textbf{Reivindicamos} que se \textbf{respeten} los derechos de todas las minorías.} & « Nous revendiquons que les droits de toutes les minorités soient respectés. » \\
\hline
 & \esp{\textbf{Exigimos} que las \textbf{leyes} se \textbf{apliquen} sin distinción.} & « Nous exigeons que les lois s'appliquent sans distinction. » \\
\hline
 & \esp{\textbf{Es urgente que} el Estado \textbf{garantice} el sufragio efectivo para todos.} & « Il est urgent que l'État garantisse le suffrage effectif pour tous. » \\
\hline
\cle{Appeler à l'action} & \esp{¡Unámonos para construir una sociedad más justa!} & « Unissons-nous pour construire une société plus juste ! » (Impératif \emph{nosotros}) \\
\hline
 & \esp{Es hora de actuar: firmemos la petición / manifestémonos pacíficamente.} & « Il est temps d'agir : signons la pétition / manifestons pacifiquement. » \\
\hline
\end{tabularx}
\end{center}

\begin{aretenir}[Le piège du subjonctif après « Es + adjectif + que »]
\emph{Siempre} subjonctif si l'adjectif exprime une nécessité, une émotion, un jugement de valeur, un doute ou une subjectivité : \cle{Es necesario que}, \cle{Es urgente que}, \cle{Es injusto que}, \cle{Es fundamental que}, \cle{Me parece mal que}.
\par\medskip
\esp{Es necesario que la ley \textbf{proteja} (subj.) a las víctimas.} $\neq$ \esp{Es verdad que la ley \textbf{protege} (ind.) a las víctimas.}
\end{aretenir}

\courssub{Plan guidé pas à pas : de l'idée au paragraphe}

Suivez ce canevas pour structurer votre brouillon. Chaque étape correspond à un paragraphe (ou demi-paragraphe) de votre texte final.

\begin{methode}[Rédiger un texto reivindicativo en 4 étapes]
\begin{enumerate}
    \item \textbf{Présenter le problème (Gancho + Contexto).}
    \begin{itemize}
        \item Accroche : \esp{En el Día Internacional de la Mujer / En nuestra comunidad...}
        \item Constat choc avec lexique : \esp{la discriminación}, \esp{la brecha}, \esp{la exclusión}, \esp{los derechos humanos}.
        \item \emph{Connecteurs :} \esp{Actualmente}, \esp{Por desgracia}, \esp{A pesar de los avances...}
    \end{itemize}
    \item \textbf{Illustrer par un exemple concret (Evidencia).}
    \begin{itemize}
        \item Contexte Cameroun : \esp{En el Norte de Camerún, muchas mujeres mbororos no acceden a la tierra.}
        \item Contexte Monde hispanophone : \esp{En América Latina, el feminicidio sigue siendo una lacra.}
        \item \emph{Connecteurs :} \esp{Por ejemplo}, \esp{Un caso claro es...}, \esp{En concreto...}
    \end{itemize}
    \item \textbf{Formuler la revendication centrale (Exigencia + Subjuntivo).}
    \begin{itemize}
        \item Verbe fort : \esp{Reivindicamos}, \esp{Exigimos}, \esp{Demandamos}.
        \item Structure : \esp{que + SUJET + SUBJONCTIF PRÉSENT}.
        \item Contenu : \esp{igualdad de oportunidades}, \esp{aplicación de la ley}, \esp{fin de la impunidad}.
    \end{itemize}
    \item \textbf{Lancer l'appel à l'action (Movilización).}
    \begin{itemize}
        \item Impératif \emph{nosotros} ou \emph{ustedes} / Infinitif d'affiche.
        \item \esp{¡Luchemos juntos!}, \esp{Firmemos la carta}, \esp{Exijamos justicia ya!}
    \end{itemize}
\end{enumerate}
\end{methode}

\courssub{Mise en œuvre : modèle de début de production}

Voici un modèle respectant exactement le plan et les contraintes (lexique, subjonctif, opinion). Utilisez-le comme « amorce » pour votre propre sujet.

\begin{textebox}[Modèle de début de production — Sujet : Journée des droits des femmes]
En nuestra sociedad, \textbf{todavía existe la discriminación} contra las mujeres y las \textbf{minorías}. \textbf{En mi opinión}, la \textbf{brecha salarial} y la violencia de género violan los \textbf{derechos humanos} fundamentales. \textbf{Por ejemplo}, en muchas zonas rurales de Camerún, las mujeres no tienen \textbf{igualdad de oportunidades} para acceder a la propiedad de la tierra. \textbf{Por eso, reivindicamos} que se \textbf{respete} la \textbf{ley} y que el Estado \textbf{garantice} la justicia real. \textbf{Es necesario que} la sociedad civil se \textbf{movilice}. \textbf{¡Unámonos para exigir igualdad ya!}
\end{textebox}

\begin{center}
\begin{tabularx}{\linewidth}{|X|X|}
\hline
\rowcolor{popgreenL}
\textbf{Élément requis} & \textbf{Réalisation dans le modèle} \\
\hline
Lexique (6+) & \esp{discriminación, minorías, brecha salarial, derechos humanos, igualdad de oportunidades, ley, reivindicamos, garantice, movilice} (9 mots) \\
\hline
Opinion (2+) & \esp{En mi opinión...} ; \esp{Es necesario que...} (impersonnel = opinion forte) \\
\hline
Revendication + Subj. & \esp{reivindicamos que se respete} ; \esp{Exigimos que el Estado garantice} (2 occurrences) \\
\hline
Plan (4 étapes) & 1. Problème (L1-2) $\rightarrow$ 2. Exemple Cameroun (L3-4) $\rightarrow$ 3. Revendication (L4-5) $\rightarrow$ 4. Appel (L6) \\
\hline
Nombre de mots & $\approx$ 85 mots (dans la fourchette 80-120) \\
\hline
\end{tabularx}
\end{center}

\courssub{Technique d'enrichissement : la règle du « 1+1+1 »}

Un paragraphe argumentatif solide ne se contente pas d'aligner des phrases. Il articule une idée, une preuve et un lien logique.

\begin{methode}[Enrichir un écrit : 1 idée + 1 exemple + 1 connecteur = 1 paragraphe efficace]
\begin{enumerate}
    \item \textbf{Idée directrice (Topic sentence) :} Une phrase claire qui annonce le thème du paragraphe. \par \esp{La discriminación laboral frena el desarrollo del país.}
    \item \textbf{Exemple / Preuve (Evidencia) :} Un fait précis, un chiffre (approximatif), une référence géographique ou légale. \par \esp{Según la OIT, las mujeres ganan un 20\% menos en el sector agrícola camerunés.}
    \item \textbf{Connecteur logique (Conector) :} Montre la relation (cause, conséquence, opposition, but). \par \esp{Por consiguiente / Por eso / Por tanto}, \textbf{reivindicamos} que se \textbf{aplique} la \textbf{ley} de igualdad salarial.
    \item \textbf{Résultat :} Un paragraphe cohérent, dense et persuasif.
\end{enumerate}
\end{methode}

\courssub{Vigilance grammaticale : les accords au féminin pluriel}

C'est l'erreur n°1 des francophones : en français, « les femmes discriminées » (accord), en espagnol, l'adjectif s'accorde en genre ET en nombre avec le nom qu'il qualifie.

\begin{attention}[Accords des adjectifs avec \esp{mujeres} et \esp{minorías} — Piège majeur]
\esp{Mujer} et \esp{minoría} sont des noms \textbf{féminins}. Au pluriel : \esp{las mujeres}, \esp{las minorías}. \textbf{Tout adjectif ou participe passé qui s'y rapporte DOIT être au féminin pluriel (\cle{-as})}.
\par\medskip
\begin{center}
\begin{tabular}{|l|l|l|}
\hline
\rowcolor{poporangeL}
\textbf{Français (piège)} & \textbf{Espagnol CORRECT} & \textbf{Erreur fréquente} \\
\hline
Les femmes discriminées & \esp{Las mujeres \textbf{discriminadas}} & \esp{*Las mujeres discriminad\emph{os}} \\
\hline
Les minorités exclues & \esp{Las minorías \textbf{excluidas}} & \esp{*Las minorías exclu\emph{idos}} \\
\hline
Les femmes vulnérables & \esp{Las mujeres \textbf{vulnerables}} (invariable) & \esp{*Las mujeres vulnerables\emph{s}} \\
\hline
Les lois appliquées & \esp{Las leyes \textbf{aplicadas}} & \esp{*Las leyes aplicad\emph{os}} \\
\hline
\end{tabular}
\end{center}
\par\medskip
\emph{Astuce :} Entourez le nom noyau (\esp{mujeres}, \esp{minorías}, \esp{leyes}) et faites « rebondir » la finale \textbf{-as} sur tous les mots du groupe nominal (déterminants, adjectifs, participes).
\end{attention}

\courssub{Exercice résolu : du brouillon à la version finale}

Sujet : \emph{« Escribe un texto reivindicativo (90-110 palabras) para el "Día Internacional de las Minorías" denunciando la exclusión de los pueblos indígenas en América Latina y exigiendo sus derechos territoriales. »}

\begin{exoresolu}[Rédaction guidée : Día de las Minorías]
\textbf{Étape 1 : Brouillon (Idées brutes + Lexique ciblé)}
\begin{itemize}
    \item \cle{Lexique choisi :} \esp{minorías indígenas, discriminación histórica, tierras ancestrales, derechos humanos, reivindicar, ley, consulta previa, exclusión}.
    \item \cle{Structures :} \esp{Considero que...} (opinion), \esp{Exigimos que se respete...} (revendication + subj).
    \item \cle{Plan :} 1. Constat historique. 2. Exemple (Amazone/Guatemala). 3. Exigence légale (Consulta previa). 4. Appel.
\end{itemize}

\textbf{Étape 2 : Rédaction du premier jet (sans compter les mots)}
\esp{Hoy, en el Día de las Minorías, recordamos la discriminación histórica que sufren los pueblos indígenas. En mi opinión, la exclusión social y el robo de sus tierras ancestrales son crímenes invisibles. Por ejemplo, en la Amazonía y en Guatemala, las empresas mineras contaminan ríos sin consulta previa. Por eso, reivindicamos que los gobiernos garanticen el consentimiento libre, previo e informado. Es urgente que la ley proteja sus territorios. ¡Exijamos justicia y dignidad para todas las minorías ya!}

\textbf{Étape 3 : Relecture / Comptage / Corrections (Grille d'autoévaluation)}
\begin{center}
\begin{tabular}{|l|c|c|}
\hline
\rowcolor{popblueL}
Critère & Cible & Vérification \\
\hline
Mots & 80-120 & \textbf{98 mots} \checkmark \\
\hline
Lexique leçon (6+) & 6 & \esp{discriminación, exclusión, tierras, derechos, reivindicamos, ley, minorías} (7) \checkmark \\
\hline
Opinion (2) & 2 & \esp{En mi opinión}, \esp{Es urgente que} \checkmark \\
\hline
Revendication + Subj. & 1 & \esp{reivindicamos que... garanticen}, \esp{Es urgente que... proteja} \checkmark \\
\hline
Accords (mujeres/minorías) & OK & \esp{minorías indígenas, tierras ancestrales, minorías} \checkmark \\
\hline
Plan (4 étapes) & Respecté & 1. Hist. 2. Ex. 3. Rev. 4. Appel \checkmark \\
\hline
\end{tabular}
\end{center}

\textbf{Étape 4 : Version finale (polissage connecteurs, ponctuation)}
\tcblower
\textbf{Solution détaillée (Version finale validée)}

\begin{textebox}[Version finale — Día de las Minorías (98 palabras)]
Hoy, en el Día Internacional de las Minorías, denunciamos la \textbf{discriminación histórica} que padecen los pueblos \textbf{indígenas}. \textbf{En mi opinión}, la \textbf{exclusión} sistemática y el despojo de sus \textbf{tierras ancestrales} atentan contra sus \textbf{derechos humanos} fundamentales. \textbf{Por ejemplo}, en la Amazonía y en Guatemala, industrias extractivas destruyen ecosistemas sin \textbf{consulta previa}. \textbf{Por consiguiente, reivindicamos} que los Estados \textbf{garanticen} el consentimiento libre, previo e informado. \textbf{Es urgente que} la \textbf{ley} blinde sus territorios y su autonomía. \textbf{¡Unámonos para exigir justicia y dignidad para todas las \textbf{minorías} ya!}
\end{textebox}
\end{exoresolu}

\courssub{Relecture collective : grille d'autoévaluation}

Avant de rendre votre copie, échangez votre texte avec un camarade et cochez cette grille. C'est la méthode utilisée aux corrections du Bac pour gagner les points de « correction de la langue » et « cohérence ».

\begin{definition}[Grille d'autoévaluation — Texto Reivindicativo (Coefficient 1 par item)]
\begin{center}
\begin{tabularx}{\linewidth}{|X|c|c|c|}
\hline
\rowcolor{poppurpleL}
\textbf{Critère d'évaluation} & \textbf{Oui (1 pt)} & \textbf{Partiel (0,5)} & \textbf{Non (0)} \\
\hline
\textbf{Contenu / Tâche} & & & \\
Respect du nombre de mots (80-120) & & & \\
Respect du plan imposé (4 étapes visibles) & & & \\
Exemple concret (Cameroun ou Hispanophone) cité & & & \\
\textbf{Lexique / Vocabulaire} & & & \\
$\ge$ 6 mots du lexique de la leçon employés correctement & & & \\
Variété lexicale (pas de répétitions \esp{discriminación} 4 fois) & & & \\
\textbf{Grammaire / Structures} & & & \\
2 structures d'opinion différentes (\esp{En mi opinión}, \esp{Considero que}...) & & & \\
1 structure revendication + \textbf{Subjonctif présent} correct (\esp{Exijamos que...}) & & & \\
Accords adjectifs \textbf{Féminin Pluriel} (\esp{mujeres discriminad\emph{as}}) vérifiés & & & \\
Conjugaison subjonctif présente sans faute (\esp{garanticen, respete, proteja}) & & & \\
\textbf{Cohérence / Style} & & & \\
Connecteurs logiques variés (\esp{Por tanto, Por ejemplo, Sin embargo}) & & & \\
Ton adapté (formel, engagé, pas familier) & & & \\
Ponctuation espagnole correcte (\esp{¡ ! ¿ ?}) & & & \\
\hline
\textbf{TOTAL / 12} & \multicolumn{3}{c|}{Objectif Bac : $\ge$ 10/12} \\
\hline
\end{tabularx}
\end{center}
\end{definition}

\courssub{Complément Bac : adapter le plan à un sujet d'expression écrite}

Au Baccalauréat (Épreuve écrite, 2h), le sujet d'expression peut être : \emph{« Redacta un artículo de opinión / una carta al director / un discurso sobre la igualdad de género / la protección de minorías »}. Le plan 4 étapes reste \textbf{la colonne vertébrale} universelle.

\begin{exemplebox}[Adaptation du plan selon le format Bac]
\begin{center}
\begin{tabularx}{\linewidth}{|X|X|X|}
\hline
\rowcolor{popblueL}
\textbf{Format demandé} & \textbf{Adaptation du plan} & \textbf{Formules types (Salutations/Closure)} \\
\hline
\textbf{Carta al director} (Lettre au directeur) & 1. Salutation \esp{Estimado Director:} \newline 2. Plan 4 étapes \newline 3. Formule de politesse \esp{Atentamente,} \newline 4. Signature \esp{Un ciudadano preocupado} & \esp{Estimado Sr. Director:} ... \esp{Atentamente,} \esp{María López, Yaoundé} \\
\hline
\textbf{Artículo de opinión} (Article d'opinion) & 1. \textbf{Titre accrocheur} \newline 2. Plan 4 étapes (développés) \newline 3. Conclusion percutante (phrase choc) & \textbf{Título:} \esp{¡Basta de impunidad!} ... \esp{En definitiva, la justicia no puede esperar.} \\
\hline
\textbf{Discurso} (Discours oral/écrit) & 1. \textbf{Salutation public} \esp{Queridos amigos, compañeras:} \newline 2. Plan 4 étapes (rhétorique : anaphores, questions) \newline 3. \textbf{Slogan final / Cri de ralliement} & \esp{Amigas, amigos:} ... \esp{¡Por la igualdad, hoy y siempre!} \\
\hline
\end{tabularx}
\end{center}
\end{exemplebox}

\begin{aretenir}[Stratégie « Gain de temps » au Bac]
\begin{enumerate}
    \item \textbf{5 min :} Analyse du sujet + Choix du camp + Brainstorming lexique (liste 8 mots sur brouillon).
    \item \textbf{5 min :} Remplissage du squelette (Plan 4 étapes) en espagnol télégraphique.
    \item \textbf{30 min :} Rédaction soignée en phrases complètes + insertion connecteurs + vérification subjonctifs/accords.
    \item \textbf{10 min :} Relecture \textbf{à l'envers} (dernière phrase $\rightarrow$ première) pour traquer les fautes d'accords (\esp{-as/-os}) et accents.
\end{enumerate}
\end{aretenir}

\courssub{Visualisation du processus de production}

Cette frise résume le cheminement cognitif et linguistique pour produire un texte revendicatif abouti, du brouillon à la copie finale.

\begin{popfigure}
\begin{tikzpicture}[
    node distance=1.2cm and 2.5cm,
    every node/.style={font=\small, align=center},
    phase/.style={rectangle, rounded corners, draw=popdark, thick, minimum width=3.2cm, minimum height=1.8cm, text width=3cm},
    arrow/.style={-Stealth, thick, popblue},
    check/.style={circle, draw=popgreen, fill=popgreenL, minimum size=0.6cm, font=\tiny}
]
% Nœuds principaux
\node[phase, fill=poporangeL, align=center] (brouillon){\textbf{1. BROUILLON}\\\emph{5-10 min}\\\textcolor{popdark}{• Brainstorming lexique (8 mots)\\• Choix exemple (Cmr/AmLat)\\• Squelette Plan 4 étapes}};
\node[phase, fill=popblueL, right=of brouillon, align=center] (redaction){\textbf{2. RÉDACTION}\\\emph{25-30 min}\\\textcolor{popdark}{• Phrases complètes\\• Insertion connecteurs\\• \textbf{Subjonctif} : \esp{Exijamos que...}\\• Accords \textbf{-as} (\esp{minorías})}};
\node[phase, fill=popgreenL, right=of redaction, align=center] (relecture){\textbf{3. RELECTURE}\\\emph{10 min}\\\textcolor{popdark}{• Grille autoévaluation\\• Comptage mots (80-120)\\• Vérif. accents \esp{áéíóú}\\• Ponctuation \esp{¿¡}}};
\node[phase, fill=poppurpleL, right=of relecture, align=center] (final){\textbf{4. VERSION FINALE}\\\emph{Prête pour le Bac}\\\textcolor{popdark}{• Cohérent \& Argumenté\\• Lexique riche \& Varié\\• Grammaire maîtrisée\\• Format respecté}};

% Flèches
\draw[arrow] (brouillon.east) -- node[above, font=\tiny, popblue] {Structure} (redaction.west);
\draw[arrow] (redaction.east) -- node[above, font=\tiny, popblue] {Contrôle} (relecture.west);
\draw[arrow] (relecture.east) -- node[above, font=\tiny, popblue] {Validation} (final.west);

% Petits check-points visuels
\node[check, anchor=west] at (brouillon.north east) {✓};
\node[check, anchor=west] at (redaction.north east) {✓};
\node[check, anchor=west] at (relecture.north east) {✓};
\node[check, anchor=west] at (final.north east) {★};

% Légende intégrée TikZ
\node[below=0.5cm of relecture, font=\footnotesize\itshape, text=popmuted, text width=\linewidth, align=center] {Processus itératif : on peut revenir du 3 au 2 si le compte de mots ou le subjonctif échoue.};
\end{tikzpicture}
\legende{Frise chronologique des 4 étapes de production écrite (Brouillon $\rightarrow$ Rédaction $\rightarrow$ Relecture $\rightarrow$ Version finale).}
\end{popfigure}

\courssub{Exercice d'entraînement : à vous de jouer !}

\begin{exercice}[Production autonome — Sujet Bac Blanc]
\textbf{Consigne :} Rédigez un \textbf{texto reivindicativo} de \textbf{90 à 110 mots} pour une \emph{« Carta al Director »} d'un journal camerounais (ex: \esp{Mutations}, \esp{Le Jour} version hispanophone fictive).
\textbf{Sujet :} \emph{« Denuncie la dificultad de acceso a la educación secundaria para las niñas en la zona septentrional de Camerún (Extrême-Nord, Nord, Adamaoua) y exija la aplicación efectiva de la ley de escolarización obligatoria. »}

\textbf{Contraintes rappelées :}
\begin{itemize}
    \item Plan 4 étapes (Problème $\rightarrow$ Exemple Nord-Cameroun $\rightarrow$ Revendication + Subj $\rightarrow$ Appel).
    \item Minimum 6 mots du lexique leçon (\cle{derechos humanos, minoría, discriminación, sufragio, igualdad de oportunidades, ley, reivindicar, excluida, brecha, empoderamiento}).
    \item 2 structures d'opinion + 1 revendication \textbf{subjonctif}.
    \item Vérification accords \esp{niñas, leyes, oportunidades, minorías} $\rightarrow$ \textbf{-as}.
\end{itemize}
\end{exercice}

\begin{corrige}[Exercice 1 — Proposition de correction type (102 palabras)]
\begin{textebox}[Carta al Director: Educación de las niñas en el Norte]
\textbf{Estimado Director:}

Le escribo para denunciar la \textbf{discriminación} estructural que impide el acceso a la educación secundaria de las \textbf{niñas} en el Septentrión de Camerún. \textbf{En mi opinión}, la \textbf{brecha} educativa en la región del Extremo Norte viola sus \textbf{derechos humanos} y frena el \textbf{empoderamiento} de toda una generación. \textbf{Por ejemplo}, en el departamento del Logone y Chari, muchas adolescentes son \textbf{excluidas} del sistema escolar por matrimonios precoces. \textbf{Por consiguiente, reivindicamos} que el Estado \textbf{aplique} con rigor la \textbf{ley} de escolarización obligatoria. \textbf{Es fundamental que} se \textbf{garantice} la \textbf{igualdad de oportunidades} real. \textbf{¡Exijamos justicia educativa ya!}

\textbf{Atentamente,} \\
\textbf{Un ciudadano comprometido, Maroua}
\end{textebox}

\textbf{Analyse de la correction :}
\begin{itemize}
    \item \textbf{Mots :} 102 (dans la fourchette).
    \item \textbf{Lexique (8) :} \esp{discriminación, niñas (contexto minoría), derechos humanos, empoderamiento, brecha, excluidas, ley, igualdad de oportunidades, reivindicamos}.
    \item \textbf{Opinion (2) :} \esp{En mi opinión}, \esp{Es fundamental que}.
    \item \textbf{Revendication + Subj (2) :} \esp{reivindicamos que el Estado aplique}, \esp{Es fundamental que se garantice}.
    \item \textbf{Accords -as :} \esp{niñas, excluidas, ley, igualdad, oportunidades, discriminación} — Tous corrects.
    \item \textbf{Format :} Salutation \esp{Estimado Director:}, cierre \esp{Atentamente}, signature lieu/nom.
\end{itemize}
\end{corrige}

\begin{savaistu}[Le saviez-vous ? — Cameroun \& Guinée équatoriale : femmes et pouvoir]
La \textbf{Guinée équatoriale} (seul pays africain hispanophone) compte des femmes à des postes de haute responsabilité (ministres, vice-ministres, parlementaires). Au \textbf{Cameroun}, bien que la Constitution garantisse l'\esp{igualdad de oportunidades}, le \cle{Code de la famille} (loi n°81-02 du 29 juin 1981, modifiée) contient encore des dispositions critiquées par le Comité CEDAW (notion de chef de famille, gestion des biens du ménage). Les associations comme \emph{RECAFEM} (Réseau des femmes parlementaires) ou \emph{ALVF} (Association de lutte contre les violences faites aux femmes) mènent le \cle{combat} (\esp{la lucha}) pour l'harmonisation des lois avec la CEDAW (Convention sur l'élimination de toutes les formes de discrimination à l'égard des femmes). \esp{¡La lucha continúa!}
\end{savaistu}

\newpage

\coursec{Activité d'intégration}

\coursec{Leçon E03 : Lectura y comentario : derechos de las mujeres y de las minorías}

\courssub{Objectifs de l'activité}

À l'issue de cette activité d'intégration, l'élève sera capable de :

\begin{itemize}
\item réemployer le \cle{lexique des droits humains} (\esp{derechos humanos, minoría, discriminación, sufragio, igualdad de oportunidades, ley, reivindicar}) dans une prise de parole orale argumentée ;
\item mobiliser les \cle{structures d'opinion et de revendication} étudiées dans la leçon, en particulier \esp{es necesario que} + subjonctif, \esp{hay que} + infinitif, \esp{reivindicamos que}, \esp{(no) estoy de acuerdo porque} ;
\item dégager les idées directrices d'un texte court et le point de vue de son auteur pour préparer son propre discours ;
\item réagir à l'oral aux arguments d'autrui de manière polie et structurée ;
\item produire, en groupe, un \esp{texto reivindicativo} bref, correct et convaincant, digne d'être affiché dans la salle de classe.
\end{itemize}

\courssub{Situation}

Le lycée bilingue de Yaoundé célèbre sa \esp{Semana de la Ciudadanía}. À cette occasion, la direction organise un \esp{Foro de debate} intitulé \esp{« Por una sociedad más justa »}. Des élèves de Terminale, des enseignants et quelques invités (une journaliste d'une radio locale, un représentant d'une association de quartier) assisteront au débat. Chaque classe de Terminale doit envoyer quatre groupes de délégués. Votre classe a été choisie pour ouvrir le forum : chacun de vos groupes incarnera un acteur de la société civile et défendra son point de vue sur l'égalité et les droits des minorités.

Pour préparer les débats, le comité d'organisation a distribué à tous les participants le texte d'appel suivant, publié dans le journal du lycée :

\begin{textebox}[Appel du comité d'organisation — journal du lycée]
¡Tu voz cuenta! El próximo viernes, nuestro liceo celebra el Foro « Por una sociedad más justa ». En muchos países, las mujeres y las minorías todavía sufren discriminación: salarios desiguales, acceso difícil a la educación, exclusión de la vida política. Sin embargo, los derechos humanos nos pertenecen a todos, sin distinción de sexo, de lengua o de origen. Es necesario que los jóvenes reivindiquen la igualdad de oportunidades y que participen en la vida ciudadana. \textbf{Vengan} con vuestras ideas, respetad las opiniones de los demás y defended vuestros argumentos con respeto. ¡La sociedad de mañana se construye hoy, con ustedes!
\end{textebox}

\textbf{Glossaire :} \esp{tu voz cuenta} : ta voix compte ; \esp{salarios desiguales} : salaires inégaux ; \esp{acceso difícil} : accès difficile ; \esp{exclusión} : exclusion ; \esp{sin distinción de} : sans distinction de ; \esp{los demás} : les autres ; \esp{se construye} : se construit ; \esp{vuestras} : vos (adjectif possessif, \emph{vosotros}) ; \esp{respetad / defended} : respectez / défendez (impératif \emph{vosotros}) ; \emph{Note : en Amérique latine et en Guinée équatoriale, on préférera} \esp{traigan / respeten / defiendan} \emph{(ustedes).}

\courssub{Compréhension détaillée du texte}

\begin{exercice}[Questions de compréhension]
Répondez en espagnol, par des phrases complètes.
\begin{enumerate}
\item ¿Qué problemas de discriminación menciona el texto?
\item ¿A quién pertenecen los derechos humanos, según el texto?
\item ¿Qué estructura gramatical se usa después de \esp{Es necesario que}? Déduisez le mode du verbe.
\item ¿Cuál es la función del último párrafo? (Llamamiento, amenaza, información pura...)
\end{enumerate}
\end{exercice}

\begin{corrige}[Exercice Compréhension]
\begin{enumerate}
\item \esp{El texto menciona salarios desiguales, acceso difícil a la educación y exclusión de la vida política.}
\item \esp{Los derechos humanos nos pertenecen a todos, sin distinción de sexo, de lengua o de origen.}
\item \esp{Después de « Es necesario que » se usa el subjuntivo (reivindiquen, participen). Expresa necesidad/obligación impersonal.}
\item \esp{Es un llamamiento (llamada a la acción) a la participación ciudadana de los jóvenes.}
\end{enumerate}
\end{corrige}

\courssub{Lexique thématique : los derechos humanos}

\begin{aretenir}[Lexique des droits à réemployer]
\begin{center}
\begin{tabularx}{\linewidth}{|X|X|X|}
\hline
\rowcolor{popblueL} \textbf{Español} & \textbf{Français} & \textbf{Exemple / Collocation} \\
\hline
\esp{los derechos humanos} & les droits humains & \esp{defender los derechos humanos} \\
\hline
\esp{la minoría} & la minorité & \esp{las minorías étnicas / lingüísticas} \\
\hline
\esp{la discriminación} & la discrimination & \esp{sufrir discriminación / discriminar a alguien} \\
\hline
\esp{el sufragio} & le suffrage, le droit de vote & \esp{el sufragio universal / femenino} \\
\hline
\esp{la igualdad de oportunidades} & l'égalité des chances & \esp{garantizar la igualdad de oportunidades} \\
\hline
\esp{la ley} & la loi & \esp{una ley justa / promulgar una ley} \\
\hline
\esp{reivindicar} & revendiquer & \esp{reivindicar un derecho / una ley} \\
\hline
\esp{la ciudadanía} & la citoyenneté & \esp{ejercer la ciudadanía / ciudadanía activa} \\
\hline
\esp{la voz} & la voix & \esp{alzar la voz / tu voz cuenta} \\
\hline
\esp{la brecha salarial} & l'écart salarial & \esp{reducir la brecha salarial} \\
\hline
\esp{el empoderamiento} & l'autonomisation / l'empowerment & \esp{el empoderamiento de las mujeres} \\
\hline
\end{tabularx}
\end{center}
\end{aretenir}

\courssub{Exprimer l'opinion et la revendication}

\begin{propriete}[Structures d'opinion et de revendication]
\begin{itemize}
\item \textbf{Opinion affirmative :} \esp{creo que / pienso que / opino que / para mí / a mi juicio} + \textbf{indicatif}.
      \esp{Creo que la igualdad es fundamental.} (« Je crois que l'égalité est fondamentale. »)
\item \textbf{Opinion négative / doute :} \esp{no creo que / no pienso que / dudo que / no estoy seguro de que} + \textbf{subjonctif}.
      \esp{No creo que sea justo.} (« Je ne crois pas que ce soit juste. »)
\item \textbf{Revendication / Nécessité impersonnelle :} \esp{es necesario que / es importante que / es urgente que / es justo que / es imprescindible que} + \textbf{subjonctif}.
      \esp{Es necesario que el gobierno actúe.} (« Il faut que le gouvernement agisse. »)
\item \textbf{Obligation générale (impersonnelle) :} \esp{hay que} + \textbf{infinitif}.
      \esp{Hay que respetar los derechos de todos.} (« Il faut respecter les droits de tous. »)
\item \textbf{Revendication personnelle (groupe) :} \esp{reivindicamos / exigimos / pedimos que} + \textbf{subjonctif}.
      \esp{Reivindicamos que se apruebe la ley.} (« Nous revendiquons que la loi soit approuvée. »)
\item \textbf{Accord et désaccord :}
      \begin{itemize}
      \item \esp{Estoy de acuerdo porque...} / \esp{Totalmente de acuerdo.}
      \item \esp{No estoy de acuerdo porque...} / \esp{No estoy nada de acuerdo.}
      \item \esp{Tienes razón, pero...} / \esp{Entiendo tu punto, sin embargo...}
      \end{itemize}
\end{itemize}
\end{propriete}

\begin{attention}[Pièges fréquents des francophones]
\begin{itemize}
\item \textbf{Subjonctif obligatoire :} \esp{Es necesario que la ley proteja} (correct) \textbf{\neq} \esp{Es necesario que la ley protege} (faux). Après expression de nécessité/volonté/émotion au principal $\rightarrow$ subjonctif au subordonné.
\item \textbf{Genre de \emph{derecho} :} \esp{los derechos humanos} (masculin pluriel). Jamais \esp{las derechos}.
\item \textbf{\emph{Reivindicar} $\neq$ Revendre :} \esp{Reivindicar} = revendiquer un droit. \esp{Revender} = revendre. Faux ami phonétique.
\item \textbf{\emph{Estar} de acuerdo :} \esp{Estoy de acuerdo}. Jamais \esp{Soy de acuerdo} (calque de \emph{être} d'accord avec \esp{ser}). \esp{Estar} exprime un état, une position.
\item \textbf{Accents toniques (règle des mots oxytones en -\emph{n}, -\emph{s}, voyelle) :} \esp{discriminación, oportunidades, minoría, opinión, igualdad, sociedad, acción, corazón}. L'accent marque l'accentuation sur l'avant-dernière syllabe (paroxytones) ou la dernière (oxytones en -\emph{n}/\emph{s}).
\item \textbf{\emph{Por} / \emph{Para} :} \esp{Luchamos \textbf{por} la igualdad} (cause, motif) vs \esp{Esta ley es \textbf{para} proteger a las mujeres} (but, destination).
\end{itemize}
\end{attention}

\courssub{Méthode du commentaire de texte}

\begin{methode}[Commenter un texte argumentatif (Bac LV2)]
\textbf{Étape 1 : Lecture active \& identification} (10 min)
\begin{enumerate}
\item Lisez le texte deux fois. Soulignez le \textbf{thème} (droits femmes/minorités), la \textbf{thèse} (les jeunes doivent revendiquer l'égalité), le \textbf{ton} (engagé, mobilisateur).
\item Repérez la \textbf{source} (journal du lycée), la \textbf{date} (prochain vendredi), le \textbf{public visé} (élèves, professeurs, invités).
\item Relevez les \textbf{mots-clés} (champ lexical des droits) et les \textbf{connecteurs} (\esp{sin embargo, por eso, y}).
\end{enumerate}

\textbf{Étape 2 : Analyse structurée} (15 min)
\begin{enumerate}
\item \textbf{Introduction :} Accroche (citation, actualité) $\rightarrow$ Présentation du texte (titre, auteur, source, date, thème) $\rightarrow$ Annonce du plan (ex: I. Constat des inégalités ; II. Appel à l'action citoyenne).
\item \textbf{Développement (2 ou 3 parties) :} Une idée directrice par paragraphe. Citez le texte (\esp{« ... »}) et analysez : \emph{Qu'est-ce que l'auteur dit ? Comment le dit-il ? (vocabulaire, syntaxe, figures de style) Pourquoi ? (intention persuasive)}.
\item \textbf{Conclusion :} Bilan des arguments $\rightarrow$ Ouverture (lien avec l'actualité camerounaise, Guinée équatoriale, ODD 5 et 10 de l'ONU).
\end{enumerate}

\textbf{Étape 3 : Rédaction soignée} (15 min)
\begin{itemize}
\item Utilisez les \textbf{connecteurs logiques} : \esp{En primer lugar, además, por un lado... por otro lado, por tanto, en conclusión}.
\item Alternez \textbf{citation} et \textbf{paraphrase}. Montrez votre \emph{visión crítica} (votre avis personnel argumenté) dans la conclusion ou une partie dédiée.
\item Relisez : accords, subjonctifs, accents, \esp{ser/estar}, \esp{por/para}.
\end{itemize}
\end{methode}

\courssub{Production guidée : un texto reivindicativo (Foro de debate)}

\courssub{Consignes}

Travail en groupes de quatre ou cinq élèves. Durée totale : environ une heure (préparation) et une heure (forum et bilan).

\begin{enumerate}
\item \textbf{Lecture et compréhension (10 min).} Lisez l'appel du comité. Relevez dans le texte : (a) deux problèmes de discrimination cités ; (b) deux verbes qui expriment la revendication ; (c) la phrase qui contient \esp{es necesario que} et identifiez le mode du verbe qui suit.
\item \textbf{Répartition des rôles (5 min).} Chaque groupe tire au sort son personnage :
      \begin{itemize}
      \item[(a)] une \esp{asociación de mujeres} de Yaoundé ;
      \item[(b)] un \esp{representante de las minorías} (linguistiques ou culturelles, ex. communauté Baka, anglophones) ;
      \item[(c)] un \esp{periodista} qui présente et modère le forum ;
      \item[(d)] un \esp{diputado} (député) qui défend une nouvelle loi sur l'égalité.
      \end{itemize}
\item \textbf{Préparation du discours (25 min).} Rédigez un discours de 5 à 6 phrases qui respecte le plan suivant : salutation et présentation du rôle ; constat d'un problème concret ; revendication principale ; proposition de solution ; phrase de conclusion.
      \textbf{Obligations de langue :} utilisez au moins quatre mots du lexique de la leçon, au moins une structure \esp{es necesario que} + subjonctif, et une structure d'opinion (\esp{creemos que}, \esp{opinamos que}, \esp{para nosotros}...).
\item \textbf{Le forum (30 min).} Le groupe « journaliste » ouvre la séance et donne la parole à chaque acteur (2 minutes par discours). Après chaque discours, les autres groupes réagissent avec \esp{(no) estoy de acuerdo porque...} suivi d'un argument. Le député conclut en résumant les propositions.
\item \textbf{Vote et affichage (10 min).} La classe vote pour le discours le plus convaincant (clarté, lexique, correction grammaticale, force des arguments). Le discours gagnant est recopié au propre et affiché dans la salle.
\item \textbf{Trace écrite individuelle (10 min).} Chaque élève rédige dans son cahier un court paragraphe (4 phrases) : \esp{¿Qué derecho te parece más importante y por qué?}
\end{enumerate}

\courssub{Ressources à mobiliser}

\begin{aretenir}[Lexique des droits à réemployer — Rappel]
\esp{los derechos humanos} ; \esp{la minoría} ; \esp{la discriminación} ; \esp{el sufragio} ; \esp{la igualdad de oportunidades} ; \esp{la ley} ; \esp{reivindicar} ; \esp{la ciudadanía} ; \esp{la voz} ; \esp{la brecha salarial} ; \esp{el empoderamiento}.
\end{aretenir}

\begin{propriete}[Structures d'opinion et de revendication — Rappel synthétique]
\begin{center}
\begin{tabularx}{\linewidth}{|X|X|X|}
\hline
\rowcolor{popblueL} \textbf{Fonction} & \textbf{Structure} & \textbf{Mode / Forme} \\
\hline
Opinion affirmative & \esp{Creo / Opino / Pienso que...} & Indicatif \\
\hline
Opinion négative / Doute & \esp{No creo / Dudo que...} & Subjonctif \\
\hline
Nécessité / Revendication & \esp{Es necesario / importante / justo / urgente que...} & Subjonctif \\
\hline
Obligation générale & \esp{Hay que + infinitif} & Infinitif \\
\hline
Revendication directe & \esp{Reivindicamos / Exigimos que...} & Subjonctif \\
\hline
Accord & \esp{Estoy de acuerdo porque...} & Indicatif \\
\hline
Désaccord poli & \esp{No estoy de acuerdo porque... / Tienes razón, pero...} & Indicatif \\
\hline
\end{tabularx}
\end{center}
\end{propriete}

\begin{savaistu}[Cameroun, Guinée équatoriale et droits des femmes]
La \textbf{Guinée équatoriale}, seul pays africain hispanophone, a ratifié la CEDAW (Convention sur l'élimination de toutes les formes de discrimination à l'égard des femmes) en 1984. Toutefois, des défis persistent : représentation politique des femmes (quoique en progression), violences basées sur le genre, accès à la terre. Au \textbf{Cameroun}, pays bilingue voisin, la Constitution garantit l'égalité, mais les pratiques coutumières (héritage, mariage précoce dans certaines régions) freinent l'effectivité des droits. Le \esp{Foro} simulé ici reflète la dynamique réelle de la \esp{Sociedad Civil} (ONG, associations de femmes comme \esp{ASAFE} au Cameroun ou \esp{Plataforma de Mujeres} en Guinée équatoriale) qui \esp{reivindican} l'application effective des \esp{leyes}. L'\emph{Objectif de Développement Durable 5} (ODD 5) de l'ONU : \esp{« Lograr la igualdad de género y empoderar a todas las mujeres y las niñas »} cadre parfaitement ce thème.
\end{savaistu}

\courssub{Proposition de production et corrigé commenté}

\begin{exoresolu}[Discours modèle du groupe « asociación de mujeres »]
Rédiger le discours d'ouverture du forum (5 à 6 phrases) pour le groupe qui incarne une association de femmes de Yaoundé, avec le lexique de la leçon et au moins une structure \esp{es necesario que} + subjonctif.
\tcblower
\textbf{Discours modèle :}

\esp{Señoras y señores, queridos compañeros: hablo en nombre de la asociación de mujeres de nuestro barrio. Hoy, en nuestra sociedad, muchas mujeres todavía sufren discriminación en el trabajo y en la educación. Creemos que la igualdad de oportunidades es un derecho humano, no un regalo. Por eso reivindicamos una ley que garantice el mismo salario para el mismo trabajo. Es necesario que el gobierno proteja a las mujeres y que las escuelas eduquen a los niños y a las niñas por igual. Juntos, podemos construir una sociedad más justa. ¡Gracias!}

\textbf{Traduction :} « Mesdames et Messieurs, chers camarades : je parle au nom de l'association des femmes de notre quartier. Aujourd'hui, dans notre société, beaucoup de femmes souffrent encore de discrimination au travail et dans l'éducation. Nous croyons que l'égalité des chances est un droit humain, pas un cadeau. C'est pourquoi nous revendiquons une loi qui garantisse le même salaire pour le même travail. Il faut que le gouvernement protège les femmes et que les écoles éduquent les garçons et les filles de la même manière. Ensemble, nous pouvons construire une société plus juste. Merci ! »

\textbf{Commentaire des choix de langue :}
\begin{itemize}
\item \textbf{Lexique réemployé (6 items) :} \esp{discriminación, igualdad de oportunidades, derecho humano, reivindicamos, ley, sociedad más justa}. Objectif atteint ($>$ 4 mots).
\item \textbf{Subjonctifs corrects :}
      \begin{itemize}
      \item \esp{garantice} (subj. présent, 3e pers. sg. de \esp{garantizar}) après \esp{una ley que} (antécédent indéfini / but recherché). \emph{Note orthographique :} \esp{z $\rightarrow$ c} devant \esp{e}.
      \item \esp{proteja} (subj. présent, 3e pers. sg. de \esp{proteger}) et \esp{eduquen} (subj. présent, 3e pers. pl. de \esp{educar}) après \esp{Es necesario que}.
      \end{itemize}
\item \textbf{Structure d'opinion :} \esp{Creemos que} + indicatif (\esp{es}), car l'opinion est affirmative et certaine pour le locuteur.
\item \textbf{Registre adapté :} salutations formelles (\esp{Señoras y señores}), \emph{nosotros} de représentation collective (\esp{hablo, creemos, reivindicamos, podemos}), formule de politesse finale.
\item \textbf{Progression argumentative :} Salutation $\rightarrow$ Constat (\esp{Hoy... sufren}) $\rightarrow$ Principe éthique (\esp{es un derecho, no un regalo}) $\rightarrow$ Revendication concrète (\esp{reivindicamos una ley}) $\rightarrow$ Mesures institutionnelles (\esp{Es necesario que...}) $\rightarrow$ Appel à l'unité (\esp{Juntos...}).
\end{itemize}

\textbf{Exemple de réaction d'un autre groupe (ex: minorías) :}
\esp{Estoy de acuerdo porque la igualdad salarial es una cuestión de justicia, pero pienso que también hay que cambiar las mentalidades y luchar contra los estereotipos desde la escuela.} « Je suis d'accord car l'égalité salariale est une question de justice, mais je pense qu'il faut aussi changer les mentalités et lutter contre les stéréotypes dès l'école. »
\end{exoresolu}

\courssub{Autres modèles de discours (pour inspiration)}

\begin{exemplebox}[Groupe « Representante de las minorías »]
\esp{Buenos días. Represento a las minorías lingüísticas y culturales de nuestro país. La discriminación no solo es económica: muchos jóvenes no pueden estudiar en su lengua materna y se sienten excluidos del sistema. Opinamos que la diversidad cultural es una riqueza, no un obstáculo. Exigimos que la ley garantice la educación bilingüe y el respeto a nuestras identidades. Es imprescindible que el Estado nos escuche. ¡Nada sobre nosotros sin nosotros!}
\end{exemplebox}

\begin{exemplebox}[Groupe « Diputado / Diputada »]
\esp{Estimados ciudadanos, como diputado me comprometo a presentar un proyecto de ley para la igualdad real. Los datos son claros: persiste la brecha salarial y la violencia de género. Es urgente que el Parlamento apruebe medidas concretas: cuotas de paridad, licencias parentales iguales, sanciones firmes. Hay que legislar con valentía. Cuenten con mi voto y mi voz. ¡La justicia no admite demora!}
\end{exemplebox}

\courssub{Bilan de l'activité}

Cette activité a permis de réinvestir, dans une situation de communication authentique et ancrée dans la réalité camerounaise (Yaoundé, \esp{Semana de la Ciudadanía}), l'ensemble des acquis de la leçon E03.

\begin{itemize}
\item \textbf{Compréhension écrite/orale :} Lecture d'un texte d'appel authentique, identification du thème (discriminations), de la thèse (participation citoyenne) et des procédés d'argumentation (impératif, subjonctif de nécessité).
\item \textbf{Lexique :} Réemploi actif et contextualisé de 10 termes du champ des droits humains (\esp{derechos humanos, minoría, discriminación, sufragio, igualdad de oportunidades, ley, reivindicar, ciudadanía, brecha salarial, empoderamiento}) selon les rôles assignés.
\item \textbf{Grammaire :} Consolidation du point dur pour francophones : \esp{es necesario que} + \textbf{subjonctif} (vs indicatif), distinction \esp{creo que} (ind.) / \esp{no creo que} (subj.), \esp{estar de acuerdo} (vs \esp{ser}), \esp{por/para}.
\item \textbf{Expression orale/écrite :} Production d'un \esp{texto reivindicativo} structuré (salutation, constat, thèse, argument, proposition, conclusion), interaction orale codifiée (accord/désaccord argumenté), gestion du temps de parole.
\item \textbf{Compétence citoyenne :} Débat contradictoire respectueux, prise de conscience des enjeux d'égalité (genre, minorités) au Cameroun et dans l'espace hispanophone (Guinée équatoriale, Espagne, Amérique latine).
\item \textbf{Trace écrite collective :} L'affichage du discours vainqueur sert de modèle pour l'épreuve de \textbf{commentaire de texte} et de \textbf{rédaction argumentative} (essai) du Baccalauréat.
\end{itemize}

\begin{exercice}[Devoir / Approfondissement personnel]
\begin{enumerate}
\item \textbf{Version :} Traduisez en espagnol : « Je ne pense pas que la nouvelle loi soit suffisante pour changer les mentalités. Il est nécessaire que les professeurs éduquent les élèves à l'égalité dès l'école primaire. Nous revendiquons des manuels scolaires non sexistes. »
\item \textbf{Thème :} Traduisez en français : \esp{« Las minorías étnicas en América Latina han sufrido discriminación histórica. Hoy, los pueblos indígenas reivindican sus derechos territoriales y lingüísticos. Es justo que los gobiernos protejan su cultura. »}
\item \textbf{Rédaction (120 mots) :} Rédigez un article pour le journal du lycée : \esp{« La igualdad de género en mi entorno : avances y retos »}. Utilisez au moins 6 mots du lexique, 2 subjonctifs de nécessité, 1 structure d'opinion négative.
\end{enumerate}
\end{exercice}

\begin{corrige}[Devoir / Approfondissement]
\begin{enumerate}
\item \esp{No pienso que la nueva ley sea suficiente para cambiar las mentalidades. Es necesario que los profesores eduquen a los alumnos en la igualdad desde la escuela primaria. Reivindicamos libros de texto no sexistas.}
\item « Les minorités ethniques en Amérique latine ont souffert d'une discrimination historique. Aujourd'hui, les peuples autochtones revendiquent leurs droits territoriaux et linguistiques. Il est juste que les gouvernements protègent leur culture. »
\item \emph{Production libre. Vérification attendue :} Lexique (ex: \esp{brecha salarial, empoderamiento, discriminación, ley, igualdad, minoría}), Subjonctifs (ex: \esp{es necesario que... cambie, es urgente que... garantice}), Opinion négative (ex: \esp{No creo que la situación haya mejorado mucho}).
\end{enumerate}
\end{corrige}

\newpage

\coursec{Méthodes et exercices résolus}

\courssub{Méthodes et exercices résolus}

\begin{methode}[Première lecture : identifier le genre, le thème et la thèse]
\textbf{Objectif :} Saisir le sens global du texte en 5 minutes chrono (temps Bac).
\begin{enumerate}
    \item \textbf{Lire le titre et la source} (journal, blog, discours, essai) : ils orientent le genre \cle{argumentatif}, \cle{informativo} ou \cle{reivindicativo}.
    \item \textbf{Lire la première et la dernière phrase} de chaque paragraphe : elles contiennent souvent les \cle{idées directrices}.
    \item \textbf{Repérer les mots-clés récurrents} (lexique des droits : \esp{igualdad, discriminación, ley, minoría, reivindicar}) et les connecteurs logiques (\esp{sin embargo, por tanto, en consecuencia}).
    \item \textbf{Formuler la thèse} en une phrase : « L'auteur défend l'idée que... » / « Le texte dénonce... et propose... ».
    \item \textbf{Noter le ton} : indigné, espoir, objectif, ironique ? Cela guide le commentaire.
\end{enumerate}
\textbf{Reflexe Bac :} Ne t'arrête pas sur un mot inconnu ; infère le sens par le contexte. Surligne \textbf{uniquement} les mots-clés du sujet.
\end{methode}

\begin{exoresolu}[Compréhension globale : « Voces que rompen el silencio »]
\textbf{Texte support :}
\begin{textebox}[Article d'opinion original -- Contexte Cameroun / Guinée équatoriale / Espagne]
« Voces que rompen el silencio »

En Yaoundé, como en Malabo o Madrid, el camino hacia la \cle{igualdad real} sigue siendo una cuesta arriba. Aunque las constituciones garantizan la \cle{no discriminación} por razón de sexo u origen étnico, la brecha entre la \cle{ley escrita} y la \cle{vida cotidiana} permanece abismal. Las mujeres, que representan más de la mitad de la población, siguen chocando contra el \cle{techo de cristal} en el acceso a los puestos de decisión: apenas un 28\% de escaños parlamentarios en la región CEMAC son ocupados por ellas.

Pero la lucha no es solo numérica. Las \cle{minorías étnicas} --como los \cle{pueblos indígenas} Baka en el sur de Camerún o los \cle{annoboneses} en Guinea Ecuatorial-- sufren una \cle{doble discriminación}: por su identidad y, a menudo, por su pobreza. Sus tierras ancestrales son concesionadas sin \cle{consentimiento previo, libre e informado}, violando el Convenio 169 de la OIT.

La \cle{violencia de género}, esa pandemia silenciosa, no entiende de fronteras. En España, la \cle{ley orgánica 1/2004} fue un hito, pero los \cle{feminicidios} no cesan. En América Latina, el movimiento \esp{« Ni una menos »} ha puesto el acento en la \cle{impunidad} estructural. \cle{Reivindicar} derechos no es pedir favores: es exigir el cumplimiento de la \cle{Declaración Universal de los Derechos Humanos}.

La \cle{igualdad de oportunidades} no cae del cielo. Se construye con \cle{educación no sexista}, \cle{presupuestos con perspectiva de género} y \cle{paridad} obligatoria en las listas electorales. El silencio ya no es una opción. \textbf{¡Basta ya!}
\end{textebox}

\textbf{Traduction / Glossaire :} \esp{cuesta arriba} = pente raide / parcours difficile ; \esp{techo de cristal} = plafond de verre ; \esp{CEMAC} = Communauté Économique et Monétaire de l'Afrique Centrale ; \esp{consentimiento previo, libre e informado} = consentement préalable, libre et éclairé (principe juridique) ; \esp{hito} = jalon, étape majeure ; \esp{presupuestos con perspectiva de género} = budgets genrés (gender budgeting).

\textbf{Énoncé :} Après une première lecture rapide (3 min), répondez en espagnol :
\begin{enumerate}
    \item ¿Cuál es el \textbf{género} del texto? Justifica con un elemento textual.
    \item Identifica el \textbf{tema central} y la \textbf{tesis} del autor en una sola frase.
    \item Señala \textbf{tres grupos humanos} mencionados como víctimas de discriminación.
    \item ¿Qué \textbf{conector} introduce la contra-argumentación en el párrafo 2? ¿Y el de consecuencia en el párrafo 4?
    \item ¿Cuál es el \textbf{tono} predominante? Cita una frase final que lo confirma.
\end{enumerate}

\tcblower
\textbf{Solution détaillée :}
\begin{enumerate}
    \item \textbf{Género :} Texto \cle{argumentativo} / \cle{de opinión} (artículo de opinión). \textbf{Justificación :} Presencia de tesis defendida (\esp{« La igualdad de oportunidades no cae del cielo »}), uso de la primera persona del plural inclusiva (\esp{« nuestra »} implícito), verbos de opinión (\esp{sufren, sigue siendo, hay que construir}) y llamada a la acción final (\esp{¡Basta ya!}). No es solo informativo porque toma partido.
    \item \textbf{Tema :} La persistente desigualdad y discriminación hacia las mujeres y minorías étnicas a pesar de los marcos legales. \textbf{Tesis :} \esp{« Las leyes existen pero no se aplican; hace falta acción política (educación, presupuestos, paridad) para lograr la igualdad real ».}
    \item \textbf{Grupos :} 1. \esp{Las mujeres} (párr. 1 y 3). 2. \esp{Pueblos indígenas / Baka} (párr. 2). 3. \esp{Annoboneses} (minoría en Guinea Ecuatorial, párr. 2). (Aceptable: \esp{mujeres en América Latina / España} como subgrupos).
    \item \textbf{Conectores :} Párrafo 2: \esp{\textbf{Pero}} (adversativo / contra-argumentación frente al marco legal del párrafo 1). Párrafo 4: No hay un conector de consecuencia clásico; la relación lógica se expresa mediante la estructura \esp{\textbf{Se construye con}} (voz pasiva refleja + preposición \esp{con} que indica el medio/instrumento necesario para lograr el fin enunciado en la frase anterior).
    \item \textbf{Tono :} \cle{Indignado / Combativo / Exigente}. \textbf{Frase final :} \esp{« El silencio ya no es una opción. ¡Basta ya! »} (Imperativo exhortativo + frase nominal elíptica de fuerza performativa).
\end{enumerate}
\textbf{Conseil :} Pour la thèse, utilise la structure \esp{« Aunque [réalité légale], [réalité factuelle] ; por tanto, [solution proposée] »}.
\end{exoresolu}

\begin{methode}[Compréhension détaillée : repérer arguments, données et lexique thématique]
\textbf{Objectif :} Répondre aux questions précises (Bac : 6-8 pts) en citant le texte.
\begin{enumerate}
    \item \textbf{Analyser la question} : repérer le mot-clé (cause, conséquence, solution, opinion, chiffre, définition contextuelle).
    \item \textbf{Localiser l'information} : numérote les paragraphes. Utilise les connecteurs (\esp{porque, ya que, debido a} = cause ; \esp{por lo tanto, así que} = conséquence ; \esp{propone, sugiere, exige} = solution).
    \item \textbf{Citer textuellement} (entre guillemets \esp{« ... »}) pour les définitions, chiffres, noms propres. \textbf{Paraphraser} (avec tes mots) pour les explications « \emph{Expliquez...} ».
    \item \textbf{Traiter le lexique} : Pour « \emph{Déduisez le sens de...} », donne un synonyme espagnol (\esp{sinónimo}) ou une définition en espagnol (\esp{definición}) + le contexte qui t'a aidé (indices : préposition, verbe régissant, opposition).
    \item \textbf{Vérifier l'accord} : Si la question porte sur « \emph{les femmes} », ta réponse doit accorder les adjectifs/participes au féminin pluriel.
\end{enumerate}
\textbf{Structure réponse type :} \esp{« Según el texto, [résponse] + [citation/paraphrase]. »}
\end{methode}

\begin{exoresolu}[Compréhension détaillée + Lexique en contexte]
\textbf{Énoncé :} Reprenons le texte \esp{« Voces que rompen el silencio »}.
\begin{enumerate}
    \item \textbf{Données chiffrées :} Quel pourcentage exact de sièges parlementaires les femmes occupent-elles dans la zone CEMAC selon le texte ? Quelle conclusion l'auteur en tire-t-il implicitement ?
    \item \textbf{Argument juridique :} Quelle convention internationale cite l'auteur au paragraphe 2 ? Quel droit spécifique cette convention protège-t-elle pour les peuples indigènes ?
    \item \textbf{Vocabulaire en contexte :} Déduisez le sens de \cle{« techo de cristal »} (p.1) et de \cle{« impunidad »} (p.3) à partir de leur environnement lexical. Donnez un synonyme espagnol pour chacun.
    \item \textbf{Structure argumentative :} Le paragraphe 4 propose des solutions. Citez \textbf{trois mesures concrètes} préconisées par l'auteur (copiez les groupes nominaux exacts).
    \item \textbf{Portée du texte :} En vous appuyant sur les exemples géographiques cités (Yaoundé, Malabo, Madrid, Amérique Latine), quelle est la portée (universelle, locale, régionale) de la problématique traitée ? Justifiez.
\end{enumerate}

\tcblower
\textbf{Solution détaillée :}
\begin{enumerate}
    \item \textbf{Chiffre :} \esp{« Apenas un 28\% de escaños parlamentarios »}. \textbf{Conclusion implicite :} La sous-représentation massive (loin de la parité 50\% ou même du seuil critique de 30\% souvent visé) prouve l'existence du \esp{techo de cristal} mentionné juste avant. L'adverbe \esp{apenas} (à peine) marque l'insuffisance.
    \item \textbf{Convention :} \esp{« Convenio 169 de la OIT »} (Organisation Internationale du Travail). \textbf{Droit protégé :} Le \esp{« consentimiento previo, libre e informado »} (CPLI) concernant l'exploitation de leurs terres ancestrales. C'est un droit collectif spécifique aux peuples indigènes et tribaux.
    \item \textbf{Lexique :}
        \begin{itemize}
            \item \esp{Techo de cristal} : Contexte : \esp{« chocan contra el techo de cristal en el acceso a los puestos de decisión »}. Verbe \esp{chocar} (heurter) + obstacle invisible. \textbf{Sens :} Barrière invisible mais infranchissable empêchant l'accès aux hautes responsabilités. \textbf{Synonyme :} \esp{barrera invisible, límite infranqueable}.
            \item \esp{Impunidad} : Contexte : \esp{« puesto el acento en la impunidad estructural »} après avoir cité les féminicides. \textbf{Sens :} Absence de sanction pour les coupables ; fait que les crimes restent impunis. \textbf{Synonyme :} \esp{falta de castigo, inmunidad de hecho}.
        \end{itemize}
    \item \textbf{Mesures (p.4) :} 1. \esp{educación no sexista} 2. \esp{presupuestos con perspectiva de género} 3. \esp{paridad obligatoria en las listas electorales}. (Ce sont des SN prépositionnels introduits par \esp{con}).
    \item \textbf{Portée :} \cle{Universelle / Mondiale}. \textbf{Justification :} Le texte cite des capitales de trois continents (Afrique : Yaoundé, Malabo ; Europe : Madrid ; Amérique Latine : mouvement \esp{Ni una menos}) et invoque la \esp{Declaración Universal de los Derechos Humanos}. L'auteur montre que le phénomène transcende les frontières nationales.
\end{enumerate}
\textbf{Point méthode :} Pour la Q3, la citation du contexte (\esp{chocan contra...}) est obligatoire pour justifier l'inférence.
\end{exoresolu}

\begin{methode}[Exprimer l'opinion et la revendication : structures grammaticales clés]
\textbf{Objectif :} Maîtriser les structures pour \cle{argumenter}, \cle{nuancer} et \cle{revendiquer} (Bac : Thème, Expression écrite, Oral).
\begin{enumerate}
    \item \textbf{Opinion personnelle (sujet + indicatif) :} \esp{Creo que / Pienso que / Considero que / A mi juicio / Desde mi punto de vista + INDICATIF} (certitude).
    \item \textbf{Opinion subjective / sentiment / doute (subjonctif obligé) :} \esp{Me parece \textbf{injusto} que + SUBJ.} / \esp{Es \textbf{fundamental / urgente / necesario / imprescindible} que + SUBJ.} / \esp{No tolero que + SUBJ.} / \esp{Exijo / Reclamo / Pido que + SUBJ.}
    \item \textbf{Revendication / Exigence (impersonnel + subjonctif ou infinitif) :} \esp{Hay que / Es preciso / Es menester / Urge + INFINITIF} (général). \esp{Es \textbf{inadmisible} que + SUBJ.} / \esp{No se puede permitir que + SUBJ.}
    \item \textbf{Concession / Nuance (subjonctif) :} \esp{Aunque + SUBJ. (fait futur/hypothétique) / INDIC. (fait réel avéré)}. \esp{Pese a que + SUBJ.} \esp{Por mucho que + SUBJ.}
    \item \textbf{Marqueurs du discours revendicatif :} \esp{¡Basta ya de + INF / SUST!} / \esp{¡No más + SUST!} / \esp{Es hora de que + SUBJ. (imparfait)} / \esp{Ya es hora de + INF.}
\end{enumerate}
\textbf{Tableau récapitulatif :}
\begin{center}
\begin{tabularx}{\linewidth}{|X|X|X|}\hline
\rowcolor{popblueL} \textbf{Fonction} & \textbf{Structure type (Espagnol)} & \textbf{Mode / Temps} \\ \hline
Affirmer / Constater & \esp{Es un hecho que / La realidad muestra que} & Indicatif Présent \\ \hline
Exiger / Revendiquer & \esp{Exijo que / Reivindico que / Es imperativo que} & \textbf{Subjonctif Présent} \\ \hline
Proposer une solution & \esp{Propongo que / Sugiero que / Hay que} & Subj. Présent / Infinitif \\ \hline
Dénoncer (indignation) & \esp{Es inaceptable que / Denuncio que} & \textbf{Subjonctif Présent} \\ \hline
Concéder un point & \esp{Admito que / Reconozco que} & Indicatif (fait réel) \\ \hline
Concéder (hypothétique) & \esp{Aunque sea cierto que / Por muy grave que sea} & \textbf{Subjonctif Présent} \\ \hline
\end{tabularx}
\end{center}
\textbf{Attention :} \esp{Creo que} $\rightarrow$ Indicatif. \esp{No creo que} $\rightarrow$ \textbf{Subjonctif}.
\end{methode}

\begin{exoresolu}[Thème / Transformation : du français vers l'espagnol (Structures d'opinion)]
\textbf{Énoncé :} Traduisez en espagnol les phrases suivantes en utilisant les structures de l'opinion et de la revendication vues ci-dessus. Soignez les modes (Indicatif / Subjonctif) et le vocabulaire thématique (\cle{igualdad, discriminación, minoría, ley, reivindicar}).
\begin{enumerate}
    \item Je considère \textbf{qu'il est fondamental} que les femmes accèdent aux postes de pouvoir. (Utiliser \esp{considerar + que + subj.} + \esp{acceder a}).
    \item Bien que la loi interdise la discrimination, la réalité est tout autre. (Utiliser \esp{Aunque + subj.} pour hypothétique/futur ou \esp{ind.} pour fait avéré ? Justifiez votre choix).
    \item Nous exigeons que le gouvernement applique la loi sur la parité \textbf{immédiatement}. (Utiliser \esp{Exigir que + subj.} + \esp{de inmediato}).
    \item Il est temps que les minorités ethniques soient consultées avant tout projet minier. (Utiliser \esp{Es hora de que + subj. imparfait} ou \esp{Ya es hora de + inf.}).
    \item Assez de violence contre les femmes ! (Utiliser \esp{¡Basta ya de + inf. / sust!}).
\end{enumerate}

\tcblower
\textbf{Solution détaillée :}
\begin{enumerate}
    \item \esp{Considero que es fundamental que las mujeres \textbf{accedan} a los puestos de poder.}
        \begin{itemize}
            \item \esp{Considero que} (opinion) + \esp{es fundamental que} (valoration) $\rightarrow$ \textbf{Subjonctif présent} : \esp{accedan} (verbe \esp{acceder}, 3e pers. pl., changement de radical \esp{c$\rightarrow$qu} devant \esp{a} : \esp{accedan}). \esp{A} \esp{puestos de} \esp{poder} (collocation).
        \end{itemize}
    \item \esp{\textbf{Aunque la ley prohíbe} la discriminación, la realidad es otra.}
        \begin{itemize}
            \item \textbf{Choix du mode :} \esp{Aunque + INDICATIF} car « la loi interdit » est un \textbf{fait avéré, connu, réel} (loi existante). Si c'était une loi hypothétique future : \esp{Aunque la ley prohíba}. \esp{Prohíbe} (indic. présent, irrégulier \esp{prohibir} $\rightarrow$ \esp{prohíbe}). \esp{La realidad es otra} (formule fixe).
        \end{itemize}
    \item \esp{Exigimos que el gobierno \textbf{aplique} la ley de paridad \textbf{de inmediato}.}
        \begin{itemize}
            \item \esp{Exigimos que} (volonté/exigence) $\rightarrow$ \textbf{Subjonctif présent} : \esp{aplique} (verbe \esp{aplicar}, 3e pers. sg., subj. \esp{-e}). \esp{De inmediato} = immédiatement (plus formel que \esp{inmediatamente} à l'oral, mais les deux corrects).
        \end{itemize}
    \item \esp{\textbf{Ya es hora de que las minorías étnicas \textbf{fueran} consultadas antes de todo proyecto minero.}}
        \begin{itemize}
            \item \esp{Es hora de que} / \esp{Ya es hora de que} $\rightarrow$ \textbf{Subjonctif imparfait} (\esp{pretérito imperfecto de subjuntivo}) pour exprimer qu'une action \textbf{devrait déjà avoir lieu} (valeur de regret/retard présent). \esp{Fueran consultadas} (passif : \esp{ser} subj. imparf. \esp{fueran} + participe \esp{consultadas} accordé f. pl.). \esp{Antes de} + infinitif (\esp{todo proyecto minero} = tout projet minier).
            \item \textit{Alternative acceptée (moins nuancée) :} \esp{Ya es hora de consultar a las minorías étnicas...} (Infinitif sujet impersonnel).
        \end{itemize}
    \item \esp{\textbf{¡Basta ya de violencia contra las mujeres!} } ou \esp{¡Basta ya de maltrato a las mujeres!}
        \begin{itemize}
            \item \esp{Basta} (verbe \esp{bastar}, 3e pers. sg. présent indicatif, valeur impersonnelle « ça suffit »). \esp{Ya} = déjà / maintenant. \esp{De} + nom/infinitif. \esp{Violencia contra} (collocation standard).
        \end{itemize}
\end{enumerate}
\textbf{Point de grammaire clé :} La \cle{séquence des temps} avec \esp{Es hora de que} impose le subjonctif imparfait (pas le plus-que-parfait) pour exprimer un retard ou un regret présent. C'est un classique du Bac.
\end{exoresolu}

\begin{methode}[Méthode du commentaire de texte (Épreuve écrite Bac : 10-12 pts)]
\textbf{Objectif :} Produire un commentaire structuré en 3 parties (+ intro/conclusion) en 1h30.
\begin{enumerate}
    \item \textbf{INTRODUCTION (10-15 lignes) :}
        \begin{itemize}
            \item \textbf{Accroche} : Contexte large (actualité, citation, chiffre clé) lié au thème.
            \item \textbf{Présentation du texte} : Auteur (si connu), source, date, titre, genre, thème général.
            \item \textbf{Thèse / Problématique} : « Dans ce texte, l'auteur soutient que... » $\rightarrow$ \textbf{Problématique} : « Dans quelle mesure / Comment l'auteur démontre-t-il que... ? »
            \item \textbf{Annonce du plan} : « Nous analyserons d'abord... puis... enfin... » (2 ou 3 axes logiques : Constats $\rightarrow$ Causes $\rightarrow$ Solutions / ou : Femmes $\rightarrow$ Minorités $\rightarrow$ Moyens d'action).
        \end{itemize}
    \item \textbf{DÉVELOPPEMENT (2 ou 3 parties équilibrées, $\approx$ 1 page chacune) :}
        \begin{itemize}
            \item \textbf{Idée directrice} (phrase thèse du paragraphe).
            \item \textbf{Arguments du texte} : Citer / paraphraser (\esp{« Cita textual : "..." »} / \esp{« El autor señala que... »}).
            \item \textbf{Analyse linguistique / stylistique} : Repérer figures de style (anaphore, métaphore \esp{techo de cristal}, antithèse \esp{ley escrita vs vida cotidiana}, champs lexicaux, modes verbaux, connecteurs).
            \item \textbf{Apport personnel / Contexte} : Lier à l'actualité (Cameroun, Guinée Équatoriale, Espagne, \esp{Ni una menos}, OIT 169).
            \item \textbf{Transition} vers la partie suivante.
        \end{itemize}
    \item \textbf{CONCLUSION (5-8 lignes) :}
        \begin{itemize}
            \item \textbf{Bilan synthétique} : Réponse directe à la problématique.
            \item \textbf{Ouverture} : Perspective future, autre texte, autre pays, rôle de la jeunesse, éducation.
        \end{itemize}
    \item \textbf{RELECTURE (10 min) :} Accords (genre/nombre), accents, subjonctifs, connecteurs, cohérence, \textbf{pas de français} dans la copie.
\end{enumerate}
\textbf{Connecteurs de commentaire :} \esp{En primer lugar / En cuanto a / Por un lado... por otro / Cabe destacar que / Esto evidencia que / A modo de conclusión.}
\end{methode}

\begin{exoresolu}[Commentaire guidé : rédaction de l'introduction et d'un paragraphe de développement]
\textbf{Énoncé :} À partir du texte \esp{« Voces que rompen el silencio »}, rédigez en espagnol :
\begin{enumerate}
    \item Une \textbf{introduction complète} (accroche, présentation, thèse, problématique, plan).
    \item Le \textbf{premier paragraphe de développement} (axe : « Le constat d'un fossé entre droit formel et réalité vécue »).
\end{enumerate}
\textbf{Contraintes :} 150-200 mots au total. Vocabulaire précis. Zéro faute d'accord. Subjonctifs maîtrisés.

\tcblower
\textbf{Proposition de corrigé (niveau Bac 16-18/20) :}

\textbf{Introducción}
En un mundo donde la \cle{Declaración Universal de los Derechos Humanos} cumple más de siete décadas, la \cle{brecha} entre el \cle{marco jurídico} y la \cle{realidad efectiva} sigue siendo la gran asignatura pendiente. El artículo de opinión \esp{« Voces que rompen el silencio »}, publicado en un contexto de reivindicación global, aborda la persistente \cle{desigualdad} que sufren las mujeres y las \cle{minorías étnicas} en espacios tan diversos como Camerún, Guinea Ecuatorial, España o América Latina. La \cle{tesis central} del autor es clara: la existencia de leyes no garantiza la \cle{igualdad real}; se requiere una voluntad política concreta —educación, presupuestos, paridad— para derribar los muros de la \cle{discriminación} y la \cle{impunidad}. \textbf{¿De qué manera el texto articula la denuncia de esta brecha y qué soluciones propone para transformar el derecho escrito en derecho vivido?} Analizaremos primero el diagnóstico de la desigualdad estructural, luego la especificidad de la \cle{doble discriminación} de las minorías, y finalmente las vías de acción política esbozadas.

\textbf{Desarrollo (Eje 1): El abismo entre la ley escrita y la vida cotidiana}
El autor abre el texto estableciendo una \cle{antítesis} contundente: \esp{« Aunque las constituciones garantizan la no discriminación..., la brecha... permanece abismal »}. El uso de la concesiva \esp{aunque} + \textbf{indicativo} (\esp{garantizan}) reconoce el avance legal, pero el adjetivo \esp{abismal} y la metáfora del \esp{« techo de cristal »} (p.1) visualizan la opacidad y la solidez de la barrera que frena el acceso femenino al poder (\esp{« apenas un 28\% »}). Este dato cuantitativo, lejos de ser anecdótico, funciona como \cle{argumento de autoridad} empírico: la \cle{subrepresentación} en la zona CEMAC prueba que la \cle{paridad} es una quimera sin \cle{medidas correctoras} obligatorias. Asimismo, la oposición léxica \esp{« ley escrita » / « vida cotidiana »} subraya que el derecho positivo, si no va acompañado de \cle{políticas públicas} efectivas, se convierte en \cle{letra muerta}. En el contexto camerunés o ecuatoguineano, donde el \cle{costumbre} a veces prima sobre la ley estatal, esta brecha adquiere dimensiones dramáticas para las mujeres rurales.

\textbf{Analyse du corrigé :}
\begin{itemize}
    \item \textbf{Intro :} Accroche sur DUDH. Présentation complète. Thèse claire. Problématique opératoire (Comment / De quelle manière). Plan en 3 axes logiques.
    \item \textbf{Développement :} Idée directrice (antithèse loi/réalité). Citation intégrée. Analyse grammaticale (\esp{aunque + ind.}) + stylistique (métaphore, adjectif, chiffre). Lien contexte Cameroun/Guinée. Vocabulaire précis (\cle{subrepresentación, letra muerta, medidas correctoras, costumbre}). Connecteurs (\esp{Asimismo, En cuanto a}).
\end{itemize}
\end{exoresolu}

\begin{methode}[Production écrite : rédiger un « texto reivindicativo » (Lettre ouverte, Discours, Article)]
\textbf{Objectif :} Écrire un texte court (150-200 mots) engagé, structuré et linguistiquement riche (Expression écrite Bac : 8-10 pts).
\begin{enumerate}
    \item \textbf{Analyser le sujet} : Format (Lettre au ministre, Discours ONU, Article blog, Pétition) $\rightarrow$ Détermine le \cle{registre} (formel / solennel / journalistique) et le \cle{destinataire} (\esp{tú} / \esp{usted} / \esp{ustedes} / impersonnel).
    \item \textbf{Structurer en 3 temps :}
        \begin{itemize}
            \item \textbf{Exordium (Accroche) :} Interpellation directe (\esp{Señor Ministro / Queridos compañeros / Ciudadanos}), énoncé du problème choc (chiffre, fait divers, citation).
            \item \textbf{Argumentatio (Corps) :} 2 arguments forts. Structure : \esp{Argumento 1 + Ejemplo/Dato + Conector + Argumento 2 + Consecuencia}. Utiliser \cle{subjonctif} pour exigences (\esp{Exijo que...}), \cle{conditionnel} pour propositions (\esp{Sería necesario que...}), \cle{futur} pour engagement (\esp{Seguiremos luchando}).
            \item \textbf{Peroratio (Conclusion / Appel à l'action) :} Formule choc, slogan, signature collective (\esp{Atentamente, Las mujeres de... / Firmado: Red de Jóvenes por la Igualdad}).
        \end{itemize}
    \item \textbf{Soigner la langue :} \cle{Champ lexical des droits} (obligatoire). \cle{Variété syntaxique} (passif \esp{se exige}, impersonnel \esp{es urgente}, gérondif \esp{exigiendo}, participes). \cle{Cohésion} (connecteurs : \esp{por tanto, además, sin embargo, en definitiva}).
    \item \textbf{Relire :} Accords (genre/nombre), accents (\esp{acción, nación, corazón}), \esp{ser/estar} (identité vs état), \esp{por/para} (cause/but).
\end{enumerate}
\textbf{Astuce :} Prépare 3 « phrases-passe-partout » adaptables : \esp{« No podemos permitir que... »}, \esp{« Es inaceptable que en el siglo XXI... »}, \esp{« Exigimos el cumplimiento inmediato de... »}.
\end{methode}

\begin{exoresolu}[Rédaction : Lettre ouverte à la Ministre de la Promotion de la Femme (Cameroun)]
\textbf{Énoncé :} Vous êtes délégué(e) d'une association lycéenne « Jeunesse pour l'Égalité » à Douala. Écrivez une \textbf{lettre ouverte formelle} à Madame la Ministre de la Promotion de la Femme et de la Famille (MINPROFF) pour réclamer l'application effective de la loi sur la parité dans les conseils d'administration des grandes entreprises publiques (SNEC, CAMTEL, CDC, etc.). Longueur : $\approx$ 180 mots. Registre formel (\esp{usted}).

\tcblower
\textbf{Proposition de corrigé :}

\begin{textebox}[Lettre ouverte -- Registre formel (Usted)]
Douala, 15 de octubre de 2025

\textbf{Señora Ministra de la Promoción de la Mujer y la Familia,}

Le dirijo la presente en nombre de la asociación \esp{« Juventud por la Igualdad »} para \cle{denunciar} la \cle{brecha inaceptable} que persiste entre la \cle{ley de paridad} (Ley N° 2024/...), votada hace dos años, y la \cle{composición real} de los consejos de administración de nuestras empresas públicas emblemáticas: SNEC, CAMTEL, CDC, SONARA.

\textbf{Es un hecho constatable} que, a día de hoy, la \cle{presencia femenina} en esos órganos de decisión no supera el 15\%, muy lejos del 30\% mínimo legal y del 50\% de justicia democrática. \textbf{No se trata de un problema de « falta de candidatas cualificadas »}, como a veces se arguye, sino de una \cle{falta de voluntad política} para hacer cumplir la norma. \esp{« Exigimos que se apliquen sanciones efectivas a las entidades que incumplan los cupos »} y que se publique anualmente un \cle{informe de transparencia} con datos desagregados por sexo.

\textbf{Es imperativo que el Ministerio lidere} con el ejemplo: \textbf{condicione} las subvenciones públicas al respeto de la paridad, \textbf{instaure} la formación obligatoria en \cle{liderazgo femenino} y \textbf{visibilice} los referentes. \esp{« No queremos promesas, queremos resultados medibles »}.

La \cle{igualdad de oportunidades} no es una concesión, es un \cle{derecho humano}. \textbf{El silencio de las instituciones ya no es una opción.}

Atentamente, \\
\textit{Ngo Mbock Chantal} \\
Portavoz, Juventud por la Igualdad -- Douala
\end{textebox}

\textbf{Analyse linguistique du corrigé :}
\begin{itemize}
    \item \textbf{Forme :} En-tête lieu/date. Formule d'appel \esp{Señora Ministra...}. Signature avec titre.
    \item \textbf{Structures d'opinion/revendication :} \esp{Le dirijo... para denunciar} (infinitif de but). \esp{Es un hecho constatable que + INDIC.} \esp{No se trata de... sino de...} (structure corrective). \esp{Exigimos que + SUBJ. (se apliquen)}. \esp{Es imperativo que + SUBJ. (lidere, condicione, instaure, visibilice)} $\leftarrow$ \textbf{Subjonctif présent 3e pers. sg. (verbes en -er/-ir : -a ; -ar : -e)}. \esp{No queremos promesas, queremos resultados} (parallélisme/antithèse).
    \item \textbf{Lexique thématique :} \cle{brecha, ley de paridad, consejos de administración, cupos, informe de transparencia, subvenciones, liderazgo femenino, derecho humano}.
    \item \textbf{Grammaire :} Accords parfaits (\esp{presencia femenina, entidades, sanciones efectivas, candidatas cualificadas}). \esp{Ser/Estar} : \esp{Es un hecho} (identité/definition), \esp{persiste} (verbe d'état). \esp{Por/Para} : \esp{para denunciar} (but), \esp{por la Igualdad} (cause).
\end{itemize}
\end{exoresolu}

\begin{attention}[Les 5 erreurs qui coûtent des points au Bac -- Spécial Droits / Opinion]
\begin{enumerate}
    \item \textbf{Subjonctif oublié après \cle{exigir / reclamar / es urgente que / es fundamental que}.}
        \esp{\textcolor{red}{Exijo que el gobierno *aplica* la ley}} $\rightarrow$ \esp{\textcolor{popgreen}{Exijo que el gobierno \textbf{aplique} la ley}}. \textbf{Règle :} Volonté / Émotion / Doute / Nécessité $\rightarrow$ Subjonctif.
    \item \textbf{Faux ami : \cle{« Minorité » $\neq$ \esp{minoría} (seul) pour « minorité d'âge ».} }
        \esp{Minoría} = groupe minoritaire (ethnique, politique). « Minorité d'âge » = \esp{menor de edad}. « Être minoritaire » = \esp{estar en minoría}.
    \item \textbf{Accord du participe passé / adjectif avec \cle{mujer / minoría / ley} (Féminin Singulier).}
        \esp{\textcolor{red}{La ley *escrito* / Las mujeres *discriminado* / Una minoría *discriminado*}} $\rightarrow$ \esp{\textcolor{popgreen}{La ley \textbf{escrita} / Las mujeres \textbf{discriminadas} / Una minoría \textbf{discriminada}}}. \textbf{Relisez tous les accords F/Sg et F/Pl.}
    \item \textbf{\cle{Ser vs Estar} avec \cle{igualdad / discriminación / derecho}.}
        \esp{\textcolor{red}{La igualdad *está* garantizada}} (Correct : état résultant d'une action). MAIS \esp{\textcolor{red}{La discriminación *es* estructural}} (Caractéristique essentielle, identité). \esp{El derecho *es* humano} (définition). \textbf{Test :} « C'est dans sa nature » $\rightarrow$ \esp{Ser}. « C'est un état/résultat » $\rightarrow$ \esp{Estar}.
    \item \textbf{Anglicisme / Calque : \cle{« Opportunité » $\neq$ \esp{oportunidad} dans « égalité des opportunités ».} }
        \esp{Igualdad de \textbf{oportunidades}} (Pluriel ! Locution figée). \esp{Oportunidad} au singulier = « occasion », « chance ». Aussi : \esp{« Lutter pour »} = \esp{luchar por} (correct) mais \esp{« Revendiquer quelque chose »} = \esp{reivindicar \textbf{algo}} (transitif direct, pas de préposition). \esp{Reivindicar \textbf{por} algo} $\rightarrow$ \textcolor{red}{FAUX}.
\end{enumerate}
\textbf{Dernier conseil :} Apprenez par cœur 5 phrases-types « passe-partout » avec subjonctif (\esp{Es inadmisible que..., Exijo que..., Propongo que..., Es hora de que..., No tolero que...}) : elles sauvent la note à l'oral et à l'écrit.
\end{attention}

\newpage

\coursec{Exercices}

\courssub{Je vérifie mes connaissances}

\begin{exercice}[QCM : comprensión del texto]
Pour chaque question, entoure la lettre de la bonne réponse.
\begin{enumerate}
\item El texto estudia sobre todo :
\begin{enumerate}
\item la historia del fútbol femenino ; \item los derechos de las mujeres y de las minorías ; \item la economía de América Latina.
\end{enumerate}
\item « Reivindicar » significa :
\begin{enumerate}
\item renunciar a un derecho ; \item pedir con firmeza algo que se considera justo ; \item castigar a los culpables.
\end{enumerate}
\item El sufragio es :
\begin{enumerate}
\item el derecho a votar ; \item el derecho a trabajar ; \item el derecho a la educación.
\end{enumerate}
\item La « igualdad de oportunidades » significa que :
\begin{enumerate}
\item todos ganan el mismo sueldo siempre ; \item todos tienen las mismas posibilidades de acceder al empleo y a la educación ; \item las leyes son iguales en todos los países.
\end{enumerate}
\end{enumerate}
\end{exercice}

\begin{exercice}[Vrai ou faux ? Justifie]
Indique si las frases siguientes son verdaderas (V) o falsas (F) y justifica tu respuesta con una cita del texto support, traducida al francés.
\begin{enumerate}
\item La discriminación contra las mujeres ha desaparecido por completo. (V / F)
\item Las minorías también sufren exclusión social. (V / F)
\item Las leyes bastan siempre para garantizar la igualdad real. (V / F)
\item El autor del texto es pesimista sobre el futuro. (V / F)
\end{enumerate}
\end{exercice}

\begin{exercice}[Vocabulaire : los derechos humanos]
Trouve l'équivalent français des mots suivants, puis utilise chacun d'eux dans une phrase espagnole de ta composition : \esp{derechos humanos}, \esp{minoría}, \esp{discriminación}, \esp{sufragio}, \esp{igualdad de oportunidades}, \esp{ley}, \esp{reivindicar}.
\end{exercice}

\begin{exercice}[Grammaire : opinion et subjonctif]
Complète les phrases avec la forme correcte du verbe entre parenthèses (indicatif ou subjonctif).
\begin{enumerate}
\item Creo que las mujeres \dots{} (tener) los mismos derechos que los hombres.
\item No creo que la ley \dots{} (ser) suficiente para acabar con la discriminación.
\item Es cierto que las minorías \dots{} (sufrir) exclusión.
\item No es verdad que todos \dots{} (poder) acceder a la educación.
\item Pienso que el sufragio universal \dots{} (ser) una conquista fundamental.
\item Dudo que el gobierno \dots{} (querer) cambiar la situación.
\end{enumerate}
\end{exercice}

\courssub{Je m'entraîne}

\begin{exercice}[Traduction (thème)]
Traduis en espagnol les phrases suivantes.
\begin{enumerate}
\item Les femmes réclament l'égalité des chances dans le monde du travail.
\item Je crois que la discrimination est une violation des droits de l'homme.
\item Il faut que la loi protège les minorités.
\item Au Cameroun, beaucoup de jeunes filles ne peuvent pas aller à l'école.
\item Nous ne pensons pas que cette situation soit juste.
\end{enumerate}
\end{exercice}

\begin{exercice}[Transformation : de l'indicatif au subjonctif]
Transforme les phrases suivantes : remplace \esp{Creo que} par \esp{No creo que} et fais tous les changements nécessaires.
\begin{enumerate}
\item Creo que las mujeres tienen las mismas oportunidades.
\item Creo que la ley castiga la discriminación.
\item Creo que las minorías pueden participar en la vida política.
\item Creo que el gobierno respeta los derechos humanos.
\item Creo que la sociedad cambia rápidamente.
\end{enumerate}
\end{exercice}

\begin{exercice}[Texte à trous : réemploi lexical]
Complète le paragraphe suivant avec les mots de la liste : \esp{derechos humanos}, \esp{minoría}, \esp{discriminación}, \esp{sufragio}, \esp{igualdad de oportunidades}, \esp{ley}, \esp{reivindicar}.

\begin{textebox}[Un párrafo para completar]
En muchos países, las mujeres y las \dots{} étnicas siguen sufriendo \dots{} en el empleo y en la educación. Aunque el \dots{} universal fue conquistado hace tiempo, la \dots{} real todavía no existe. Por eso, miles de personas salen a la calle para \dots{} una nueva \dots{} que garantice los \dots{} de todos los ciudadanos.
\end{textebox}
\end{exercice}

\begin{exercice}[Appariements : expressions de revendication]
Associe chaque expression espagnole (1-6) à son équivalent français (a-f).
\begin{center}
\begin{tabularx}{\linewidth}{|X|X|}
\hline
\rowcolor{popblueL} Español & Français \\
\hline
1. Exigimos que se respeten nuestros derechos. & a. Nous dénonçons toute forme de discrimination. \\
\hline
2. Denunciamos toda forma de discriminación. & b. Il est indispensable que la loi change. \\
\hline
3. Es indispensable que la ley cambie. & c. Nous exigeons que nos droits soient respectés. \\
\hline
4. Pedimos la igualdad de oportunidades. & d. Nous luttons pour un monde plus juste. \\
\hline
5. Luchamos por un mundo más justo. & e. Nous demandons l'égalité des chances. \\
\hline
6. Es necesario que todos participemos. & f. Il est nécessaire que nous participions tous. \\
\hline
\end{tabularx}
\end{center}
\end{exercice}

\courssub{Je me prépare au Bac}

\begin{exercice}[Compréhension de texte type Bac]
Lis le texte suivant, puis réponds en espagnol aux questions.

\begin{textebox}[Las hijas del mercado]
En el gran mercado de la ciudad, Rosa vende frutas desde hace veinte años. Cada mañana, antes del amanecer, prepara su puesto mientras sus hijas estudian a la luz de una lámpara. « Yo no pude ir a la escuela », cuenta Rosa, « pero mis hijas serán médicas o maestras. Es mi manera de luchar ».

Sin embargo, la situación no es fácil. Muchas familias del barrio todavía piensan que una niña no necesita estudiar, y las jóvenes que terminan sus estudios encuentran a menudo puertas cerradas en el mundo del trabajo. « Nos piden experiencia, pero no nos dan la primera oportunidad », explica Amina, la hija mayor de Rosa, que sueña con ser abogada para defender a las mujeres de su comunidad.

A pesar de todo, Rosa no se rinde. Con otras vendedoras del mercado, ha creado una asociación que ayuda a las niñas a continuar sus estudios. « Solas somos débiles », dice, « pero juntas podemos cambiar las cosas. La educación es el derecho que nadie nos puede quitar ».
\end{textebox}

\textbf{Glossaire :} \esp{amanecer} = lever du jour ; \esp{rendirse} = abandonner ; \esp{puertas cerradas} = portes fermées ; \esp{juntas} = ensemble.

\textbf{Questions :}
\begin{enumerate}
\item ¿A qué se dedica Rosa y desde cuándo ? (2 pts)
\item ¿Por qué quiere Rosa que sus hijas estudien ? Justifica con una cita del texto. (2 pts)
\item ¿Qué problema encuentran las jóvenes en el mundo del trabajo ? Cita el texto y traduce la cita al francés. (3 pts)
\item ¿Qué ha hecho Rosa para cambiar la situación ? (2 pts)
\item ¿Cuál es la idea principal del texto ? Exprésala en dos frases con tus propias palabras. (3 pts)
\item ¿Qué significa la frase « La educación es el derecho que nadie nos puede quitar » ? ¿Estás de acuerdo ? Justifica tu respuesta. (4 pts)
\item Grammaire : dans la phrase \esp{« Muchas familias piensan que una niña no necesita estudiar »}, transforme en phrase négative d'opinion (\esp{No creo que\dots}) avec le subjonctif. (2 pts)
\item Grammaire : trouve dans le texte deux verbes au présent de l'indicatif et conjuge-les à toutes les personnes. (2 pts)
\end{enumerate}
\end{exercice}

\begin{exercice}[Thème et version]
\textbf{A. Thème.} Traduis en espagnol le texte suivant. (10 pts)

« Les droits de l'homme concernent tout le monde : les hommes, les femmes et les minorités. Malheureusement, la discrimination existe encore dans beaucoup de pays. Je crois que l'éducation est la meilleure arme pour lutter contre l'injustice. Il faut que les gouvernements fassent des lois pour garantir l'égalité des chances. Nous devons tous réclamer un monde plus juste. »

\textbf{B. Version.} Traduis en français le texte suivant. (10 pts)

\begin{textebox}[Una conquista histórica]
El sufragio femenino fue una conquista larga y difícil. Durante muchos años, las mujeres no podían votar ni participar en la vida política. Gracias a la lucha de muchas personas valientes, hoy las mujeres pueden elegir a sus representantes. Sin embargo, todavía queda mucho por hacer: la igualdad real exige cambios en las leyes, pero también en las mentalidades.
\end{textebox}
\end{exercice}

\begin{exercice}[Expression écrite type Bac]
\textbf{Sujet.} Has leído en un periódico que en tu ciudad una empresa ha rechazado a una joven ingeniera solo porque es mujer. Indignado/a, decides escribir una carta al director del periódico para denunciar esta discriminación y reivindicar la igualdad.

Escribe la carta en español (\textbf{100-120 palabras}). Debes :
\begin{enumerate}
\item presentar brevemente los hechos que denuncias ;
\item expresar tu opinión con al menos dos estructuras de opinión (\esp{me parece que}, \esp{estoy convencido/a de que}, \esp{no creo que} + subjuntivo\dots) ;
\item formular al menos dos reivindicaciones (\esp{exigimos que}, \esp{es necesario que}, \esp{pedimos}\dots) ;
\item utilizar al menos cuatro palabras del lexique de la lección : \esp{derechos humanos}, \esp{discriminación}, \esp{igualdad de oportunidades}, \esp{ley}, \esp{reivindicar}, \esp{minoría}.
\end{enumerate}

\textbf{Barème indicatif :} respect de la consigne et de la longueur (4 pts) ; correction grammaticale (6 pts) ; richesse du vocabulaire (4 pts) ; organisation et cohérence (4 pts) ; présentation de la lettre (2 pts).
\end{exercice}

\newpage

\coursec{Corrigés des exercices}

\courssub{Corrigés des exercices d'application}

\begin{corrige}[Exercice 1 — QCM : Lexique et structures]
\textbf{Réponses correctes :}
\begin{enumerate}
    \item \textbf{B.} \esp{Reivindicar} signifie « revendiquer, exiger quelque chose comme un droit ». Exemple : \esp{Las mujeres reivindican la igualdad salarial.} « Les femmes revendiquent l'égalité salariale. » \textbf{Attention au faux ami :} \esp{reivindicar} ne signifie jamais « revendiquer » au sens commercial (réclamer un bien), ni « réhabiliter » (ce serait \esp{rehabilitar}).
    \item \textbf{B.} \esp{haga}. Après \esp{no creo que} (opinion niée, doute), l'espagnol exige le \cle{subjonctif présent} : \esp{No creo que el gobierno haga suficiente\ldots} « Je ne crois pas que le gouvernement fasse assez\ldots » La forme \esp{hace} (indicatif) serait correcte seulement après \esp{creo que} affirmatif.
    \item \textbf{C.} \esp{sufragio}. Le \esp{sufragio} est le droit de vote : \esp{sufragio universal} « suffrage universel », \esp{sufragio femenino} « droit de vote des femmes ».
    \item \textbf{B.} \esp{Exijo que + subjonctif}. Les verbes de volonté et d'exigence (\esp{exigir, pedir, reclamar, querer}) sont suivis de \esp{que} + \cle{subjonctif} : \esp{Exijo que se respeten los derechos de las minorías.} « J'exige que les droits des minorités soient respectés. »
\end{enumerate}
\textbf{Barème indicatif :} 1 point par réponse correcte, soit 4 points.
\end{corrige}

\begin{corrige}[Exercice 2 — Vrai ou Faux ?]
\begin{enumerate}
    \item \textbf{Faux.} Après \esp{creo que} à la forme affirmative, on emploie l'\cle{indicatif} car le locuteur présente son idée comme un fait certain : \esp{Creo que la ley protege a las minorías.} C'est à la forme négative (\esp{no creo que}) que le \cle{subjonctif} apparaît, car la réalité du fait devient douteuse ou non affirmée.
    \item \textbf{Vrai.} La \esp{discriminación} (traitement inégal fondé sur le sexe, l'origine, la langue, le handicap\ldots) est précisément ce que l'\esp{igualdad de oportunidades} cherche à éliminer : les deux notions s'opposent dans le champ lexical des droits humains.
    \item \textbf{Faux.} Une \esp{minoría} est tout groupe numériquement minoritaire dans une société donnée : elle peut être ethnique, mais aussi linguistique (par exemple les locuteurs d'une langue minoritaire), religieuse ou sociale (par exemple les personnes en situation de handicap).
    \item \textbf{Faux.} La \esp{ley} est un instrument juridique indispensable, mais non suffisant : il faut aussi l'éducation, la prise de conscience collective (\esp{la concienciación}) et l'action des associations de la société civile pour que les droits soient réellement effectifs.
\end{enumerate}
\textbf{Barème indicatif :} 0,5 point pour le choix V/F + 0,5 point pour la justification argumentée (règle grammaticale citée ou exemple espagnol), soit 4 points au total.
\textbf{Point de vigilance :} une justification sans exemple espagnol ni règle grammaticale explicite ne vaut que la moitié des points.
\end{corrige}

\begin{corrige}[Exercice 3 — Vocabulaire : Complétion et définitions]
\textbf{Propositions de correction} (les formes attendues sont les mots de la liste, accordés en genre et en nombre si nécessaire) :
\begin{enumerate}
    \item \esp{La \textbf{ley} garantiza los derechos de todos los ciudadanos.} « La loi garantit les droits de tous les citoyens. »
    \item \esp{Los \textbf{derechos humanos} son universales e inalienables.} « Les droits humains sont universels et inaliénables. »
    \item \esp{Las mujeres luchan \textbf{por} la \textbf{igualdad de oportunidades} en el trabajo.} « Les femmes luttent pour l'égalité des chances au travail. »
    \item \esp{El \textbf{sufragio} femenino fue una conquista histórica.} « Le droit de vote des femmes fut une conquête historique. »
    \item \esp{Cada \textbf{minoría} tiene derecho a conservar su lengua y su cultura.} « Chaque minorité a le droit de conserver sa langue et sa culture. »
    \item \esp{La \textbf{discriminación} racial está prohibida por la Constitución.} « La discrimination raciale est interdite par la Constitution. »
    \item \esp{Los ciudadanos salen a la calle \textbf{para} \textbf{reivindicar} sus derechos.} « Les citoyens descendent dans la rue pour revendiquer leurs droits. »
\end{enumerate}
\textbf{Points de vigilance :}
\begin{itemize}
    \item \esp{Reivindicar} se conjugue comme \esp{explicar} : changement orthographique \esp{c} $\rightarrow$ \esp{qu} devant \esp{e} / \esp{i} pour garder le son [k] (\esp{yo reivindico} $\rightarrow$ \esp{yo reivindique} au subjonctif présent).
    \item Ne pas confondre \esp{la ley} (la loi, texte voté) et \esp{el derecho} (le droit, prérogative) : on dit \esp{respetar la ley} mais \esp{defender los derechos}.
    \item \esp{Igualdad} est féminin : \esp{la igualdad}, jamais \esp{*el igualdad}.
    \item Prépositions : \esp{luchar \textbf{por} algo} (cause, motif) vs \esp{salir \textbf{para} + infinitif} (but, finalité).
\end{itemize}
\textbf{Barème indicatif :} 1 point par phrase correctement complétée (mot juste + forme correcte + préposition le cas échéant), soit 7 points.
\end{corrige}

\newpage

\coursec{Fiche bilan}

\begin{popfigure}
\begin{tikzpicture}[
  mindmap,
  concept color=popblue,
  level 1 concept/.append style={level distance=3.5cm, sibling angle=360/7},
  level 2 concept/.append style={level distance=2.5cm},
  every node/.style={font=\small, align=center},
  branch1/.style={concept color=poporange},
  branch2/.style={concept color=popgreen},
  branch3/.style={concept color=poppurple},
  branch4/.style={concept color=poppink},
  branch5/.style={concept color=popgold},
  branch6/.style={concept color=popblue!70!popgreen},
  branch7/.style={concept color=popmuted}
]
\node[concept, text=white, font=\bfseries\large, align=center]{Derechos:\\Mujeres y Minorías}
  child[branch1] { node[concept, text=white] {Texto soporte} }
  child[branch2] { node[concept, text=white] {Léxico clave} }
  child[branch3] { node[concept, text=white, align=center]{Estructuras\\opinión/reivindicación} }
  child[branch4] { node[concept, text=white, align=center]{Método\\comentario} }
  child[branch5] { node[concept, text=white, align=center]{Producción\\escrita} }
  child[branch6] { node[concept, text=white] {Contexto cultural} }
  child[branch7] { node[concept, text=white, align=center]{Errores\\frecuentes} };
\end{tikzpicture}
\legende{Carte mentale récapitulative de la leçon E03}
\end{popfigure}

\begin{bilan}

\textbf{Léxico clave : los derechos humanos y las minorías}

\begin{center}
\begin{tabularx}{\linewidth}{|X|X|l|}
\hline
\rowcolor{popblueL}
\textbf{Español} & \textbf{Français} & \textbf{Categoría} \\
\hline
\cle{derechos humanos} & droits de l'homme & Sust. pl. \\
\cle{minoría} & minorité & Sust. f. \\
\cle{discriminación} & discrimination & Sust. f. \\
\cle{sufragio} & suffrage, droit de vote & Sust. m. \\
\cle{igualdad de oportunidades} & égalité des chances & Loc. nom. f. \\
\cle{ley} & loi & Sust. f. \\
\cle{reivindicar} & revendiquer & Verbe \\
\cle{exigir} & exiger & Verbe \\
\cle{denunciar} & dénoncer & Verbe \\
\cle{brecha de género} & écart/fossé entre les sexes & Loc. nom. f. \\
\cle{empoderamiento} & autonomisation, empowerment & Sust. m. \\
\cle{paridad} & parité & Sust. f. \\
\cle{estereotipo} & stéréotype & Sust. m. \\
\cle{violencia de género} & violence sexiste / basée sur le genre & Loc. nom. f. \\
\hline
\end{tabularx}
\end{center}

\textbf{Grammaire à connaître par cœur : Structures pour exprimer l'opinion et la revendication}

\begin{center}
\begin{tabularx}{\linewidth}{|X|X|X|}
\hline
\rowcolor{poporangeL}
\textbf{Fonction} & \textbf{Structure (Espagnol)} & \textbf{Mode / Note} \\
\hline
Opinion personnelle (certitude) & \esp{Creo que / Pienso que / Opino que} + \textbf{Indicatif} & Indicatif (réalité) \\
\hline
Opinion négative / Doute & \esp{No creo que / No pienso que / Dudo que} + \textbf{Subjonctif} & Subjonctif (subjectivité) \\
\hline
Valoración impersonnelle & \esp{Es justo / injusto / inadmisible / urgente / necesario que} + \textbf{Subjonctif} & Subjonctif (jugement de valeur) \\
\hline
Revendication / Exigence & \esp{Exijo / Reivindico / Demando / Pido que} + \textbf{Subjonctif} & Subjonctif (volonté) \\
\hline
Proposition / Solution & \esp{Propongo / Sugiero / Recomiendo que} + \textbf{Subjonctif} & Subjonctif \\
\hline
Concession & \esp{Aunque / A pesar de que} + \textbf{Subjonctif} (fait futur/hypothétique) / \textbf{Indicatif} (fait réel) & Choix selon réalité \\
\hline
Conséquence & \esp{Por tanto / Por consiguiente / En consecuencia / Así que} + \textbf{Indicatif} & Connecteurs logiques \\
\hline
Cause & \esp{Porque / Ya que / Puesto que / Debido a que} + \textbf{Indicatif} & Connecteurs causaux \\
\hline
\end{tabularx}
\end{center}

\textbf{Méthode express : Commentaire de texte (4 étapes)}

\begin{enumerate}
\item \textbf{Lectura global :} Identifier le \cle{tipo de texto} (artículo de opinión, manifiesto, discurso, informe), el \cle{tema} central, la \cle{tesis} del autor y el \cle{tono} (denuncia, propuesta, informativo).
\item \textbf{Análisis detallado :} Estructura (introducción, desarrollo, conclusión), argumentos (datos, ejemplos, citas), conectores, léxico específico, estructuras gramaticales clave (subjuntivo en reivindicaciones).
\item \textbf{Síntesis :} Redactar un resumen objetivo (1/3 del texto) usando \emph{style indirecto} y conectores (\esp{en primer lugar, además, finalmente}).
\item \textbf{Comentario crítico personal :} Situar el texto en su contexto (histórico, social, camerunés/ecuatoguineano), valorar la argumentación, aportar ejemplos propios y concluir con una reflexión abierta.
\end{enumerate}

\end{bilan}

\begin{aretenir}[Checklist antes del Bac]
$\square$ Dominar el léxico de los derechos (\emph{derechos humanos, minoría, discriminación, sufragio, igualdad, reivindicar}).\\
$\square$ Distinguir \textbf{Indicativo} (certeza: \esp{creo que}) vs \textbf{Subjuntivo} (duda/voluntad/valoración: \esp{no creo que, exijo que, es urgente que}).\\
$\square$ Conjugar sin error los verbos clave: \emph{reivindicar, exigir, denunciar, proponer, lograr, conseguir} (Presente, Pretérito indefinido, Subjuntivo presente).\\
$\square$ Identificar la \cle{tesis} y los \cle{argumentos} del texto soporte en menos de 10 minutos.\\
$\square$ Estructurar el comentario en 4 partes claras: Introducción, Resumen, Análisis, Opinión personal.\\
$\square$ Usar conectores lógicos variados (\esp{por tanto, sin embargo, además, en consecuencia, por un lado... por otro}).\\
$\square$ Redactar un \emph{texto reivindicativo} coherente: tesis clara + 2-3 argumentos + propuesta concreta + conclusión fuerte.\\
$\square$ Conocer hitos clave: Sufragio femenino (España 1931), CEDAW (1979), ODS 5 (Agenda 2030), Ley de paridad.\\
$\square$ Integrar referencias al contexto \textbf{Camerún / Guinea Ecuatorial} (leyes locales, ONGs, realidades del mercado, bendskins, cacao).\\
$\square$ Evitar \textbf{falsos amigos}: \esp{realizar} (hacer, no realizar), \esp{discriminar} (distinguir/tratardiferente, connotación negativa), \esp{pretender} (tener intención, no pretender).\\
$\square$ Respetar \textbf{ortografía y puntuación española}: acentos (léxico, gramática), signos de interrogación/exclamación de apertura (\esp{¿? ¡!}), diéresis (\esp{pingüino, vergüenza}).\\
$\square$ Gestionar el tiempo: 45 min lectura/análisis, 45 min redacción comentario, 30 min producción escrita.
\end{aretenir}

\findecours

\end{document}
